\documentclass[11pt,a4paper,oneside]{report}

% ==============================================================================
% PAQUETS FONDAMENTAUX & TYPOGRAPHIE ACADÉMIQUE (CHARTE ENSAM RABAT)
% ==============================================================================
\usepackage[utf8]{inputenc}
\usepackage[T1]{fontenc}
\usepackage{lmodern} % Police vectorielle Latin Modern
\usepackage[french]{babel}
\usepackage{geometry}
\geometry{top=2.4cm,bottom=2.4cm,left=2.6cm,right=2.4cm}
\usepackage{setspace}
\setstretch{1.13}

\usepackage{microtype,graphicx,float,booktabs,tabularx,array,longtable,multirow}
\usepackage{amsmath,amssymb,amsfonts,amsthm}
\usepackage{xcolor,tikz}
\usetikzlibrary{arrows.meta,positioning,shapes.geometric,fit,calc,trees}
\usepackage{pgfgantt}
\usepackage{tcolorbox}
\tcbuselibrary{skins,breakable}
\usepackage{listings}
\usepackage{titlesec,fancyhdr,hyperref,caption}
\usepackage{subcaption}
\usepackage{xurl}

% ==============================================================================
% CHARTE GRAPHIQUE ET COULEURS INSTITUTIONNELLES
% ==============================================================================
\definecolor{EnsamBlue}{RGB}{21,60,85}       % Bleu ENSAM officiel
\definecolor{bleu}{HTML}{153C55}            % Bleu nuit profond
\definecolor{vert}{HTML}{087F73}            % Vert canard / Teal santé
\definecolor{gris}{HTML}{F0F4F5}            % Gris clair fond encadré
\definecolor{primary}{RGB}{21,60,85}        % Alias couleur primaire
\definecolor{secondary}{HTML}{087F73}       % Alias couleur secondaire
\definecolor{accent}{RGB}{197, 48, 48}       % Rouge alerte / contrainte
\definecolor{darkgray}{RGB}{50, 50, 50}
\definecolor{lightgray}{HTML}{F0F4F5}
\definecolor{placeholderbg}{HTML}{F8FAFC}

\hypersetup{
    colorlinks=true,
    linkcolor=black,
    urlcolor=vert,
    citecolor=EnsamBlue
}

% ==============================================================================
% FORMATAGE DES TITRES (TITRESEC - STYLE ENSAM)
% ==============================================================================
\titleformat{\chapter}[display]
{\normalfont\bfseries\color{black}}
{\Large Chapitre \thechapter}
{9pt}
{\Huge}
\titlespacing*{\chapter}{0pt}{-5pt}{17pt}

\titleformat{\section}
{\large\bfseries\color{black}}
{\thesection}{0.7em}{}

\titleformat{\subsection}
{\normalsize\bfseries\color{black}}
{\thesubsection}{0.6em}{}

\titleformat{\subsubsection}
{\small\bfseries\color{black}}
{\thesubsubsection}{0.5em}{}

\setlength{\parindent}{0pt}
\setlength{\parskip}{5pt}
\captionsetup{font=small,labelfont=bf}

% Répertoires de recherche des figures
\graphicspath{{figures/}{./}}
\DeclareGraphicsExtensions{.png,.jpg,.jpeg,.pdf}

% ==============================================================================
% EN-TÊTES ET PIEDS DE PAGE (FANCYHDR - CONFORME AU MODÈLE ENSAM)
% ==============================================================================
\pagestyle{fancy}
\fancyhf{}
\fancyhead[L]{\small\color{gray}Planification Prédictive des Ressources Hospitalières}
\fancyhead[R]{\small\color{gray}ENSAM Rabat -- INDIA}
\fancyfoot[C]{\thepage}
\renewcommand{\headrulewidth}{0.4pt}
\renewcommand{\footrulewidth}{0pt}

% ==============================================================================
% MACROS ET ENCADRÉS POUR CAPTURES ET NOTES
% ==============================================================================
% Boîte de note académique
\newtcolorbox{notebox}[1]{
    enhanced,
    colback=gris,
    colframe=EnsamBlue,
    fonttitle=\bfseries\small,
    title=#1,
    boxrule=0.8pt,
    arc=2pt
}

% Encadré pour réservations de captures d'écran
\newtcolorbox{imageplaceholder}[1]{
    enhanced,
    colback=gris,
    colframe=vert,
    fonttitle=\bfseries\small,
    title=Emplacement réservé : #1,
    boxrule=0.9pt,
    borderline={0.6pt}{0pt}{vert,dashed},
    arc=3pt,
    breakable
}

% Macro d'insertion de capture simplifiée
\newcommand{\capture}[2]{%
    \fcolorbox{gray!45}{gris}{%
        \parbox[c][3.2cm][c]{0.88\linewidth}{%
            \centering\textbf{\color{EnsamBlue}Capture à insérer : #1}\\[4pt]%
            \small\color{darkgray}Remplacer cet encadré par une capture réelle du projet.%
        }%
    }%
}

\newcommand{\captureimg}[2][0.95\textwidth]{%
  \IfFileExists{figures/#2.png}{\includegraphics[width=#1,height=0.45\textheight,keepaspectratio]{figures/#2.png}}{%
  \IfFileExists{#2.png}{\includegraphics[width=#1,height=0.45\textheight,keepaspectratio]{#2.png}}{%
  \IfFileExists{figures/#2.jpg}{\includegraphics[width=#1,height=0.45\textheight,keepaspectratio]{figures/#2.jpg}}{%
  \IfFileExists{#2.jpg}{\includegraphics[width=#1,height=0.45\textheight,keepaspectratio]{#2.jpg}}{%
  \fcolorbox{gray!45}{gris}{%
      \parbox[c][3.2cm][c]{0.86\linewidth}{%
          \centering\textbf{\color{EnsamBlue}Schéma / Image à insérer : #2}\\[4pt]%
          \small\color{darkgray}Déposer le fichier image \texttt{#2.png} dans le dossier \texttt{figures/}.%
      }%
  }%
  }}}}}

\newcommand{\capturebox}[2]{%
    \begin{figure}[H]%
        \centering%
        \capture{#1}{#2}%
        \caption{#2}%
    \end{figure}%
}

% Configuration des listings de code
\lstdefinestyle{code}{
    basicstyle=\ttfamily\footnotesize\color{darkgray},
    backgroundcolor=\color{gris},
    frame=single,
    rulecolor=\color{gray!35},
    breaklines=true,
    showstringspaces=false,
    columns=fullflexible,
    keywordstyle=\color{EnsamBlue}\bfseries,
    commentstyle=\color{vert}\itshape,
    stringstyle=\color{accent},
    numbers=left,
    numberstyle=\tiny\color{gray},
    tabsize=2
}
\lstset{style=code}

% Définition du langage JSON
\lstdefinelanguage{json}{
    basicstyle=\normalfont\ttfamily\footnotesize,
    numbers=left,
    numberstyle=\scriptsize\color{gray},
    stepnumber=1,
    numbersep=8pt,
    showstringspaces=false,
    breaklines=true,
    frame=single,
    backgroundcolor=\color{gris}
}

% Chiffres du rapport, générés par scripts/generate_tables.py
% Généré par rapport/scripts/generate_tables.py : ne pas modifier à la main.
\newcommand{\AbsentsSt}{35}
\newcommand{\AgentsCent}{464}
\newcommand{\CapacitePostes}{1\,175}
\newcommand{\CibleSz}{2\,322}
\newcommand{\CouvDemSqGreedy}{91{,}0\,\%}
\newcommand{\CouvDemSzGreedy}{99{,}0\,\%}
\newcommand{\CouvGreedySc}{50{,}6\,\%}
\newcommand{\CouvGreedySd}{44{,}6\,\%}
\newcommand{\CouvGreedySq}{38{,}4\,\%}
\newcommand{\CouvGreedySt}{41{,}7\,\%}
\newcommand{\CouvGreedySu}{47{,}9\,\%}
\newcommand{\CouvGreedySz}{48{,}3\,\%}
\newcommand{\CouvMaxTous}{60\,\%}
\newcommand{\CouvMinTous}{33\,\%}
\newcommand{\CouvNuitInfUrgCp}{16\,\%}
\newcommand{\CouvSzCp}{51{,}1\,\%}
\newcommand{\CouvSzCpChi}{70\,\%}
\newcommand{\CouvSzCpMed}{25\,\%}
\newcommand{\CouvSzCpRea}{93\,\%}
\newcommand{\CouvSzCpUrg}{45\,\%}
\newcommand{\CouvSzGreedy}{48{,}3\,\%}
\newcommand{\CouvSzGreedyChi}{64\,\%}
\newcommand{\CouvSzGreedyMed}{20\,\%}
\newcommand{\CouvSzGreedyRea}{90\,\%}
\newcommand{\CouvSzGreedyUrg}{48\,\%}
\newcommand{\CouvSzMilp}{50{,}6\,\%}
\newcommand{\CouvSzMilpChi}{86\,\%}
\newcommand{\CouvSzMilpMed}{26\,\%}
\newcommand{\CouvSzMilpRea}{91\,\%}
\newcommand{\CouvSzMilpUrg}{30\,\%}
\newcommand{\CouvSzMilpUrgAccueil}{48\,\%}
\newcommand{\CouvSzNaive}{41{,}2\,\%}
\newcommand{\CouvSzNaiveChi}{51\,\%}
\newcommand{\CouvSzNaiveMed}{35\,\%}
\newcommand{\CouvSzNaiveRea}{49\,\%}
\newcommand{\CouvSzNaiveUrg}{36\,\%}
\newcommand{\DateDebut}{2023-01-01}
\newcommand{\DateFin}{2025-12-30}
\newcommand{\DeficitCpAccSq}{1\,541}
\newcommand{\DeficitCpAccSz}{1\,147}
\newcommand{\DeficitCpEqSq}{1\,526}
\newcommand{\DeficitCpEqSz}{1\,123}
\newcommand{\DeficitCpMaxAcc}{5\,\%}
\newcommand{\DeficitCpMaxEq}{10\,\%}
\newcommand{\DeficitCpMaxPer}{21\,\%}
\newcommand{\DeficitCpMinAcc}{1\,\%}
\newcommand{\DeficitCpMinEq}{2\,\%}
\newcommand{\DeficitCpMinPer}{9\,\%}
\newcommand{\DeficitCpPerSq}{1\,421}
\newcommand{\DeficitCpPerSz}{1\,045}
\newcommand{\DeficitGreedySc}{1\,096}
\newcommand{\DeficitGreedySd}{1\,408}
\newcommand{\DeficitGreedySq}{1\,564}
\newcommand{\DeficitGreedySt}{1\,353}
\newcommand{\DeficitGreedySu}{1\,218}
\newcommand{\DeficitSzCp}{1\,123}
\newcommand{\DeficitSzGreedy}{1\,200}
\newcommand{\DeficitSzMilp}{1\,147}
\newcommand{\DeficitSzNaive}{1\,365}
\newcommand{\EquiteGreedySq}{5}
\newcommand{\EquiteGreedySt}{6}
\newcommand{\GainMaxAcc}{3{,}1\,\%}
\newcommand{\GainMaxEq}{5{,}0\,\%}
\newcommand{\GainMaxPer}{7{,}9\,\%}
\newcommand{\GainMinAcc}{0{,}4\,\%}
\newcommand{\GainMinEq}{1{,}4\,\%}
\newcommand{\GainMinPer}{2{,}8\,\%}
\newcommand{\GapMaxAccueilCp}{0{,}5\,\%}
\newcommand{\GapMaxAccueilMilp}{0{,}1\,\%}
\newcommand{\GapMilpPerSt}{19{,}2\,\%}
\newcommand{\GapMilpPerSz}{16{,}8\,\%}
\newcommand{\LignesDispo}{771\,975}
\newcommand{\LimiteMed}{17\,\%}
\newcommand{\MaeValUrg}{4{,}91}
\newcommand{\MilpPerSt}{$+$12{,}9\,\%}
\newcommand{\MilpPerSz}{$+$6{,}7\,\%}
\newcommand{\NaiveMax}{13{,}9\,\%}
\newcommand{\NaiveMin}{4{,}6\,\%}
\newcommand{\NbAffectations}{7\,193}
\newcommand{\NbAgents}{235}
\newcommand{\NbAides}{75}
\newcommand{\NbCas}{80}
\newcommand{\NbContraintes}{8\,136}
\newcommand{\NbCpMeilleur}{17}
\newcommand{\NbDeficits}{252}
\newcommand{\NbEndpoints}{16}
\newcommand{\NbInfirmiers}{110}
\newcommand{\NbJours}{1\,095}
\newcommand{\NbLits}{300}
\newcommand{\NbMedecins}{50}
\newcommand{\NbPrevTest}{28}
\newcommand{\NbPrevVal}{84}
\newcommand{\NbReprisesCbc}{1}
\newcommand{\NbResSens}{160}
\newcommand{\NbResSensKo}{40}
\newcommand{\NbResSensOk}{120}
\newcommand{\NbTestsNavigateur}{4}
\newcommand{\NbTestsPython}{95}
\newcommand{\NbVariables}{7\,563}
\newcommand{\NuitsMax}{3}
\newcommand{\ObjSzCp}{604\,800}
\newcommand{\ObjSzGreedy}{636\,379}
\newcommand{\ObjSzMilp}{609\,870}
\newcommand{\ObjSzNaive}{718\,885}
\newcommand{\PaliersCinqEq}{6}
\newcommand{\PaliersCinqPer}{11}
\newcommand{\PaliersDiffEq}{9}
\newcommand{\PaliersDiffPer}{15}
\newcommand{\PartCibleChi}{20\,\%}
\newcommand{\PartCibleMed}{37\,\%}
\newcommand{\PartCibleRea}{18\,\%}
\newcommand{\PartCibleUrg}{26\,\%}
\newcommand{\PartDemMax}{23\,\%}
\newcommand{\PartDemMin}{2\,\%}
\newcommand{\PartRhMax}{95\,\%}
\newcommand{\PartRhMin}{75\,\%}
\newcommand{\PlafondSz}{50{,}6\,\%}
\newcommand{\PostesMax}{5}
\newcommand{\QualifChi}{50\,\%}
\newcommand{\QualifMed}{26\,\%}
\newcommand{\QualifRea}{50\,\%}
\newcommand{\QualifUrg}{24\,\%}
\newcommand{\RecensementMed}{113}
\newcommand{\RefusCpAccSq}{73}
\newcommand{\RefusCpAccSz}{4}
\newcommand{\RefusCpEqSq}{76}
\newcommand{\RefusCpEqSz}{9}
\newcommand{\RefusCpPerSq}{113}
\newcommand{\RefusCpPerSz}{27}
\newcommand{\RefusGreedySc}{26}
\newcommand{\RefusGreedySd}{73}
\newcommand{\RefusGreedySq}{73}
\newcommand{\RefusGreedySt}{4}
\newcommand{\RefusGreedySu}{29}
\newcommand{\RefusMilpPerSt}{221}
\newcommand{\RefusMilpPerSz}{169}
\newcommand{\RefusSzCp}{9}
\newcommand{\RefusSzGreedy}{4}
\newcommand{\RefusSzMilp}{4}
\newcommand{\RefusSzNaive}{4}
\newcommand{\ReposMin}{11}
\newcommand{\SaturesSq}{16 sur 28}
\newcommand{\SensCouvAbs}{39{,}9\,\%}
\newcommand{\SensCouvBase}{48{,}3\,\%}
\newcommand{\SensCouvFermee}{50{,}6\,\%}
\newcommand{\SensCouvPire}{35{,}8\,\%}
\newcommand{\SensNaiveAbs}{34{,}0\,\%}
\newcommand{\SensNaiveBase}{41{,}2\,\%}
\newcommand{\SensRefusBase}{4}
\newcommand{\SensRefusCq}{188}
\newcommand{\SensRefusDx}{434}
\newcommand{\SensRefusFermee}{26}
\newcommand{\SensRefusVc}{73}
\newcommand{\TempsSolveur}{50}
\newcommand{\ThreadsSolveur}{8}
\newcommand{\VivierA}{117}
\newcommand{\VivierB}{118}
\newcommand{\WapeTestChi}{28{,}7\,\%}
\newcommand{\WapeTestMed}{19{,}9\,\%}
\newcommand{\WapeTestRea}{35{,}4\,\%}
\newcommand{\WapeTestUrg}{15{,}8\,\%}
\newcommand{\WapeValChi}{25{,}78\,\%}
\newcommand{\WapeValMed}{17{,}25\,\%}
\newcommand{\WapeValRea}{37{,}26\,\%}
\newcommand{\WapeValUrg}{14{,}08\,\%}


% ==============================================================================
% DÉBUT DU DOCUMENT
% ==============================================================================
\begin{document}

% Page de garde académique conforme au modèle ENSAM
\begin{titlepage}
\thispagestyle{empty}

% ==================================================
% LOGOS
% ==================================================
\vspace*{-0.5cm}

\noindent
\begin{minipage}[c]{0.49\textwidth}
    \raggedright
    \includegraphics[height=2.8cm,keepaspectratio]{logo-Ensam-rabat-.png}
\end{minipage}
\begin{minipage}[c]{0.49\textwidth}
    \raggedleft
    \includegraphics[height=2.8cm,keepaspectratio]{Mohammed_V_Logo.png}
\end{minipage}

\vspace{1cm}

% ==================================================
% EN-TÊTE
% ==================================================
\begin{center}

{\fontsize{17}{21}\selectfont
\bfseries PROJET ACADÉMIQUE}\\[0.25cm]

{\fontsize{14}{18}\selectfont
\color{gray}
MODULE : SANTÉ DIGITALE\\
ET INGÉNIERIE DES SYSTÈMES DE SANTÉ}\\[0.3cm]

{\fontsize{14}{16}\selectfont
Filière : Ingénierie Numérique en Data Science
et Intelligence Artificielle}

\vspace{1cm}

% ==================================================
% TITRE DU PROJET
% ==================================================
{\color{EnsamBlue}\rule{\textwidth}{2pt}}

\vspace{0.6cm}

\begin{minipage}{0.96\textwidth}
\centering
{\fontsize{24}{32}\selectfont
\bfseries\color{EnsamBlue}
Système d'Aide à la Décision pour la
Planification Prédictive des Ressources
Hospitalières}
\end{minipage}

\vspace{0.6cm}

{\color{EnsamBlue}\rule{\textwidth}{2pt}}

\end{center}

\vspace{1.2cm}

% ==================================================
% ENCADREMENT
% ==================================================
\begin{center}

{\fontsize{12}{15}\selectfont
\bfseries\itshape\color{EnsamBlue}
Encadré par :}\\[0.4cm]

{\fontsize{13}{17}\selectfont
Pr. ZEMZAMI Maria}

\vspace{1.1cm}

% ==================================================
% RÉALISÉ PAR
% ==================================================
{\fontsize{12}{15}\selectfont
\bfseries\itshape\color{EnsamBlue}
Réalisé par :}

\vspace{0.5cm}

\begin{tabular}{ccc}
\fontsize{12}{16}\selectfont\bfseries ABDELLAOUI Hajar &
\fontsize{12}{16}\selectfont\bfseries AFIF Mohamed &
\fontsize{12}{16}\selectfont\bfseries CHAOUKI Basma \\[0.45cm]

\fontsize{12}{16}\selectfont\bfseries EL FEKAK Hafsa &
\fontsize{12}{16}\selectfont\bfseries EL HOUR Maroua &
\fontsize{12}{16}\selectfont\bfseries LACHGAR Ridouan
\end{tabular}

\end{center}

\vfill

% ==================================================
% ANNÉE UNIVERSITAIRE
% ==================================================
\begin{center}
    {\color{EnsamBlue}\rule{0.55\textwidth}{0.8pt}}\\[0.35cm]

    {\fontsize{12}{15}\selectfont
    \bfseries\color{gray}
    Année universitaire 2026 -- 2027}
\end{center}

\end{titlepage}

\clearpage
\thispagestyle{empty}
\mbox{}
\clearpage


% Liminaires (numérotation romaine)
\clearpage
\pagenumbering{roman}
\chapter*{Remerciements}
\addcontentsline{toc}{chapter}{Remerciements}

Ce travail de projet a été mené au sein de l'École Nationale Supérieure d'Arts et Métiers (ENSAM Rabat), dans le cadre du module \textbf{Santé Digitale} de la filière d'ingénieur \textbf{INDIA}.

Nous tenons à exprimer notre profonde gratitude à l'ensemble du corps professoral et des intervenants pédagogiques de la filière INDIA pour leur accompagnement bienveillant, la richesse des enseignements dispensés et l'exigence méthodologique qui a guidé nos travaux tout au long de ce semestre.

Nos vifs remerciements s'adressent particulièrement à notre professeure encadrante du module Santé Digitale, le \textbf{Prof. Maria ZEMZAMI}, dont les orientations méthodologiques, les précieux conseils et le suivi rigoureux ont grandement enrichi notre démarche de modélisation et guidé la réalisation de ce travail.

Nous remercions également la direction de l'ENSAM Rabat pour la mise à disposition des infrastructures technologiques et des environnements numériques qui ont rendu possible la conduite de nos expérimentations algorithmiques.

Enfin, nous adressons nos sincères remerciements à nos collègues étudiants de la filière INDIA pour l'esprit d'émulation, la solidarité et la richesse des échanges intellectuels partagés au cours de ce projet.


\clearpage
\chapter*{Résumé}
\addcontentsline{toc}{chapter}{Résumé}

La gestion des lits et du personnel d'un hôpital est un problème décisionnel difficile : les admissions sont incertaines, les règles de travail sont nombreuses et les services dépendent les uns des autres. Réalisé dans le cadre du module \textbf{Santé digitale et ingénierie des systèmes de santé}, ce projet présente la conception, l'implémentation et l'évaluation d'un outil d'aide à la décision reproductible pour un hôpital synthétique de \textbf{\NbLits{} lits répartis sur quatre services (Urgences, Médecine, Réanimation, Chirurgie)} et \NbAgents{} professionnels, sur un horizon de \textbf{7 jours ($J+1$ à $J+7$)}. Aucune donnée patient réelle n'est utilisée.

La chaîne de traitement comprend :
\begin{enumerate}
    \item La génération et le contrôle de données synthétiques couvrant \NbJours{} jours ;
    \item La comparaison de cinq modèles de prévision des admissions, validés chronologiquement (WAPE de validation de \WapeValUrg{} aux Urgences pour la régression à retards retenue) ;
    \item Une formulation unique de l'ouverture des lits et de l'affectation individuelle du personnel sur trois postes, avec repos minimal de \ReposMin{}\,h, \PostesMax{} postes par semaine au plus et dotation minimale ;
    \item Quatre méthodes comparées sur les mêmes instances (stratégie naïve, heuristique Greedy, MILP/CBC, CP-SAT) et trois profils de pondération (équilibré, priorité accueil, priorité personnel), soit 72 plans contrôlés par un validateur indépendant ;
    \item Six scénarios ($S_0$ nominal à $S_5$ fermeture de lits) et une grille de sensibilité de \NbCas{} cas ;
    \item Une API FastAPI et un tableau de bord Next.js de dix vues, dont un simulateur de scénario personnalisé ;
    \item Un cadre réglementaire, normatif et organisationnel : loi 09-08 et CNDP, HL7 FHIR, ISO/IEC 27001 et ISO 27799.
\end{enumerate}

Les 72 plans respectent toutes les contraintes dures. CP-SAT est la meilleure méthode, ou ex æquo, dans \NbCpMeilleur{} cas sur 18 ; avec le profil équilibré, la meilleure méthode améliore l'objectif de Greedy de \GainMinEq{} à \GainMaxEq{}. Le résultat principal est cependant un \textbf{arbitrage entre accueil et personnel} : réduire le déficit de personnel conduit à refuser davantage de patients. L'hôpital étudié est en sous-effectif structurel : en situation nominale, Greedy accepte \CouvDemSzGreedy{} de la demande mais ne couvre que \CouvSzGreedy{} des postes cibles. Respecter les contraintes ne signifie donc pas disposer d'un effectif suffisant, et l'outil affiche toujours les deux indicateurs.

\vspace{0.5cm}
\noindent\textbf{Mots-clés :} planification hospitalière, recherche opérationnelle, MILP, CP-SAT, profils de pondération, séries temporelles, données synthétiques, loi 09-08, HL7 FHIR, santé digitale, FastAPI, Next.js.

\vspace{1cm}
\selectlanguage{english}
\chapter*{Abstract}
\addcontentsline{toc}{chapter}{Abstract}

Planning hospital beds and staff is a difficult decision problem: admissions are uncertain, working rules are numerous and wards depend on one another. Carried out within the \textit{Santé digitale et ingénierie des systèmes de santé} module (digital health and health systems engineering), this project presents the design, implementation and evaluation of a reproducible decision-support tool for a synthetic hospital with \NbLits{} beds across four departments (Emergency, Medicine, Intensive Care, Surgery) and \NbAgents{} professionals, over a 7-day horizon ($D+1$ to $D+7$). No real patient data is used.

The pipeline includes synthetic data covering \NbJours{} days; five forecasting models with chronological validation (validation WAPE of \WapeValUrg{} for Emergency with the selected lagged regression); a single formulation of bed opening and individual staff assignment with an \ReposMin{}-hour minimum rest, at most \PostesMax{} shifts per week and minimum staffing; four methods (naive strategy, greedy heuristic, MILP/CBC, CP-SAT) and three weighting profiles (balanced, admission priority, staff priority), that is 72 plans checked by an independent validator; six scenarios and an \NbCas{}-case sensitivity grid; a FastAPI service and a ten-view Next.js dashboard with a custom scenario simulator; and a regulatory, normative and organisational framework (Moroccan law 09-08, HL7 FHIR, ISO/IEC 27001 and ISO 27799).

All 72 plans satisfy every hard constraint. CP-SAT is the best method, or tied for best, in \NbCpMeilleur{} of 18 cases. The main finding is a trade-off between admitting patients and protecting staff: lowering the staffing deficit leads to more rejected admissions. The hospital under study is structurally understaffed: in the nominal scenario the greedy plan accepts \CouvDemSzGreedy{} of the demand but covers only \CouvSzGreedy{} of the target staff positions. Satisfying the constraints is therefore not the same as having enough staff, and the tool always reports both indicators.

\vspace{0.5cm}
\noindent\textbf{Keywords:} hospital resource planning, operations research, mixed-integer linear programming, CP-SAT, weighting profiles, time-series forecasting, synthetic data, law 09-08, HL7 FHIR, digital health, FastAPI, Next.js.
\selectlanguage{french}


\clearpage
\tableofcontents

\clearpage
\listoffigures

\clearpage
\listoftables

\clearpage
\chapter*{Liste des Acronymes}
\addcontentsline{toc}{chapter}{Liste des Acronymes}

\begin{description}
    \item[API] Application Programming Interface
    \item[CBC] COIN-OR Branch and Cut (solveur libre d'optimisation en nombres entiers)
    \item[CHI] Service de Chirurgie
    \item[CNDP] Commission nationale de contrôle de la protection des données à caractère personnel
    \item[CORS] Cross-Origin Resource Sharing
    \item[CP-SAT] Constraint Programming -- Satisfiability (solveur Google OR-Tools)
    \item[DMS] Durée moyenne de séjour
    \item[ENSAM] École Nationale Supérieure d'Arts et Métiers
    \item[FHIR] Fast Healthcare Interoperability Resources (standard HL7)
    \item[HL7] Health Level Seven International
    \item[INDIA] Ingénierie Numérique en Data Science et Intelligence Artificielle
    \item[ISO] Organisation internationale de normalisation
    \item[KPI] Key Performance Indicator (indicateur clé de performance)
    \item[MAE] Mean Absolute Error (erreur absolue moyenne)
    \item[MED] Service de Médecine
    \item[MILP] Mixed-Integer Linear Programming (programmation linéaire en nombres entiers)
    \item[OBS] Organizational Breakdown Structure (structure organisationnelle)
    \item[RACI] Réalisateur, approbateur, consulté, informé
    \item[REA] Service de Réanimation
    \item[RH] Ressources humaines
    \item[RMSE] Root Mean Squared Error (racine de l'erreur quadratique moyenne)
    \item[SIH] Système d'information hospitalier
    \item[URG] Service des Urgences
    \item[WAPE] Weighted Absolute Percentage Error (erreur absolue pondérée en pourcentage)
    \item[WBS] Work Breakdown Structure (décomposition des travaux)
\end{description}


% Corps du rapport de projet (6 chapitres conformes au modèle ENSAM)
\clearpage
\pagenumbering{arabic}
\chapter*{Introduction Générale}
\addcontentsline{toc}{chapter}{Introduction Générale}

\section*{Le Contexte de Tension Hospitalière}
L'hôpital accueille des patients à toute heure, avec des ressources limitées. Sa planification est soumise à des contraintes qui se contredisent souvent :
\begin{itemize}
    \item \textbf{Une demande variable} : une partie des admissions, en particulier aux urgences, n'est pas programmée et dépend du calendrier et d'événements imprévus ;
    \item \textbf{Une capacité physique fixe} : le nombre de lits d'un service ne s'ajuste pas à court terme, et un patient admis occupe un lit pendant plusieurs jours ;
    \item \textbf{Un personnel limité} : chaque lit ouvert demande des médecins, des infirmiers et des aides-soignants qualifiés, jour et nuit, dans le respect des temps de repos ;
    \item \textbf{Des aléas} : absences, pics d'activité et fermetures de lits obligent à replanifier.
\end{itemize}

\section*{Une Planification Souvent Cloisonnée}
Les travaux de recherche opérationnelle distinguent deux familles de décisions~\cite{hulshof2012} : la gestion des lits~\cite{hallbeds} et l'affectation du personnel~\cite{burke2004}. Les traiter séparément pose une difficulté simple : un lit ouvert sans personnel ne permet pas de soigner, et un personnel affecté sans lit disponible est mal employé. Le projet part de cette observation et traite les deux décisions dans un même modèle, à partir d'une prévision de la demande.

\section*{Objectifs et Positionnement du Projet}
Réalisé au sein de la filière INDIA de l'ENSAM Rabat, dans le cadre du module \textbf{Santé Digitale}, ce projet conçoit, implémente et évalue un \textbf{prototype d'aide à la décision pour la planification des lits et du personnel}. Il doit :
\begin{enumerate}
    \item \textbf{Prévoir} les admissions par service de $J+1$ à $J+7$ ;
    \item \textbf{Traduire} cette prévision en besoins de lits et de personnel, en suivant les séjours des patients ;
    \item \textbf{Planifier conjointement} l'ouverture des lits et les postes individuels, avec repos minimal, qualifications et nombre maximal de postes ;
    \item \textbf{Comparer} quatre méthodes (naïve, Greedy, MILP/CBC, CP-SAT) et trois profils de pondération, en contrôlant chaque plan par un validateur indépendant ;
    \item \textbf{Évaluer} la réponse de l'hôpital à six scénarios, dont des cas dégradés (pic d'activité, absences, crise combinée, fermeture de lits), et sa sensibilité sur \NbCas{} cas ;
    \item \textbf{Restituer} les résultats dans une API et un tableau de bord de dix vues, avec un simulateur de scénario personnalisé ;
    \item \textbf{Situer} l'outil dans son cadre réglementaire, normatif et organisationnel.
\end{enumerate}

\begin{notebox}{Nature des données}
Toutes les données sont \textbf{synthétiques} : aucune donnée de patient ou de professionnel réel n'est utilisée. Les effectifs, cibles, durées de séjour et poids sont des hypothèses de démonstration, pas des normes cliniques ou réglementaires.
\end{notebox}

\section*{Organisation du Rapport de Projet}
Le rapport comprend six chapitres, une conclusion générale, une bibliographie et des annexes :
\begin{itemize}
    \item \textbf{Chapitre 1 -- Contexte et analyse des besoins :} hôpital étudié, problématique, besoins et organisation du projet ;
    \item \textbf{Chapitre 2 -- Cadre réglementaire, normatif et organisationnel :} loi 09-08, CNDP, HL7 FHIR, ISO/IEC 27001 et ISO 27799, utilisateurs et processus de décision ;
    \item \textbf{Chapitre 3 -- Étude technique et choix technologiques :} méthodes d'optimisation, modèles de prévision et outils logiciels ;
    \item \textbf{Chapitre 4 -- Analyse et conception du système :} modèle mathématique, objectif, profils, validateur et interface ;
    \item \textbf{Chapitre 5 -- Réalisation et implémentation :} code, API, tableau de bord et simulateur ;
    \item \textbf{Chapitre 6 -- Tests, validation et analyse des résultats :} tests, prévisions, comparaison des méthodes, arbitrages, sensibilité et limites.
\end{itemize}

\chapter{Contexte et Analyse des Besoins}
\label{ch:contexte_analyse_besoins}

% ------------------------------------------------------------------------------
\section*{Introduction}
\addcontentsline{toc}{section}{Introduction}
% ------------------------------------------------------------------------------
Planifier les lits et le personnel d'un hôpital demande d'anticiper la demande, de respecter des règles de travail et de faire des choix lorsque les ressources manquent. Un outil d'aide à la décision peut rendre ces choix explicites et comparables. Ce premier chapitre présente le contexte du projet, réalisé au sein de la filière INDIA de l'ENSAM Rabat dans le module Santé digitale et ingénierie des systèmes de santé, puis l'hôpital synthétique étudié. Il formule ensuite la problématique et ses enjeux, recense les besoins fonctionnels et non fonctionnels, et décrit l'organisation de l'équipe (OBS, WBS, RACI et planning prévisionnel).

% ------------------------------------------------------------------------------
\section{Contexte Général de la Santé Digitale et de la Planification Hospitalière}
% ------------------------------------------------------------------------------

\subsection{De la Donnée à la Décision}
La \textbf{santé digitale} associe données, modèles et outils logiciels au service des organisations de soins. Pour la planification des ressources, quatre apports sont recherchés :
\begin{itemize}
    \item \textbf{L'anticipation} : prévoir les admissions avant que la tension n'apparaisse ;
    \item \textbf{La cohérence} : décider ensemble des lits ouverts et du personnel affecté ;
    \item \textbf{La résilience} : mesurer l'effet d'une crise avant qu'elle ne survienne ;
    \item \textbf{La transparence} : montrer ce qu'une décision améliore et ce qu'elle dégrade.
\end{itemize}

\subsection{Deux Familles de Problèmes}
La littérature distingue les décisions de planification en santé selon leur niveau et leur objet~\cite{hulshof2012}. Deux familles concernent ce projet :
\begin{enumerate}
    \item \textbf{La gestion des lits}~\cite{hallbeds} : ouvrir, affecter et libérer des lits en tenant compte de la durée des séjours ;
    \item \textbf{L'affectation du personnel} (\emph{nurse rostering}) : affecter des professionnels à des postes sous contraintes de disponibilité, de qualification et de repos. Ce problème combinatoire fait l'objet de revues régulières~\cite{burke2004,koukaras2026}.
\end{enumerate}

Dans le projet, les besoins en personnel ne sont pas une donnée fixe : ils découlent des lits ouverts, eux-mêmes liés à la demande prévue.

% ------------------------------------------------------------------------------
\section{Présentation du Projet}
% ------------------------------------------------------------------------------

\subsection{Périmètre Hospitalier}
Le prototype modélise un hôpital synthétique de \textbf{\NbLits{} lits} et \textbf{\NbAgents{} professionnels}, répartis sur quatre services (tableau~\ref{tab:services}). Ces paramètres sont lus dans \texttt{configs/hospital.json}.

% Généré par rapport/scripts/generate_tables.py : ne pas modifier à la main.
\begin{table}[H]
    \centering
    \caption{Services de l'hôpital synthétique (\texttt{configs/hospital.json}).}
    \label{tab:services}
    \resizebox{\textwidth}{!}{\begin{tabular}{lcccccc}
        \toprule
        \textbf{Service} & \textbf{Code} & \textbf{Lits} & \textbf{Admissions de base / jour} & \textbf{DMS (jours)} & \textbf{Effet d'événement} & \textbf{Rang Greedy} \\
        \midrule
        Urgences & \texttt{URG} & 75 & 30 & 2 & +10 & 1 \\
        Médecine & \texttt{MED} & 106 & 18 & 5 & +2 & 3 \\
        Réanimation & \texttt{REA} & 55 & 5 & 8 & +3 & 0 \\
        Chirurgie & \texttt{CHI} & 64 & 12 & 4 & +1 & 2 \\
        \midrule
        \textbf{Total} &  & \textbf{300} &  &  &  &  \\
        \bottomrule
    \end{tabular}}
    \par\vspace{3pt}
    {\footnotesize Rang Greedy : ordre de traitement du personnel par l'heuristique (0 = premier servi).}
\end{table}


\begin{itemize}
    \item \textbf{Urgences (\texttt{URG})} : flux le plus élevé, séjour court, hausse d'activité en fin de semaine ;
    \item \textbf{Médecine (\texttt{MED})} : service le plus grand, dont les lits sont presque tous occupés au début de l'horizon (\RecensementMed{} patients présents avant les sorties du premier jour) ;
    \item \textbf{Réanimation (\texttt{REA})} : peu d'admissions, mais un séjour plus long que l'horizon de planification ;
    \item \textbf{Chirurgie (\texttt{CHI})} : activité concentrée sur les jours ouvrés.
\end{itemize}

\subsection{Personnel et Postes de Travail}
Le personnel comprend \NbMedecins{} médecins, \NbInfirmiers{} infirmiers et \NbAides{} aides-soignants (tableau~\ref{tab:effectifs}). L'activité est continue, sur trois postes de 8 heures : matin ($07\text{h}$--$15\text{h}$), soir ($15\text{h}$--$23\text{h}$) et nuit ($23\text{h}$--$07\text{h}$). Chaque professionnel peut faire au plus \PostesMax{} postes par semaine.

% Généré par rapport/scripts/generate_tables.py : ne pas modifier à la main.
\begin{table}[H]
    \centering
    \caption{Effectifs et cibles de personnel (\texttt{configs/hospital.json}).}
    \label{tab:effectifs}
    \small
    \begin{tabular}{lccc}
        \toprule
        \textbf{Corps professionnel} & \textbf{Agents} & \textbf{Cible pour dix lits ouverts et par poste} & \textbf{Minimum par poste} \\
        \midrule
        Médecins & 50 & 1 & 1 \\
        Infirmiers & 110 & 2 & 1 \\
        Aides-soignants & 75 & 1 & 1 \\
        \midrule
        \textbf{Total} & \textbf{235} &  &  \\
        \bottomrule
    \end{tabular}
    \par\vspace{3pt}
    {\footnotesize La cible est identique pour tous les services et pour les trois postes ; le minimum ne s'applique qu'à un service ouvert.}
\end{table}


% ------------------------------------------------------------------------------
\section{Problématique et Enjeux}
% ------------------------------------------------------------------------------

\subsection{Problématique Centrale}
\begin{center}
\fcolorbox{EnsamBlue}{gris}{
\parbox{0.9\textwidth}{
\centering\itshape\bfseries\color{EnsamBlue}
Comment anticiper les admissions et planifier ensemble l'ouverture des lits et les postes du personnel de plusieurs services, en respectant les règles de travail, et comment rendre lisibles les arbitrages entre accueil des patients et charge du personnel en situation de crise ?
}
}
\end{center}

\subsection{Enjeux Opérationnels et Scientifiques}
\begin{itemize}
    \item \textbf{Enjeu prédictif} : prévoir les admissions à court terme sans utiliser d'information future ;
    \item \textbf{Enjeu combinatoire} : résoudre, avec un budget de \TempsSolveur{} secondes par solveur, un modèle de \NbVariables{} variables et \NbContraintes{} contraintes en $S_0$ ;
    \item \textbf{Enjeu de contrôle} : vérifier chaque plan par un code distinct de celui qui l'a produit ;
    \item \textbf{Enjeu d'interprétation} : distinguer un plan qui respecte les contraintes d'un plan correctement doté en personnel.
\end{itemize}

% ------------------------------------------------------------------------------
\section{Analyse des Besoins}
% ------------------------------------------------------------------------------

\subsection{Besoins Fonctionnels}
Le tableau~\ref{tab:besoins_fonctionnels} présente les besoins fonctionnels et la preuve disponible dans le projet.

\begin{table}[H]
    \centering
    \caption{Synthèse des besoins fonctionnels et de leur réalisation.}
    \label{tab:besoins_fonctionnels}
    \resizebox{\textwidth}{!}{
    \begin{tabular}{lp{5.2cm}p{8.6cm}}
        \toprule
        \textbf{Code} & \textbf{Besoin} & \textbf{Réalisation et preuve} \\
        \midrule
        \textbf{BF01} & Données synthétiques reproductibles & \NbJours{} jours, graine fixée, rapport de qualité (tableau~\ref{tab:qualite}). \\
        \textbf{BF02} & Prévision multi-modèles & Cinq modèles, validation chronologique et test final réservé (tableaux~\ref{tab:previsions_validation} et~\ref{tab:previsions_test}). \\
        \textbf{BF03} & Besoins de lits et de personnel & Instance opérationnelle $S_0$ dans \texttt{results/operational\_instances/S0}. \\
        \textbf{BF04} & Méthodes rapides & Stratégie naïve et heuristique Greedy (tableau~\ref{tab:performances}). \\
        \textbf{BF05} & Solveurs & MILP/CBC et CP-SAT, budget de \TempsSolveur{} s et \ThreadsSolveur{} threads. \\
        \textbf{BF06} & Contrôle indépendant & Module \texttt{validator.py}, 72 plans sur 72 sans violation. \\
        \textbf{BF07} & Scénarios et cas dégradés & Six scénarios $S_0$ à $S_5$ (tableau~\ref{tab:scenarios}). \\
        \textbf{BF08} & Profils de décision & Trois profils de pondération (tableau~\ref{tab:profils}). \\
        \textbf{BF09} & Analyse de sensibilité & Grille de \NbCas{} cas, calculée avec Naïve et Greedy. \\
        \textbf{BF10} & Simulation personnalisée & Vue Simulateur et champ \texttt{custom\_scenario} de l'API. \\
        \textbf{BF11} & Restitution & API de \NbEndpoints{} points d'accès et tableau de bord de dix vues. \\
        \textbf{BF12} & Export & Téléchargement contrôlé des artefacts d'une expérience. \\
        \bottomrule
    \end{tabular}
    }
\end{table}

\subsection{Besoins Non Fonctionnels}
\begin{itemize}
    \item \textbf{Réactivité} : les méthodes Naïve et Greedy répondent en moins de deux secondes ; les solveurs ont un budget de temps fixe et configurable ;
    \item \textbf{Reproductibilité} : graine fixée, empreinte des entrées de chaque expérience, chemins relatifs ;
    \item \textbf{Intégrité} : validation des configurations et des échanges par schémas Pydantic, contrôle des chemins téléchargés ;
    \item \textbf{Honnêteté de l'affichage} : une valeur absente donne un état vide explicite, jamais un chiffre inventé.
\end{itemize}

% ------------------------------------------------------------------------------
\section{Organisation et Planification du Projet}
% ------------------------------------------------------------------------------

Le projet a été mené par une équipe de six étudiants de la filière INDIA. Le planning prévisionnel s'étend du 21 septembre au 24 novembre 2026, date prévue de la soutenance.

\subsection{Structure Organisationnelle de l'Équipe (OBS)}
\begin{figure}[H]
    \centering
    \captureimg[0.92\textwidth]{obs}
    \caption{Structure organisationnelle et répartition des domaines techniques (OBS).}
    \label{fig:obs_team}
\end{figure}

\begin{itemize}
    \item \textbf{LACHGAR Ridouan} : données synthétiques, contrôle de qualité et évaluation des cinq modèles de prévision ;
    \item \textbf{CHAOUKI Basma} : modélisation mathématique, contraintes de personnel et modèle MILP (PuLP / CBC) ;
    \item \textbf{EL HOUR Maroua} : modèle CP-SAT (OR-Tools) et validateur indépendant ;
    \item \textbf{EL FEKAK Hafsa} : architecture logicielle, API FastAPI, jobs de calcul et sécurité des téléchargements ;
    \item \textbf{ABDELLAOUI Hajar} : tableau de bord Next.js et graphiques ;
    \item \textbf{AFIF Mohamed} : protocole expérimental, scénarios $S_0$ à $S_5$ et grille de sensibilité.
\end{itemize}

\subsection{Décomposition Hiérarchique du Projet (WBS)}
\begin{figure}[H]
    \centering
    \captureimg{wbs}
    \caption{Décomposition modulaire et hiérarchique des travaux (WBS).}
    \label{fig:wbs_diagram}
\end{figure}

\subsection{Matrice des Responsabilités (RACI)}
Le tableau~\ref{tab:raci_matrix} utilise les conventions \textbf{R} (réalisateur), \textbf{A} (approbateur), \textbf{C} (consulté) et \textbf{I} (informé).

\begin{table}[H]
    \centering
    \caption{Matrice d'attribution des responsabilités (RACI) par lot de travail.}
    \label{tab:raci_matrix}
    \resizebox{\textwidth}{!}{
    \begin{tabular}{lcccccc}
        \toprule
        \textbf{Lot de travail} & \textbf{LACHGAR} & \textbf{EL FEKAK} & \textbf{CHAOUKI} & \textbf{EL HOUR} & \textbf{ABDELLAOUI} & \textbf{AFIF} \\
        \midrule
        WP1 : Cadrage et données & R & C & C & I & I & C \\
        WP2 : Modélisation et prévision & R & I & R & C & I & I \\
        WP3 : Optimisation et contrôle & I & C & R & R & I & C \\
        WP4 : Plateforme web & I & R & I & I & R & I \\
        WP5 : Expérimentations & C & C & C & C & C & R \\
        \bottomrule
    \end{tabular}
    }
\end{table}

\subsection{Planification Temporelle du Projet (Diagramme de Gantt)}
\begin{figure}[H]
    \centering
    \captureimg{gantt}
    \caption{Planning prévisionnel du projet (21 septembre au 24 novembre 2026).}
    \label{fig:gantt_diagram}
\end{figure}

Cinq jalons rythment ce planning : données contrôlées (M1), prévision validée (M2), méthodes et validateur (M3), plateforme intégrée (M4) et soutenance (M5).

% ------------------------------------------------------------------------------
\section*{Conclusion}
\addcontentsline{toc}{section}{Conclusion}
% ------------------------------------------------------------------------------
Ce chapitre a posé le cadre du projet : un hôpital synthétique de \NbLits{} lits et \NbAgents{} professionnels répartis sur quatre services, et une problématique centrée sur l'arbitrage entre accueil des patients et charge du personnel. Douze besoins fonctionnels ont été recensés, chacun associé à une preuve disponible dans le projet, et les travaux ont été organisés en cinq lots. Un tel outil manipulerait toutefois, dans un établissement réel, des données et des décisions encadrées. Avant d'aborder la technique, le chapitre suivant examine donc les règles, les normes et l'organisation qui encadreraient son usage.

\chapter{Cadre Réglementaire, Normatif et Organisationnel}
\label{ch:cadre_reglementaire}

% ------------------------------------------------------------------------------
\section*{Introduction}
\addcontentsline{toc}{section}{Introduction}
% ------------------------------------------------------------------------------
Un outil de planification hospitalière manipulerait, dans un établissement réel, des données relatives à des patients et à des professionnels. Sa mise en service ne serait donc pas seulement une question technique. Ce chapitre situe le prototype par rapport à la loi marocaine 09-08, au standard d'interopérabilité HL7 FHIR et aux normes de sécurité ISO/IEC 27001 et ISO 27799, puis décrit les utilisateurs visés et le processus de décision hebdomadaire. Il distingue ce que le prototype fait déjà de ce qui resterait à faire ; il ne constitue pas un avis juridique.

% ------------------------------------------------------------------------------
\section{Cadre Réglementaire : Loi 09-08 et CNDP}
% ------------------------------------------------------------------------------

\subsection{Les Textes}
La loi n\textsuperscript{o}~09-08 relative à la protection des personnes physiques à l'égard du traitement des données à caractère personnel a été promulguée par le dahir n\textsuperscript{o}~1-09-15 du 18 février 2009 et publiée au Bulletin officiel n\textsuperscript{o}~5714 du 5 mars 2009~\cite{loi0908}. Son décret d'application est le décret n\textsuperscript{o}~2-09-165 du 21 mai 2009. La Commission nationale de contrôle de la protection des données à caractère personnel (CNDP) est instituée par l'article~27.

\subsection{Les Dispositions Utiles au Projet}
\begin{table}[H]
    \centering
    \caption{Articles de la loi 09-08 utiles à un outil de planification hospitalière.}
    \label{tab:loi0908}
    \resizebox{\textwidth}{!}{
    \begin{tabular}{lp{7.6cm}p{7.2cm}}
        \toprule
        \textbf{Article} & \textbf{Contenu} & \textbf{Conséquence pour l'outil} \\
        \midrule
        Art. 1, al. 1 & Donnée à caractère personnel : information concernant une personne physique identifiée ou identifiable. & Les données du prototype ne concernent aucune personne réelle. \\
        Art. 1, al. 3 & Données sensibles : elles incluent les données relatives à la santé. & Des admissions nominatives réelles seraient sensibles. \\
        Art. 3 & Finalités déterminées, données adéquates, pertinentes et non excessives, conservation limitée. & La planification n'a besoin que de comptages par service et par jour. \\
        Art. 12, al. 1-a et 21 & Autorisation préalable pour le traitement de données sensibles. & Démarche auprès de la CNDP avant tout usage de données de santé identifiantes. \\
        Art. 22 & Déclaration, par dérogation, si la seule finalité relève notamment de la gestion des services de santé, par une personne soumise au secret. & Régime à faire confirmer par la CNDP pour l'établissement concerné. \\
        Art. 23 et 24 & Mesures techniques et organisationnelles de sécurité ; contrôles renforcés pour les données sensibles. & Contrôle d'accès, des supports et de la transmission à mettre en place. \\
        Art. 26 & Secret professionnel des personnes ayant connaissance des données. & Engagement des utilisateurs de l'outil. \\
        \bottomrule
    \end{tabular}
    }
\end{table}

\subsection{Justification des Données Synthétiques}
Le choix de données \textbf{entièrement synthétiques} répond à une exigence du cahier des charges et à trois arguments :
\begin{enumerate}
    \item \textbf{Absence de personne concernée} : les admissions, l'occupation et les professionnels sont produits par un générateur (\texttt{generator.py}) à partir de paramètres et d'une graine, sans fichier source réel. Selon notre lecture de l'article~1, il ne s'agit pas de données à caractère personnel ;
    \item \textbf{Minimisation} : même en données réelles, le modèle n'utiliserait que des comptages d'admissions et de patients présents par service, jamais un dossier médical ;
    \item \textbf{Reproductibilité} : le jeu de données peut être publié et régénéré à l'identique, ce que des données de santé ne permettraient pas.
\end{enumerate}

La limite de ce choix est connue : les résultats ne décrivent aucun hôpital existant. Les identifiants de professionnels (par exemple \texttt{NUR-001}) sont eux aussi fictifs. Dans un usage réel, le planning du personnel serait nominatif et relèverait de la loi, même sans donnée de santé.

% ------------------------------------------------------------------------------
\section{Règles de Travail Retenues}
% ------------------------------------------------------------------------------

Le modèle impose un repos minimal de \ReposMin{}\,h entre deux postes, au plus \PostesMax{} postes par semaine et au plus \NuitsMax{} nuits consécutives. Ces valeurs sont des \textbf{paramètres de configuration}. Le repos de 11 heures consécutives par période de 24 heures correspond à l'article~3 de la directive européenne 2003/88/CE~\cite{directive2003}, citée comme référence connue ; le prototype ne prétend pas vérifier la conformité au droit du travail marocain ni au statut du personnel hospitalier.

% ------------------------------------------------------------------------------
\section{Cadre Normatif : Interopérabilité HL7 FHIR}
% ------------------------------------------------------------------------------

Le prototype lit des fichiers CSV et JSON. Une intégration au système d'information hospitalier (SIH) pourrait s'appuyer sur HL7 FHIR, dans sa version R4 (4.0.1)~\cite{fhir}. Le tableau~\ref{tab:fhir} propose une correspondance ; \textbf{elle n'est pas implémentée}.

\begin{table}[H]
    \centering
    \caption{Correspondance proposée entre le modèle et les ressources FHIR R4.}
    \label{tab:fhir}
    \resizebox{\textwidth}{!}{
    \begin{tabular}{lp{6.6cm}p{7.8cm}}
        \toprule
        \textbf{Ressource FHIR} & \textbf{Définition et éléments utiles} & \textbf{Usage envisagé} \\
        \midrule
        \texttt{Location} & Lieu de soins ; \texttt{physicalType} et \texttt{operationalStatus}, ce dernier visant surtout les lits et les chambres. & Capacité disponible et lits fermés par service (\texttt{capacity\_available}). \\
        \texttt{Encounter} & Interaction entre un patient et un prestataire de soins ; \texttt{status}, \texttt{period}, \texttt{location}, \texttt{hospitalization}. & Admissions, sorties et recensement initial, agrégés par service avant d'entrer dans le modèle. \\
        \texttt{Practitioner} & Personne participant aux soins ; \texttt{qualification}. & Rôles et qualifications du personnel. \\
        \texttt{Schedule} & Ensemble de créneaux pour un ou plusieurs acteurs ; \texttt{actor}, \texttt{planningHorizon}. & Disponibilités en entrée et plan hebdomadaire en sortie. \\
        \bottomrule
    \end{tabular}
    }
\end{table}

L'agrégation des ressources \texttt{Encounter} en comptages devrait avoir lieu dans le SIH, avant l'outil de planification, afin de ne pas y transférer de donnée identifiante.

% ------------------------------------------------------------------------------
\section{Cadre Normatif : Sécurité de l'Information}
% ------------------------------------------------------------------------------

ISO/IEC 27001:2022 définit les exigences d'un système de management de la sécurité de l'information~\cite{iso27001}. ISO 27799:2025 décline, pour la santé, des mesures de sécurité fondées sur ISO/IEC 27002~\cite{iso27799}. Le prototype n'est \textbf{pas certifié} et n'a pas fait l'objet d'un audit. Le tableau~\ref{tab:iso} situe son état par thème.

\begin{table}[H]
    \centering
    \caption{État du prototype par thème de sécurité.}
    \label{tab:iso}
    \resizebox{\textwidth}{!}{
    \begin{tabular}{lp{7.4cm}p{7.2cm}}
        \toprule
        \textbf{Thème} & \textbf{Présent dans le prototype} & \textbf{À faire avant un usage réel} \\
        \midrule
        Protection des données & Données synthétiques uniquement. & Registre des traitements, démarche CNDP, durée de conservation. \\
        Contrôle d'accès & Aucun : API locale sans authentification. & Authentification, rôles et droits par service. \\
        Échanges & Origines CORS explicites, le joker est refusé. & Chiffrement des échanges et des supports. \\
        Intégrité & Schémas Pydantic, empreinte SHA-256 des entrées de scénario, validateur indépendant. & Signature ou horodatage des plans validés. \\
        Fichiers & Contrôle du chemin résolu avant tout téléchargement. & Cloisonnement des droits sur les dossiers. \\
        Traçabilité & Graine, profil, poids, budget et date de chaque expérience. & Journal des accès et des décisions humaines. \\
        Disponibilité & Résultats conservés sur disque ; les jobs en mémoire sont perdus au redémarrage. & Sauvegardes, reprise et supervision. \\
        \bottomrule
    \end{tabular}
    }
\end{table}

% ------------------------------------------------------------------------------
\section{Volet Organisationnel}
% ------------------------------------------------------------------------------

Cette section est une \textbf{proposition de l'équipe} : elle n'a pas été testée auprès d'un établissement, et l'utilisabilité du tableau de bord n'a pas été mesurée.

\subsection{Utilisateurs et Rôles}
\begin{table}[H]
    \centering
    \caption{Utilisateurs visés et vues du tableau de bord associées.}
    \label{tab:utilisateurs}
    \resizebox{\textwidth}{!}{
    \begin{tabular}{lp{6.8cm}p{6.6cm}}
        \toprule
        \textbf{Rôle} & \textbf{Responsabilité} & \textbf{Vues principales} \\
        \midrule
        Gestionnaire des lits & Ouverture et fermeture de lits, refus et réorientations. & Vue d'ensemble, Allocation, Simulateur \\
        Cadre de santé & Planning du personnel et remplacements. & Allocation (couverture et planning) \\
        Direction des soins & Choix du profil et arbitrage en cas de crise. & Comparaison scénarios, Sensibilité \\
        Référent données et sécurité & Qualité des entrées, conformité et accès. & Exploration, Expériences et export \\
        \bottomrule
    \end{tabular}
    }
\end{table}

\subsection{Processus de Décision Hebdomadaire}
\begin{enumerate}
    \item \textbf{Actualiser} les données et les prévisions de $J+1$ à $J+7$ ;
    \item \textbf{Calculer} un premier plan avec Naïve et Greedy, en quelques secondes ;
    \item \textbf{Comparer} avec MILP/CBC et CP-SAT, pour un profil choisi par la direction ;
    \item \textbf{Tester} dans le simulateur les aléas redoutés : absences, hausse de demande, fermeture de lits ;
    \item \textbf{Lire ensemble} les refus, la couverture de la demande, le déficit RH et les postes couverts ;
    \item \textbf{Décider} : le plan est une proposition ; les responsables valident, corrigent ou rejettent ;
    \item \textbf{Consigner} l'expérience retenue, son profil et les écarts décidés.
\end{enumerate}

\begin{notebox}{Place de l'humain}
L'outil ne décide pas d'une prise en charge. Un « refus » du modèle signifie qu'une admission prévue ne tient pas dans le plan ; sa traduction concrète (réorientation, report, renfort) relève des professionnels.
\end{notebox}

\subsection{Conduite en Mode Dégradé}
Trois situations dégradées sont prévues et mesurées au chapitre~\ref{ch:tests_validation_resultats} : saturation ($S_2$, $S_4$), absences ($S_3$, $S_4$) et fermeture de lits ($S_5$). Côté logiciel, si un solveur ne trouve pas de solution dans son budget, son statut est \texttt{UNKNOWN} et aucun indicateur n'est affiché ; si l'API ne répond pas, le tableau de bord le signale au lieu d'afficher des valeurs anciennes. Les méthodes Naïve et Greedy servent alors de repli immédiat.

% ------------------------------------------------------------------------------
\section*{Conclusion}
\addcontentsline{toc}{section}{Conclusion}
% ------------------------------------------------------------------------------
Les données synthétiques permettent de développer et d'évaluer l'outil sans traiter de donnée personnelle. Un usage réel demanderait une démarche auprès de la CNDP, une interopérabilité avec le système d'information hospitalier, pour laquelle une correspondance FHIR est proposée, et des mesures de sécurité que le prototype n'a pas, en premier lieu l'authentification. L'organisation proposée, non encore testée, associe quatre rôles à un processus hebdomadaire et laisse la décision aux professionnels. Ce cadre étant posé, le chapitre suivant présente les méthodes et les outils techniques retenus.

\chapter{Étude Technique et Choix Technologiques}
\label{ch:etude_technique}

% ------------------------------------------------------------------------------
\section*{Introduction}
\addcontentsline{toc}{section}{Introduction}
% ------------------------------------------------------------------------------
La conception d'un système d'aide à la décision pour la planification hospitalière assemble plusieurs briques : modèles de prévision, méthodes d'optimisation, service web et interface de visualisation. Ce chapitre présente les critères de sélection, puis les quatre méthodes de planification et les cinq modèles de prévision comparés. Il décrit ensuite chaque outil retenu et la raison de son choix au regard des besoins identifiés au chapitre~\ref{ch:contexte_analyse_besoins}, et les résume dans un tableau de synthèse. Les solutions non retenues sont citées comme alternatives ; elles n'ont pas fait l'objet de mesures comparatives dans le projet.

% ------------------------------------------------------------------------------
\section{Critères de Sélection}
% ------------------------------------------------------------------------------

Les technologies ont été choisies en fonction de cinq critères, déduits des besoins du projet :
\begin{itemize}
    \item \textbf{Adéquation au problème :} prévision de séries journalières, optimisation combinatoire sous contraintes et restitution interactive ;
    \item \textbf{Exécution locale :} fonctionnement sur un poste de travail, sans serveur dédié ni base de données ;
    \item \textbf{Facilité d'intégration :} bibliothèques documentées, utilisables ensemble dans un même dépôt ;
    \item \textbf{Gratuité :} outils utilisables sans licence à acquérir ;
    \item \textbf{Contrôle :} possibilité de valider les données échangées et de tester automatiquement le calcul et l'interface.
\end{itemize}

% ------------------------------------------------------------------------------
\section{Méthodes et Modèles Comparés}
% ------------------------------------------------------------------------------

\subsection{Méthodes d'Optimisation}
Le projet compare quatre méthodes sur les mêmes instances : deux heuristiques rapides et deux solveurs. Cette comparaison permet de mesurer ce qu'apporte un solveur par rapport à une règle simple.
\begin{itemize}
    \item \textbf{Stratégie naïve} : point de comparaison bas. Chaque jour, un taux d'acceptation unique s'applique à tous les services ; le personnel est partagé proportionnellement aux cibles, sans service prioritaire. Elle est déterministe et rapide, mais ne redistribue pas les postes qu'un service ne peut pas pourvoir ;
    \item \textbf{Heuristique gloutonne (\emph{Greedy})} : référence de la comparaison. Elle accepte autant d'admissions que de lits libres, puis affecte le personnel jour par jour, dans l'ordre $\text{REA} \succ \text{URG} \succ \text{CHI} \succ \text{MED}$, en deux passes : dotation minimale, puis cible. Elle est déterministe et calculée en moins de deux secondes, mais ses choix ne sont jamais remis en cause ;
    \item \textbf{Programmation linéaire en nombres entiers (MILP)} : le problème est décrit par des variables entières et binaires et par des contraintes linéaires~\cite{wolsey2020}. Le solveur fournit un \emph{gap}, écart entre la solution trouvée et une borne sur l'optimum ;
    \item \textbf{Programmation par contraintes (CP-SAT)} : elle propage des règles logiques sur des variables entières~\cite{rossi2006}. Avec plusieurs threads et un budget de temps, deux exécutions peuvent donner deux plans différents.
\end{itemize}

\begin{notebox}{« Solveur exact » ne signifie pas « optimum trouvé »}
CBC et CP-SAT sont capables de prouver l'optimalité, mais ils sont arrêtés après \TempsSolveur{} secondes. Sauf preuve terminée dans ce budget, leur résultat est une solution réalisable accompagnée d'un gap. Le statut \texttt{OPTIMAL} n'est accordé que si la recherche est allée à son terme.
\end{notebox}

\subsection{Modèles de Prévision}
Les admissions journalières de chaque service sont prévues de $J+1$ à $J+7$. Cinq modèles sont comparés~\cite{hyndman2021} :
\begin{enumerate}
    \item \textbf{Modèle naïf} : répète la dernière valeur observée ;
    \item \textbf{Modèle naïf saisonnier} : répète la valeur du même jour de la semaine précédente (période 7) ;
    \item \textbf{Moyenne mobile} : moyenne des 7 derniers jours observés ;
    \item \textbf{Lissage exponentiel simple} : moyenne à poids décroissants dans le temps ;
    \item \textbf{Régression à retards} : régression de type Ridge à coefficients positifs, utilisant les valeurs à 1, 7 et 14 jours, les moyennes des 7 et 28 derniers jours, le jour de la semaine et les jours fériés synthétiques. Les événements exceptionnels sont exclus, car ils ne seraient pas connus au moment de prévoir.
\end{enumerate}

La \textbf{régression à retards} est retenue pour les quatre services, car elle a la plus faible erreur WAPE de validation (tableau~\ref{tab:previsions_validation}). Le chapitre~\ref{ch:tests_validation_resultats} montre que cet avantage n'est pas confirmé partout sur le test final.

% ------------------------------------------------------------------------------
\section{Étude et Choix des Outils et Technologies}
% ------------------------------------------------------------------------------

% ------------------------------------------------------------------------------
\subsection{Langage de Programmation : Python}
% ------------------------------------------------------------------------------

\textbf{Python} est un langage polyvalent, largement utilisé en traitement des données et en optimisation. Sa syntaxe claire et ses nombreuses bibliothèques facilitent le développement d'une chaîne de calcul complète.

\begin{figure}[H]
    \centering
    \includegraphics[width=0.2\textwidth]{logos/python.png}
    \caption{Logo de Python.}
    \label{fig:logo_python}
\end{figure}

\paragraph{Justification du choix.} Python a été retenu pour toute la partie calcul et pour l'API, car :
\begin{itemize}
    \item il réunit dans un même langage la préparation des données, la prévision, les deux solveurs et le service web ;
    \item les calculs lourds sont confiés à des solveurs compilés, ce qui limite l'intérêt d'un langage compilé pour le reste du code ;
    \item les annotations de type sont exploitées par Pydantic et FastAPI pour contrôler les données.
\end{itemize}

% ------------------------------------------------------------------------------
\subsection{Manipulation des Données : pandas}
% ------------------------------------------------------------------------------

\textbf{pandas} est une bibliothèque Python de manipulation de données tabulaires : lecture et écriture de fichiers CSV, jointures, agrégations et gestion des dates.

\begin{figure}[H]
    \centering
    \includegraphics[width=0.2\textwidth]{logos/pandas.png}
    \caption{Logo de pandas.}
    \label{fig:logo_pandas}
\end{figure}

\paragraph{Justification du choix.} Les données du projet sont des tables datées par service, par jour et par poste, dont la plus grande compte \LignesDispo{} lignes. pandas permet de les contrôler et de les joindre sans base de données, et de conserver des résultats dans des fichiers CSV lisibles. Le validateur l'utilise aussi pour recalculer les règles à partir des tables produites.

% ------------------------------------------------------------------------------
\subsection{Apprentissage Automatique : scikit-learn}
% ------------------------------------------------------------------------------

\textbf{scikit-learn} est une bibliothèque Python d'apprentissage automatique, qui fournit des modèles de régression et de classification sous une interface commune.

\begin{figure}[H]
    \centering
    \includegraphics[height=2.5cm]{logos/scikit_learn.png}
    \caption{Logo de scikit-learn.}
    \label{fig:logo_scikit_learn}
\end{figure}

\paragraph{Justification du choix.} La régression à retards utilise sa classe \texttt{Ridge}, dont la régularisation convient à des variables corrélées comme les retards et les moyennes mobiles, et dont l'option \texttt{positive} interdit les coefficients négatifs. Le modèle reste rapide et explicable ; un réseau de neurones n'a pas été jugé nécessaire pour ces séries courtes.

% ------------------------------------------------------------------------------
\subsection{Modèles Statistiques : statsmodels}
% ------------------------------------------------------------------------------

\textbf{statsmodels} est une bibliothèque Python de modélisation statistique, qui comprend des modèles de séries temporelles.

\begin{figure}[H]
    \centering
    \includegraphics[width=0.2\textwidth]{logos/statsmodels.png}
    \caption{Logo de statsmodels.}
    \label{fig:logo_statsmodels}
\end{figure}

\paragraph{Justification du choix.} Elle fournit le lissage exponentiel simple (\texttt{SimpleExpSmoothing}), utilisé comme modèle statistique de comparaison. S'appuyer sur une implémentation existante évite de recoder l'estimation du paramètre de lissage.

% ------------------------------------------------------------------------------
\subsection{Programmation Linéaire en Nombres Entiers : PuLP et COIN-OR CBC}
% ------------------------------------------------------------------------------

\textbf{PuLP} est une bibliothèque Python de modélisation de problèmes linéaires~\cite{mitchell2011}. \textbf{CBC} est le solveur libre en nombres entiers auquel elle transmet le modèle~\cite{cbc}. Tous deux sont hébergés par le projet COIN-OR.

\begin{figure}[H]
    \centering
    \includegraphics[height=2.5cm]{logos/coin_or.png}
    \caption{Logo de COIN-OR, projet qui héberge PuLP et CBC.}
    \label{fig:logo_coin_or}
\end{figure}

\paragraph{Justification du choix.} PuLP permet d'écrire les contraintes sous une forme proche des équations du chapitre~\ref{ch:conception_systeme}, et l'installation fournit directement CBC. Le solveur donne une solution et une borne, donc un gap, sans licence commerciale. Ses limites sur ce problème, notamment avec la priorité personnel, sont mesurées au chapitre~\ref{ch:tests_validation_resultats}.

% ------------------------------------------------------------------------------
\subsection{Programmation par Contraintes : Google OR-Tools}
% ------------------------------------------------------------------------------

\textbf{Google OR-Tools} est une suite libre d'outils d'optimisation. Son solveur \textbf{CP-SAT} traite des problèmes à variables entières et binaires~\cite{ortools}.

\begin{figure}[H]
    \centering
    \includegraphics[height=2.5cm]{logos/ortools.png}
    \caption{Logo de Google OR-Tools.}
    \label{fig:logo_ortools}
\end{figure}

\paragraph{Justification du choix.} Le modèle est dominé par des variables binaires d'affectation et par des incompatibilités entre postes, auxquelles CP-SAT est adapté. Il exploite plusieurs threads (\ThreadsSolveur{} dans le projet), accepte le plan Greedy comme indication de départ et fournit une borne inférieure. Comparer CBC et CP-SAT sur la même formulation permet en outre de contrôler l'un par l'autre.

% ------------------------------------------------------------------------------
\subsection{API REST : FastAPI}
% ------------------------------------------------------------------------------

\textbf{FastAPI} est un cadre Python de développement d'API fondé sur les annotations de type~\cite{fastapi}. Il est exécuté ici par le serveur \textbf{Uvicorn}.

\begin{figure}[H]
    \centering
    \includegraphics[width=0.2\textwidth]{logos/fastapi.png}
    \caption{Logo de FastAPI.}
    \label{fig:logo_fastapi}
\end{figure}

\paragraph{Justification du choix.} Django est complet mais orienté vers des applications à base de données, absentes du projet ; Flask est léger, mais la validation des échanges y demande des extensions. FastAPI a été retenu car :
\begin{itemize}
    \item les requêtes et les réponses sont contrôlées par des schémas partagés avec le moteur de calcul ;
    \item un calcul long est soumis comme un \emph{job} et répond immédiatement ; un \texttt{ThreadPoolExecutor} à un seul travailleur évite que deux solveurs ne se disputent le processeur ;
    \item la documentation interactive (\texttt{/docs}) est générée depuis la description OpenAPI.
\end{itemize}

% ------------------------------------------------------------------------------
\subsection{Validation des Données : Pydantic}
% ------------------------------------------------------------------------------

\textbf{Pydantic} est une bibliothèque Python qui décrit des structures de données et vérifie leur conformité à l'aide des annotations de type~\cite{pydantic}.

\begin{figure}[H]
    \centering
    \includegraphics[height=2.5cm]{logos/pydantic.png}
    \caption{Logo de Pydantic.}
    \label{fig:logo_pydantic}
\end{figure}

\paragraph{Justification du choix.} Les mêmes classes (\texttt{contracts.py}) décrivent la configuration de l'hôpital, les profils, les résultats et les requêtes de l'API. Une valeur hors bornes, comme un taux d'absence supérieur à 50\,\% dans le simulateur, ou un champ inconnu est refusé avec un message qui désigne le champ fautif, sans contrôle écrit à la main.

% ------------------------------------------------------------------------------
\subsection{Structure de l'Interface Web : Next.js}
% ------------------------------------------------------------------------------

\textbf{Next.js} est un cadre de développement web fondé sur React, qui fournit la structure du projet, le serveur de développement et les outils de compilation~\cite{nextjs}.

\begin{figure}[H]
    \centering
    \includegraphics[width=0.2\textwidth]{logos/nextjs.png}
    \caption{Logo de Next.js.}
    \label{fig:logo_nextjs}
\end{figure}

\paragraph{Justification du choix.} Le tableau de bord ne se limite pas à afficher des courbes : il lance des simulations, suit leur avancement et croise plusieurs filtres. Ces interactions justifient une application dédiée. Next.js apporte une organisation de projet et un outillage TypeScript prêts à l'emploi ; la mise en page utilise \textbf{Tailwind CSS}~\cite{tailwindcss} et une feuille de style propre au projet.

% ------------------------------------------------------------------------------
\subsection{Composants d'Interface : React}
% ------------------------------------------------------------------------------

\textbf{React} est une bibliothèque JavaScript qui permet de construire une interface à partir de composants réutilisables~\cite{react}.

\begin{figure}[H]
    \centering
    \includegraphics[height=2.5cm]{logos/react.png}
    \caption{Logo de React.}
    \label{fig:logo_react}
\end{figure}

\paragraph{Justification du choix.} Les dix vues partagent les mêmes composants et les mêmes filtres. Les courbes, les barres et les cartes de chaleur sont écrites directement en composants React produisant du SVG ou des tableaux, sans bibliothèque graphique, ce qui permet de contrôler l'affichage des valeurs absentes et des plans irréalisables.

% ------------------------------------------------------------------------------
\subsection{Typage de l'Interface : TypeScript}
% ------------------------------------------------------------------------------

\textbf{TypeScript} ajoute un typage statique à JavaScript.

\begin{figure}[H]
    \centering
    \includegraphics[height=2.5cm]{logos/typescript.png}
    \caption{Logo de TypeScript.}
    \label{fig:logo_typescript}
\end{figure}

\paragraph{Justification du choix.} Les réponses de l'API sont décrites dans \texttt{api-types.ts} et utilisées par toutes les vues. Une incohérence de format, ou l'oubli d'une valeur possiblement absente, est signalée dès la vérification des types, avant l'exécution.

% ------------------------------------------------------------------------------
\subsection{Tests du Calcul et de l'API : pytest}
% ------------------------------------------------------------------------------

\textbf{pytest} est un outil de test pour Python.

\begin{figure}[H]
    \centering
    \includegraphics[height=2.5cm]{logos/pytest.png}
    \caption{Logo de pytest.}
    \label{fig:logo_pytest}
\end{figure}

\paragraph{Justification du choix.} Les \NbTestsPython{} tests couvrent les données, les prévisions, les méthodes, le validateur, les scénarios et l'API. Les \emph{fixtures} de pytest permettent de partager la configuration de référence entre les tests, et les solveurs sont vérifiés sur de petites instances. L'analyse statique est assurée par \textbf{ruff}.

% ------------------------------------------------------------------------------
\subsection{Tests de l'Interface : Playwright}
% ------------------------------------------------------------------------------

\textbf{Playwright} est un outil qui pilote un navigateur pour tester une application web.

\begin{figure}[H]
    \centering
    \includegraphics[height=2.5cm]{logos/playwright.png}
    \caption{Logo de Playwright.}
    \label{fig:logo_playwright}
\end{figure}

\paragraph{Justification du choix.} Les \NbTestsNavigateur{} parcours vérifient l'ensemble, du navigateur à l'API : ouverture des dix vues sans calcul automatique, simulation complète, cas irréalisable et changement de profil. Ils s'exécutent sur une copie des résultats. Playwright a aussi servi à produire les captures de ce rapport.

% ------------------------------------------------------------------------------
\subsection{Gestion de Versions : Git et GitHub}
% ------------------------------------------------------------------------------

\textbf{Git} est un système de gestion de versions. \textbf{GitHub} héberge le dépôt du projet.

\begin{figure}[H]
    \centering
    \includegraphics[width=0.2\textwidth]{logos/git.png}
    \hspace{1.5cm}
    \includegraphics[height=2.5cm]{logos/github.png}
    \caption{Logos de Git et de GitHub.}
    \label{fig:logo_git_github}
\end{figure}

\paragraph{Justification du choix.} Le code, les configurations, les données synthétiques, les résultats de référence et ce rapport sont versionnés ensemble : un résultat peut ainsi être rattaché à l'état du code qui l'a produit. Les simulations personnalisées, propres à chaque poste, sont exclues du dépôt.

% ------------------------------------------------------------------------------
\section{Synthèse des Choix Technologiques}
% ------------------------------------------------------------------------------

Le tableau~\ref{tab:technologies} récapitule, pour chaque fonction du système, la technologie retenue et la raison principale de son choix. Les versions sont celles installées lors de la rédaction.

\begin{table}[H]
\centering
\small
\renewcommand{\arraystretch}{1.25}
\begin{tabularx}{\linewidth}{>{\raggedright\arraybackslash}p{3.2cm} >{\raggedright\arraybackslash}p{4.6cm} >{\raggedright\arraybackslash}X}
\toprule
\textbf{Fonction} & \textbf{Technologie retenue} & \textbf{Raison principale du choix} \\
\midrule
Calcul et API & Python 3.13.2 & Un seul langage pour les données, la prévision, l'optimisation et le service \\
Données tabulaires & pandas 2.3.3 & Tables datées, fichiers CSV lisibles, sans base de données \\
Régression & scikit-learn 1.9.1 & Ridge à coefficients positifs, rapide et explicable \\
Lissage exponentiel & statsmodels 0.15.0 & Implémentation statistique existante \\
Solveur MILP & PuLP 4.3.1 / COIN-OR CBC & Formulation lisible, gap fourni, sans licence \\
Solveur par contraintes & Google OR-Tools CP-SAT 9.15 & Variables binaires, recherche parallèle, borne inférieure \\
API REST & FastAPI 0.165.0 / Uvicorn 0.54.0 & Contrats typés, documentation OpenAPI \\
Validation & Pydantic 2.14.4 & Schémas partagés, erreurs détaillées \\
Interface & Next.js 15.5.28 / Tailwind CSS 3.4.19 & Structure de projet et outillage prêts à l'emploi \\
Composants et graphiques & React 19.3.1 (SVG) & Composants partagés, contrôle des états vides \\
Typage & TypeScript 5.9.3 & Erreurs de format détectées avant l'exécution \\
Tests Python & pytest 9.2.0 & Contrôle du calcul, du validateur et de l'API \\
Tests navigateur & Playwright 1.64.1 & Parcours complets, du navigateur à l'API \\
Versions & Git / GitHub & Code, données et résultats versionnés ensemble \\
\bottomrule
\end{tabularx}
\caption{Synthèse et justification des technologies retenues.}
\label{tab:technologies}
\end{table}

% ------------------------------------------------------------------------------
\section*{Conclusion}
\addcontentsline{toc}{section}{Conclusion}
% ------------------------------------------------------------------------------
Les outils retenus sont gratuits et s'exécutent sur un poste de travail. Deux heuristiques déterministes donnent une réponse immédiate ; deux solveurs les complètent dans un budget de temps fixé, sans garantie d'atteindre l'optimum dans ce délai. Les prévisions reposent sur des modèles simples et explicables, et l'ensemble est servi par une API typée et un tableau de bord Next.js, sans base de données. Ces choix étant fixés, le chapitre suivant formalise le problème que les méthodes doivent résoudre et décrit la conception du système.

\chapter{Analyse et Conception du Système}
\label{ch:conception_systeme}

% ------------------------------------------------------------------------------
\section*{Introduction}
\addcontentsline{toc}{section}{Introduction}
% ------------------------------------------------------------------------------
Planifier ensemble les lits et le personnel demande un modèle précis : ce que l'on décide, ce que l'on interdit et ce que l'on accepte de dégrader. Ce chapitre décrit l'architecture du système, la dynamique des patients et les contraintes de personnel, puis la fonction objectif et ses trois profils de pondération. Il présente ensuite la conception des quatre méthodes, du validateur indépendant, des scénarios et de l'interface. La formulation est celle du code ; sa version de référence est le document \texttt{docs/mathematical\_model.md}.

% ------------------------------------------------------------------------------
\section{Architecture Générale du Système}
% ------------------------------------------------------------------------------

\begin{figure}[H]
    \centering
    \captureimg{architecture_pipeline}
    \caption{Architecture fonctionnelle globale du système de planification hospitalière.}
    \label{fig:global_architecture}
\end{figure}

Le traitement s'articule en cinq étapes :
\begin{enumerate}
    \item \textbf{Données} : génération et contrôle de \NbJours{} jours de données synthétiques ;
    \item \textbf{Prévision} : admissions attendues par service de $J+1$ à $J+7$ ;
    \item \textbf{Besoins opérationnels} : conversion des prévisions en admissions demandées, occupation, capacité et cibles de personnel ; c'est l'\emph{instance} commune à toutes les méthodes ;
    \item \textbf{Optimisation et contrôle} : calcul des plans par quatre méthodes, puis contrôle par le validateur ;
    \item \textbf{Restitution} : enregistrement des expériences, API et tableau de bord.
\end{enumerate}

% ------------------------------------------------------------------------------
\section{Modélisation Métier}
% ------------------------------------------------------------------------------

\subsection{Ensembles et Paramètres}
\begin{table}[H]
    \centering
    \caption{Principales notations du modèle.}
    \label{tab:notations}
    \resizebox{\textwidth}{!}{
    \begin{tabular}{lp{10.6cm}l}
        \toprule
        \textbf{Symbole} & \textbf{Signification} & \textbf{Nature} \\
        \midrule
        $\mathcal{S}, \mathcal{T}, \mathcal{K}, \mathcal{R}, \mathcal{I}$ & Services, jours (7), postes (3), rôles (3) et professionnels actifs & Ensembles \\
        $q_{s,t}$ & Admissions demandées, arrondi supérieur de la prévision & Paramètre \\
        $o_{s,t}$ & Occupation si toute la demande est acceptée & Paramètre \\
        $c_{s,t}$, $\ell_s$ & Lits disponibles, durée de séjour fixe en jours & Paramètres \\
        $\rho_r$, $m_r$ & Cible pour dix lits ouverts et par poste, minimum par poste & Paramètres \\
        $n_i$ & Nombre maximal de postes du professionnel $i$ sur l'horizon & Paramètre \\
        $A_{s,t}$, $O_{s,t}$ & Admissions acceptées, occupation soutenue & Entiers \\
        $L_{s,t}$, $Y_{s,t}$ & Lits ouverts, service ouvert & Entier, binaire \\
        $X_{i,s,t,k}$ & Affectation, créée seulement si $i$ est disponible et qualifié & Binaire \\
        $Q_{s,t,k,r}$ & Postes manquants par rapport à la cible & Entier \\
        $W^{\min}_r$, $W^{\max}_r$ & Charges minimale et maximale du rôle $r$, en postes & Entiers \\
        \bottomrule
    \end{tabular}
    }
\end{table}

En $S_0$, le modèle compte \NbVariables{} variables, dont \NbAffectations{} variables d'affectation et \NbDeficits{} variables de déficit, et \NbContraintes{} contraintes.

\subsection{Dynamique des Séjours}
Les sorties sont traitées avant les admissions. La durée de séjour est fixe par service : un patient refusé le jour $j$ aurait occupé un lit de $j$ à $j + \ell_s - 1$. L'occupation soutenue est donc :
\begin{equation}
    O_{s,t} + \sum_{j=\max(1,\,t-\ell_s+1)}^{t} \left(q_{s,j} - A_{s,j}\right) = o_{s,t}
    \quad \forall s \in \mathcal{S}, \; \forall t \in \mathcal{T}, \qquad 0 \le A_{s,t} \le q_{s,t}
\end{equation}

La demande non soutenue, en journées-patients, n'est pas une variable supplémentaire : $D_{s,t} = o_{s,t} - O_{s,t}$.

\subsection{Capacité en Lits}
Les patients présents doivent tenir dans les lits ouverts, qui ne dépassent pas les lits disponibles :
\begin{equation}
    O_{s,t} \le L_{s,t} \le c_{s,t} \, Y_{s,t}, \qquad L_{s,t} \ge Y_{s,t}
    \quad \forall s \in \mathcal{S}, \; \forall t \in \mathcal{T}
\end{equation}

\subsection{Dotation Minimale et Cible de Personnel}
Le modèle distingue deux niveaux, dont les valeurs sont données au tableau~\ref{tab:effectifs}.

La \textbf{dotation minimale} est une contrainte dure, indépendante du nombre de lits :
\begin{equation}
    \sum_{i \in \mathcal{I}_r} X_{i,s,t,k} \ge m_r \, Y_{s,t} \quad \forall s, t, k, r
\end{equation}

La \textbf{cible} dépend des lits ouverts ; elle peut ne pas être atteinte, au prix d'un déficit $Q$ :
\begin{equation}
    10 \left( \sum_{i \in \mathcal{I}_r} X_{i,s,t,k} + Q_{s,t,k,r} \right) \ge \rho_r \, L_{s,t} \quad \forall s, t, k, r
\end{equation}

où $\mathcal{I}_r$ est l'ensemble des professionnels du rôle $r$. Les cibles sont identiques pour tous les services et tous les postes : c'est une simplification.

\subsection{Contraintes Individuelles}
\begin{enumerate}
    \item \textbf{Un poste par jour} :
    \begin{equation}
        \sum_{s \in \mathcal{S}} \sum_{k \in \mathcal{K}} X_{i,s,t,k} \le 1 \quad \forall i \in \mathcal{I}, \; \forall t \in \mathcal{T}
    \end{equation}
    \item \textbf{Nombre maximal de postes} (\PostesMax{} sur l'horizon) :
    \begin{equation}
        \sum_{s \in \mathcal{S}} \sum_{t \in \mathcal{T}} \sum_{k \in \mathcal{K}} X_{i,s,t,k} \le n_i \quad \forall i \in \mathcal{I}
    \end{equation}
    \item \textbf{Repos minimal de \ReposMin{} heures} entre deux jours consécutifs. Soit $\mathcal{F}$ l'ensemble des couples de postes séparés par un repos insuffisant :
    \begin{equation}
        \sum_{s \in \mathcal{S}} X_{i,s,t,k} + \sum_{s \in \mathcal{S}} X_{i,s,t+1,k'} \le 1 \quad \forall i, \; \forall t, \; \forall (k,k') \in \mathcal{F}
    \end{equation}
    Avec les horaires retenus, $\mathcal{F}$ contient trois couples : soir puis matin, nuit puis matin, nuit puis soir. Deux nuits consécutives sont permises.
    \item \textbf{Nuits consécutives} : au plus \NuitsMax{} nuits sur toute fenêtre de quatre jours de l'horizon.
\end{enumerate}

\subsection{Politique de Sous-Effectif}
La politique \texttt{shortage\_permitting}, utilisée pour tous les résultats, n'assouplit que la cible de personnel, par la variable $Q$. La politique \texttt{strict} impose $A = q$ et $Q = 0$. Les autres contraintes restent dures dans les deux cas.

\begin{notebox}{Conséquence à retenir}
Ouvrir un lit n'exige que la dotation minimale du service. Un plan peut donc \textbf{respecter les contraintes} tout en étant très loin de sa cible de personnel. C'est pourquoi l'outil affiche toujours, à côté, le pourcentage de \textbf{postes cibles couverts}.
\end{notebox}

\subsection{Équité}
L'équité porte sur les professionnels actifs \textbf{présents}, c'est-à-dire disponibles à au moins un poste de l'horizon. Pour chaque rôle, $W^{\min}_r \le \text{charge}_i \le W^{\max}_r$ ; l'indicateur est la somme des étendues $W^{\max}_r - W^{\min}_r$. Un agent absent sur tout l'horizon est exclu, sinon l'étendue resterait élevée dans les scénarios d'absence quelle que soit la méthode.

% ------------------------------------------------------------------------------
\section{Conception de la Chaîne de Données et de la Prévision}
% ------------------------------------------------------------------------------

\subsection{Générateur Synthétique}
Le générateur couvre \NbJours{} jours, du \DateDebut{} au \DateFin{}. Les admissions d'un service suivent une loi de Poisson dont le paramètre additionne :
\begin{itemize}
    \item un niveau de base par service ;
    \item un effet du jour de la semaine ;
    \item une tendance journalière ;
    \item un effet d'événement exceptionnel, tiré avec une probabilité de 3\,\% par jour.
\end{itemize}

Le générateur ne comporte pas de saisonnalité annuelle. Il produit aussi les fermetures temporaires de lits, les absences et la disponibilité de chaque professionnel à chaque poste (\LignesDispo{} lignes). La graine fixe le résultat.

\subsection{Contrôle de Qualité}
Le contrôle vérifie les doublons, les valeurs invalides et la conservation des patients. Il ne conclut pas à l'absence de valeurs vides : le tableau~\ref{tab:qualite} montre que deux colonnes en contiennent, et pourquoi.

% Généré par rapport/scripts/generate_tables.py : ne pas modifier à la main.
\begin{table}[H]
    \centering
    \caption{Rapport de qualité des données synthétiques (\texttt{data/synthetic/quality\_report.json}).}
    \label{tab:qualite}
    \small
    \begin{tabular}{lrrrl}
        \toprule
        \textbf{Table} & \textbf{Lignes} & \textbf{Valeurs vides} & \textbf{Doublons} & \textbf{Explication} \\
        \midrule
        \texttt{historical\_demand} & 4\,380 & 3\,285 & 0 & \texttt{emergency\_arrivals} vide hors URG \\
        \texttt{bed\_capacity} & 4\,380 & 0 & 0 & -- \\
        \texttt{occupancy} & 4\,380 & 0 & 0 & -- \\
        \texttt{staff} & 235 & 0 & 0 & -- \\
        \texttt{staff\_availability} & 771\,975 & 748\,883 & 0 & \texttt{absence\_type} vide si disponible \\
        \texttt{staff\_requirements} & 39\,420 & 0 & 0 & -- \\
        \bottomrule
    \end{tabular}
    \par\vspace{3pt}
    {\footnotesize Période couverte : 2023-01-01 au 2025-12-30. Valeurs invalides : 0. Verdict : conforme.}
\end{table}


\subsection{Protocole de Prévision et Erreur WAPE}
Les modèles sont évalués chronologiquement : apprentissage sur 730 jours au minimum, 12 origines de validation de 7 jours (soit \NbPrevVal{} prévisions par service et par modèle), puis un test final de \NbPrevTest{} jours réservé. Aucune donnée postérieure à l'origine n'est utilisée. Le modèle retenu pour un service est celui qui minimise l'erreur de validation :
\begin{equation}
    \text{WAPE} = \frac{\sum_{t=1}^{N} |y_t - \hat{y}_t|}{\sum_{t=1}^{N} y_t} \times 100\,\%
\end{equation}
Le WAPE reste défini lorsque certaines valeurs journalières sont nulles, ce qui n'est pas le cas d'une moyenne d'erreurs relatives.

% ------------------------------------------------------------------------------
\section{Fonction Objectif et Profils de Pondération}
% ------------------------------------------------------------------------------

\subsection{Fonction Objectif}
Toutes les méthodes sont évaluées avec la même fonction $\mathcal{Z}$ :
\begin{equation}
    \min \; \mathcal{Z} = c_{\text{lit}} \sum_{s,t} L_{s,t}
    + \sum_{s,t} c^{s}_{\text{dem}} \, D_{s,t}
    + \sum_{s,t} c^{s}_{\text{ref}} \, (q_{s,t} - A_{s,t})
    + \sum_{s,t,k,r} c^{s}_{\text{rh}} \, Q_{s,t,k,r}
    + c_{\text{eq}} \sum_{r} \left(W^{\max}_r - W^{\min}_r\right)
    \label{eq:objectif}
\end{equation}

La configuration de base (\texttt{hospital.json}) fixe \textbf{quatre coefficients} : $c_{\text{lit}} = 10$ par lit-jour ouvert, $c_{\text{dem}} = 1\,000$ par journée-patient non soutenue, $c_{\text{rh}} = 500$ par poste cible non couvert et $c_{\text{eq}} = 1$ par poste d'étendue de charge. Les profils peuvent y ajouter un coût par admission refusée $c_{\text{ref}}$, nul par défaut, et des multiplicateurs par service. Ces coefficients sont des points sans unité monétaire : l'objectif n'est pas un coût financier.

\subsection{Trois Profils}
% Généré par rapport/scripts/generate_tables.py : ne pas modifier à la main.
\begin{table}[H]
    \centering
    \caption{Poids de l'objectif selon le profil (\texttt{configs/profiles.json}).}
    \label{tab:profils}
    \resizebox{\textwidth}{!}{\begin{tabular}{llccc}
        \toprule
        \textbf{Coefficient} & \textbf{Champ} & \textbf{Équilibré} & \textbf{Priorité accueil} & \textbf{Priorité personnel} \\
        \midrule
        $c_{\text{lit}}$, lit-jour ouvert & \texttt{opened\_bed} & 10 & 10 & 10 \\
        $c_{\text{dem}}$, journée-patient non soutenue & \texttt{unmet\_patient\_day} & 1\,000 & 3\,000 (URG $\times$3) & 1\,000 \\
        $c_{\text{ref}}$, admission refusée & \texttt{rejected\_admission} & 0 & 5\,000 (URG $\times$3) & 0 \\
        $c_{\text{rh}}$, poste cible non couvert & \texttt{staff\_target\_shortage} & 500 & 500 (URG $\times$2) & 750 \\
        $c_{\text{eq}}$, poste d'étendue de charge & \texttt{workload\_imbalance} & 1 & 1 & 100 \\
        \bottomrule
    \end{tabular}}
\end{table}


\begin{itemize}
    \item \textbf{Équilibré} : poids de la configuration de base ;
    \item \textbf{Priorité accueil} : journées-patients et refus fortement pénalisés, surtout aux Urgences ;
    \item \textbf{Priorité personnel} : déficit de personnel et déséquilibre de charge plus pénalisés.
\end{itemize}

Ces valeurs ont été choisies par l'équipe. Les objectifs de deux profils ne sont pas comparables : on compare leurs indicateurs (refus, couverture, déficit, équité).

\subsection{Effet de Palier}
Dix lits ouverts demandent 12 postes cibles par jour : avec le profil équilibré et sans personnel disponible, un lit coûte en moyenne $10 + 1{,}2 \times 500 = 610$ points, moins que les $1\,000$ points d'une journée-patient non soutenue. Accueillir reste donc préférable en moyenne. Mais la cible est entière : passer de 5 à 6 lits ajoute 3 postes infirmiers par jour (1\,510 points avec le lit), et passer de 10 à 11 lits ajoute 9 postes (4\,510 points). Refuser un patient pour rester sous un palier devient alors avantageux pour l'objectif.

Ce mécanisme se lit dans les résultats : en $S_0$, profil équilibré, CP-SAT ouvre un nombre de lits différent de Greedy dans \PaliersDiffEq{} cas (un service, un jour) ; dans \PaliersCinqEq{} d'entre eux, il s'arrête sur un multiple de 5. Avec la priorité personnel, c'est le cas de \PaliersCinqPer{} sur \PaliersDiffPer{}.

% ------------------------------------------------------------------------------
\section{Conception des Quatre Méthodes}
% ------------------------------------------------------------------------------

\begin{table}[H]
    \centering
    \caption{Principe des quatre méthodes.}
    \label{tab:methodes}
    \resizebox{\textwidth}{!}{
    \begin{tabular}{lp{6.2cm}p{6.6cm}cc}
        \toprule
        \textbf{Méthode} & \textbf{Admissions et lits} & \textbf{Personnel} & \textbf{Lit les poids} & \textbf{Déterministe} \\
        \midrule
        Naïve & Taux d'acceptation unique par jour, dans la limite des lits de chaque service. & Dotation minimale, puis même fraction de la cible partout (plus forts restes). & Non & Oui \\
        Greedy & Autant d'admissions que de lits libres, service par service. & Dotation minimale, puis cible ; ordre date, priorité du service, poste, rôle ; premier agent admissible. & Non & Oui \\
        MILP/CBC & Formulation complète. & Formulation complète, sans solution initiale. & Oui & Non garanti \\
        CP-SAT & Formulation complète. & Formulation complète, avec le plan Greedy comme indication de départ. & Oui & Non \\
        \bottomrule
    \end{tabular}
    }
\end{table}

Naïve et Greedy ne lisent pas les poids : leur plan est identique pour les trois profils, seule son évaluation change. Aucune solution initiale n'est donnée à CBC : lors d'essais, il restait bloqué sur le plan Greedy, ce qui rendait la comparaison sans objet.

% ------------------------------------------------------------------------------
\section{Conception du Validateur Indépendant}
% ------------------------------------------------------------------------------

Le module \texttt{validator.py} recalcule les règles depuis les tables produites, sans réutiliser les contraintes des solveurs ni savoir quelle méthode a produit le plan. Il contrôle :
\begin{enumerate}
    \item la capacité et le bilan des lits de chaque jour ;
    \item les admissions : $0 \le \text{acceptées} \le \text{demandées}$, et refus cohérents ;
    \item la \textbf{dynamique des cohortes}, rejouée depuis l'instance : recensement initial, sorties et capacité sont comparés ligne à ligne ;
    \item la disponibilité et la qualification de chaque agent affecté ;
    \item l'absence de double affectation journalière ;
    \item le nombre maximal de postes ;
    \item le repos minimal et les nuits consécutives ;
    \item la dotation minimale de chaque service ouvert.
\end{enumerate}

Le troisième contrôle renforce une première version qui ne vérifiait que le bilan interne du plan : ce bilan est toujours juste pour une méthode qui en déduit elle-même ses sorties. Le déficit de cible est compté à part et n'invalide un plan que sous la politique \texttt{strict}. L'interface n'affiche donc jamais « solution valide », mais « contraintes respectées » accompagné du pourcentage de postes cibles couverts :
\begin{equation}
    \text{Postes couverts} = 1 - \frac{\sum_{s,t,k,r} \max(\text{cible} - \text{affectés},\, 0)}{\sum_{s,t,k,r} \text{cible}}
\end{equation}

Un surplus d'effectif sur un besoin ne compense pas un manque sur un autre.

% ------------------------------------------------------------------------------
\section{Conception des Scénarios}
% ------------------------------------------------------------------------------

Les six scénarios modifient l'instance de référence sans toucher aux paramètres du modèle.

% Généré par rapport/scripts/generate_tables.py : ne pas modifier à la main.
\begin{table}[H]
    \centering
    \caption{Définition des six scénarios (\texttt{configs/scenarios.json}).}
    \label{tab:scenarios}
    \resizebox{\textwidth}{!}{\begin{tabular}{llp{5.2cm}cc}
        \toprule
         & \textbf{Identifiant} & \textbf{Admissions demandées} & \textbf{Absences ajoutées} & \textbf{Lits fermés} \\
        \midrule
        $S_0$ & \texttt{S0\_BASELINE} & $\times$1{,}00, tous services & 0\,\% & 0\,\% \\
        $S_1$ & \texttt{S1\_URGENT\_PEAK} & URG $\times$1{,}50, 2 premiers jours & 0\,\% & 0\,\% \\
        $S_2$ & \texttt{S2\_GENERAL\_SURGE} & $\times$1{,}25, tous services & 0\,\% & 0\,\% \\
        $S_3$ & \texttt{S3\_STAFF\_ABSENCE} & $\times$1{,}00, tous services & 15\,\% & 0\,\% \\
        $S_4$ & \texttt{S4\_COMBINED\_CRISIS} & $\times$1{,}25, tous services & 15\,\% & 0\,\% \\
        $S_5$ & \texttt{S5\_BED\_CLOSURE} & $\times$1{,}00, tous services & 0\,\% & 10\,\% \\
        \bottomrule
    \end{tabular}}
    \par\vspace{3pt}
    {\footnotesize Absences : part des professionnels disponibles rendus indisponibles sur tout l'horizon (tirage fixé par la graine). Lits fermés : réduction de la capacité de chaque service, arrondie à l'entier inférieur.}
\end{table}


Les \textbf{scénarios personnalisés} du simulateur utilisent le même mécanisme : multiplicateur d'admissions de 1 à 2, absences et fermetures de 0 à 50\,\%, globalement ou pour un service. Ils reçoivent un identifiant \texttt{CUSTOM\_} et ne sont jamais comptés avec $S_0$ à $S_5$.

% ------------------------------------------------------------------------------
\section{Conception de l'Interface Utilisateur Web}
% ------------------------------------------------------------------------------

Le tableau de bord comporte \textbf{dix vues}, accessibles par une barre latérale :
\begin{enumerate}
    \item \textbf{Vue d'ensemble} : indicateurs d'une expérience et d'une méthode ;
    \item \textbf{Comparaison scénarios} : méthodes, scénarios et profils côte à côte ;
    \item \textbf{Exploration des données} : données synthétiques et instance ;
    \item \textbf{Prévisions} : prévisions et erreurs des modèles ;
    \item \textbf{Allocation} : lits, postes manquants, couverture RH et planning du personnel ;
    \item \textbf{Méthodes} : comparaison des méthodes d'une expérience ;
    \item \textbf{Scénarios} : définition des scénarios ;
    \item \textbf{Simulateur} : scénario personnalisé calculé à la demande ;
    \item \textbf{Sensibilité} : lecture de la grille de \NbCas{} cas ;
    \item \textbf{Expériences et export} : expériences enregistrées et téléchargements.
\end{enumerate}

Deux règles guident l'affichage : aucune vue ne lance de calcul à son ouverture, et une valeur absente donne un état vide explicite.

% ------------------------------------------------------------------------------
\section*{Conclusion}
\addcontentsline{toc}{section}{Conclusion}
% ------------------------------------------------------------------------------
Ce chapitre a fixé le modèle : cohortes de patients, capacité en lits, dotation minimale dure, cible de personnel assouplie par un déficit, règles individuelles et objectif pondéré selon trois profils. Il a aussi mis en évidence l'effet de palier, qui peut rendre un refus avantageux pour l'objectif. Le validateur est séparé des méthodes, et deux notions ne sont jamais confondues : le respect des contraintes et la couverture des postes cibles. Le chapitre suivant décrit la mise en œuvre de cette conception, des données jusqu'au tableau de bord.

\chapter{Réalisation et Implémentation}
\label{ch:realisation_implementation}

La réalisation traduit le modèle du chapitre précédent en un logiciel modulaire, testé et utilisable depuis un navigateur.

Ce chapitre décrit l'organisation du code, l'implémentation des données, des prévisions, des quatre méthodes et du validateur, puis l'API FastAPI et les principales vues du tableau de bord, dont l'allocation, le simulateur et la sensibilité.

% ------------------------------------------------------------------------------
\section{Environnement de Développement et Organisation du Projet}
% ------------------------------------------------------------------------------

Le projet réunit dans un même dépôt le moteur Python, l'API et le tableau de bord :

\begin{lstlisting}[caption={Structure arborescente du projet (fichiers principaux).}]
hospital-resource-planning/
|-- configs/                  # hospital.json, scenarios.json, profiles.json
|-- data/synthetic/           # Donnees synthetiques et rapport de qualite
|-- docs/                     # Modele, solveurs, experiences, dashboard
|-- src/hospital_planning/
|   |-- domain/               # contracts.py, enums.py (schemas Pydantic)
|   |-- config/               # loader.py
|   |-- data/                 # generator.py, occupancy.py, validation.py
|   |-- forecasting/          # baselines, statistical, machine_learning, evaluation
|   |-- operational/          # requirements.py (instance commune)
|   |-- optimization/         # naive, greedy, exact, objective, validator, availability
|   |-- experiments/          # scenarios.py, runner.py, sensitivity.py
|   |-- solvers/              # smoke.py (disponibilite des solveurs)
|   \-- api/                  # app.py, jobs.py
|-- frontend/
|   |-- app/                  # layout.tsx, page.tsx, globals.css
|   |-- components/           # dashboard, staffing, simulator, sensitivity
|   |-- lib/                  # api-client.ts, api-types.ts
|   \-- tests/                # Parcours Playwright
|-- results/                  # Previsions, instance S0, experiences
|-- scripts/                  # verify_api.ps1
\-- tests/                    # Tests pytest
\end{lstlisting}

Les chemins écrits dans les artefacts sont relatifs et utilisent le séparateur \texttt{/}, afin qu'une expérience calculée sous Windows soit lisible sous Linux ou macOS.

% ------------------------------------------------------------------------------
\section{Implémentation des Données et des Prévisions}
% ------------------------------------------------------------------------------

\subsection{Générateur de Données (\texttt{data/generator.py})}
Le générateur lit \texttt{configs/hospital.json} et écrit six tables CSV : demande, capacité, occupation, personnel, disponibilité et besoins en personnel. La graine (42) rend le résultat reproductible. Le module \texttt{occupancy.py} en déduit sorties et occupation à partir des durées de séjour, et \texttt{validation.py} produit le rapport de qualité sans corriger les données.

\subsection{Prévision (\texttt{forecasting/})}
Les cinq modèles sont répartis dans \texttt{baselines.py} (naïf, naïf saisonnier, moyenne mobile), \texttt{statistical.py} (lissage exponentiel) et \texttt{machine\_learning.py} (régression à retards). Le listing suivant montre les variables de la régression, toutes calculées à partir du passé ou du calendrier :

\begin{lstlisting}[language=Python, caption={Variables explicatives de la régression à retards (extrait de \texttt{machine\_learning.py}).}]
weekday = date.weekday()
return [
    values[-1],
    values[-7],
    values[-14],
    float(np.mean(values[-7:])),
    float(np.mean(values[-28:])),
    *(1.0 if weekday == day else 0.0 for day in range(1, 7)),
    float(date.strftime("%m-%d") in holiday_month_days),
]
\end{lstlisting}

Le module \texttt{evaluation.py} calcule MAE, RMSE, WAPE et biais, choisit le modèle de chaque service et enregistre les prévisions de $J+1$ à $J+7$. Les prévisions restent continues ; leur conversion en admissions entières (arrondi supérieur) a lieu une seule fois, dans \texttt{operational/requirements.py}.

% ------------------------------------------------------------------------------
\section{Implémentation des Quatre Méthodes}
% ------------------------------------------------------------------------------

\subsection{Stratégie Naïve (\texttt{naive.py})}
Chaque jour, les lits libres de l'hôpital sont répartis au prorata des demandes par la méthode des plus forts restes, dans la limite des lits de chaque service. Le personnel est partagé de la même façon après la dotation minimale. À reste égal, l'ordre des clés départage, ce qui rend le résultat déterministe.

\subsection{Heuristique Gloutonne (\texttt{greedy.py})}
Pour chaque service et chaque jour, l'heuristique applique les sorties, accepte autant d'admissions que de lits libres et ouvre exactement les lits occupés. Elle remplit ensuite tous les besoins minimaux, puis les cibles. Un agent n'est retenu que s'il respecte disponibilité, qualification, nombre de postes, repos et nuits consécutives.

\subsection{MILP avec PuLP et CBC (\texttt{exact.py})}
Les variables d'affectation ne sont créées que pour les couples disponibles et qualifiés, ce qui réduit la taille du modèle. L'extrait suivant montre le lien entre lits ouverts, dotation minimale et cible :

\begin{lstlisting}[language=Python, caption={Dotation minimale et cible de personnel (extrait de \texttt{exact.py}).}]
minimum = config.synthetic_data.minimum_staff_by_role[role]
ratio = config.synthetic_data.staff_per_ten_beds[role]
shortage = variables["shortage"][(service, date_index, shift, role)]
model += staffed >= minimum * open_indicator
model += 10 * (staffed + shortage) >= ratio * opened
\end{lstlisting}

Trois précautions ont été nécessaires avec CBC :
\begin{itemize}
    \item \textbf{Statut} : CBC peut annoncer un optimum alors que sa recherche a été interrompue ; le statut \texttt{OPTIMAL} n'est retenu que si le journal confirme une recherche terminée ;
    \item \textbf{Borne} : la meilleure borne est lue dans le journal, puis corrigée du terme constant de l'objectif, pour calculer le gap ;
    \item \textbf{Arrêt anormal} : certaines graines déclenchent une erreur interne de CBC ; le calcul est relancé avec la graine suivante, dans le temps restant. Cela s'est produit pour \NbReprisesCbc{} résultat sur 18.
\end{itemize}

\subsection{CP-SAT avec OR-Tools (\texttt{exact.py})}
CP-SAT résout la même formulation. Le plan Greedy lui est donné comme indication de départ (\emph{hint}), et la recherche utilise \ThreadsSolveur{} threads. La graine est fixée, mais la recherche parallèle limitée en temps n'est pas reproductible à l'identique.

\subsection{Budget des Solveurs}
Le budget (\TempsSolveur{} secondes et \ThreadsSolveur{} threads par solveur) est lu dans \texttt{hospital.json}, et peut être modifié en ligne de commande (\texttt{-{}-time-limit}, \texttt{-{}-workers}) ou dans une requête à l'API. Les deux solveurs tiennent ensemble dans deux minutes par scénario et par profil.

% ------------------------------------------------------------------------------
\section{Implémentation du Validateur de Contraintes (\texttt{validator.py})}
% ------------------------------------------------------------------------------

Le validateur lit les tables produites par une méthode. Le contrôle des cohortes rejoue la dynamique depuis l'instance :

\begin{lstlisting}[language=Python, caption={Contrôle des cohortes par le validateur (extrait de \texttt{validator.py}).}]
discharges = max(int(expected["expected_discharges"]) - rejected.get(row.date - stay, 0), 0)
consistent = (
    int(row.occupancy_start) == census
    and int(row.expected_discharges) == discharges
    and int(row.capacity_available) == int(expected["capacity_available"])
)
violations += int(not consistent)
\end{lstlisting}

Chaque famille de violations a son compteur dans \texttt{method\_result.json}. La table \texttt{staff\_coverage.csv}, construite au même moment, donne pour chaque service, jour, poste et rôle le minimum, la cible, l'effectif affecté et le déficit ; elle alimente la vue Allocation.

% ------------------------------------------------------------------------------
\section{Expériences et Profils (\texttt{experiments/})}
% ------------------------------------------------------------------------------

\begin{itemize}
    \item \textbf{Scénarios} : \texttt{scenarios.py} transforme l'instance de référence sans la modifier et calcule une empreinte SHA-256 des entrées, indépendante du système d'exploitation ;
    \item \textbf{Expériences} : \texttt{runner.py} enregistre une expérience par scénario et par profil. Relancer une méthode remplace son seul résultat ;
    \item \textbf{Sensibilité} : \texttt{sensitivity.py} calcule la grille de \NbCas{} cas et fusionne les résultats dans \texttt{sensitivity\_results.csv} ;
    \item \textbf{Traçabilité} : chaque résultat porte la graine, le profil, les poids, le budget, le statut du solveur et la date.
\end{itemize}

% ------------------------------------------------------------------------------
\section{Implémentation de l'API REST FastAPI (\texttt{api/app.py})}
% ------------------------------------------------------------------------------

\subsection{Points d'Accès}
L'API expose \NbEndpoints{} points d'accès (tableau~\ref{tab:endpoints}). Les lectures ne déclenchent jamais de calcul : seule la requête \texttt{POST} sur \texttt{/api/v1/jobs/experiments} en lance un.

% Généré par rapport/scripts/generate_tables.py : ne pas modifier à la main.
\begin{table}[H]
    \centering
    \caption{Les 16 points d'accès de l'API, extraits de sa description OpenAPI.}
    \label{tab:endpoints}
    \resizebox{\textwidth}{!}{\begin{tabular}{lp{8.2cm}p{8.4cm}}
        \toprule
        \textbf{Verbe} & \textbf{Route} & \textbf{Rôle} \\
        \midrule
        \texttt{GET} & \texttt{/\allowbreak health} & Expose l'état minimal du service sans déclencher de calcul. \\
        \texttt{GET} & \texttt{/\allowbreak api/\allowbreak v1/\allowbreak config/\allowbreak hospital} & Retourne la configuration validée de l'hôpital synthétique. \\
        \texttt{POST} & \texttt{/\allowbreak api/\allowbreak v1/\allowbreak jobs/\allowbreak experiments} & Soumet un job explicite. \\
        \texttt{GET} & \texttt{/\allowbreak api/\allowbreak v1/\allowbreak jobs/\allowbreak \{job\_id\}} & Retourne l'état réel d'un job local sans jamais le relancer. \\
        \texttt{GET} & \texttt{/\allowbreak api/\allowbreak v1/\allowbreak experiments} & Liste les expériences déjà stockées. \\
        \texttt{GET} & \texttt{/\allowbreak api/\allowbreak v1/\allowbreak experiments/\allowbreak \{experiment\_id\}} & Retourne les métadonnées et KPI d'une expérience stockée, sans calcul. \\
        \texttt{GET} & \texttt{/\allowbreak api/\allowbreak v1/\allowbreak dashboard/\allowbreak scenario-comparison} & Agrège les KPI persistés d'une seule expérience par scénario, pour un profil. \\
        \texttt{GET} & \texttt{/\allowbreak api/\allowbreak v1/\allowbreak experiments/\allowbreak \{experiment\_id\}/\allowbreak downloads/\allowbreak \{artifact\_path\}} & Télécharge un artefact existant après contrôle strict du chemin. \\
        \texttt{GET} & \texttt{/\allowbreak api/\allowbreak v1/\allowbreak dashboard/\allowbreak scenarios} & Expose les scénarios déclarés, sans les transformer ni les exécuter. \\
        \texttt{GET} & \texttt{/\allowbreak api/\allowbreak v1/\allowbreak dashboard/\allowbreak profiles} & Expose les profils de pondération. \\
        \texttt{GET} & \texttt{/\allowbreak api/\allowbreak v1/\allowbreak dashboard/\allowbreak sensitivity/\allowbreak manifest} & Expose le manifeste déjà produit, sans déclencher la grille coûteuse. \\
        \texttt{GET} & \texttt{/\allowbreak api/\allowbreak v1/\allowbreak dashboard/\allowbreak sensitivity/\allowbreak results} & Lit les KPI de sensibilité déjà calculés pour un profil, sans lancer de calcul. \\
        \texttt{GET} & \texttt{/\allowbreak api/\allowbreak v1/\allowbreak dashboard/\allowbreak forecasts} & Lit les prévisions finales et métriques enregistrées, sans les recalculer. \\
        \texttt{GET} & \texttt{/\allowbreak api/\allowbreak v1/\allowbreak dashboard/\allowbreak exploration} & Lit les séries synthétiques historiques nécessaires aux graphiques d'exploration. \\
        \texttt{GET} & \texttt{/\allowbreak api/\allowbreak v1/\allowbreak dashboard/\allowbreak operational/\allowbreak \{scenario\_id\}} & Lit une instance canonique opérationnelle, sans déclencher de décision. \\
        \texttt{GET} & \texttt{/\allowbreak api/\allowbreak v1/\allowbreak dashboard/\allowbreak experiments/\allowbreak \{experiment\_id\}/\allowbreak methods/\allowbreak \{method\_name\}} & Lit les allocations et la couverture RH d'une méthode déjà terminée. \\
        \bottomrule
    \end{tabular}}
\end{table}


\begin{figure}[H]
    \centering
    \captureimg{swagger_ui}
    \caption{Documentation interactive de l'API (Swagger UI, \texttt{/docs}).}
    \label{fig:swagger}
\end{figure}

\subsection{Gestion des Calculs Longs (\texttt{api/jobs.py})}
Un calcul soumis reçoit immédiatement une réponse \texttt{202} avec un identifiant de job. Il passe par les états \texttt{QUEUED}, \texttt{RUNNING}, puis \texttt{COMPLETED} ou \texttt{FAILED}. Un seul travailleur exécute les jobs, l'un après l'autre :

\begin{lstlisting}[language=Python, caption={File d'exécution à un seul travailleur (extrait de \texttt{jobs.py}).}]
self._executor = ThreadPoolExecutor(max_workers=1, thread_name_prefix="hospital-job")

def submit(self, task: Callable[[], list[str]]) -> JobResponse:
    job_id = f"job-{uuid4().hex}"
    record = JobRecord(job_id=job_id, created_at=datetime.now(UTC))
    with self._lock:
        self._records[job_id] = record
    self._executor.submit(self._run, record, task)
    return record.response()
\end{lstlisting}

L'état des jobs est gardé en mémoire : il est perdu au redémarrage de l'API, mais les expériences terminées restent sur disque.

\subsection{Contrôles d'Accès aux Fichiers et aux Origines}
Le chemin d'un téléchargement est résolu, puis refusé s'il sort du dossier de l'expérience :

\begin{lstlisting}[language=Python, caption={Contrôle du chemin avant téléchargement (extrait de \texttt{app.py}).}]
directory = _experiment_directory(experiment_id)
artifact = (directory / artifact_path).resolve()
if not artifact.is_relative_to(directory) or not artifact.is_file():
    raise HTTPException(status_code=404, detail="Artefact introuvable.")
return FileResponse(artifact, filename=artifact.name)
\end{lstlisting}

Les origines CORS sont explicites et configurables (\texttt{HOSPITAL\_CORS\_ORIGINS}). Un profil ou un service inconnu donne une erreur \texttt{422}. L'API n'a pas d'authentification : elle est prévue pour un usage local.

% ------------------------------------------------------------------------------
\section{Réalisation du Tableau de Bord (Next.js 15)}
% ------------------------------------------------------------------------------

Le tableau de bord est une page unique ; la barre latérale ouvre l'une des dix vues, et une barre de filtres choisit le service, le scénario, l'expérience et la méthode. Toutes les valeurs viennent de l'API.

\subsection{Vue d'Ensemble}
\begin{figure}[H]
    \centering
    \captureimg{dashboard_overview}
    \caption{Vue d'ensemble : indicateurs d'une expérience et d'une méthode.}
    \label{fig:overview}
\end{figure}

La vue affiche séparément les admissions, les refus, les lits ouverts et les postes cibles couverts. La mention « contraintes respectées » est toujours accompagnée de ce dernier indicateur.

\subsection{Vue Comparaison Scénarios}
La vue compare les quatre méthodes sur les six scénarios, pour un profil choisi dans un sélecteur et pour un indicateur parmi neuf : objectif, refus, demande non soutenue, couverture de la demande, déficit RH, postes couverts, lits ouverts, équité et temps de calcul. Un tableau complet suit le graphique, avec la meilleure méthode de chaque scénario (figure~\ref{fig:comparaison}). Les simulations personnalisées y figurent dans un tableau distinct.

\subsection{Vue Allocation : Postes Manquants, Couverture RH et Planning}
\begin{figure}[H]
    \centering
    \captureimg{allocation_heatmap}
    \caption{Vue Allocation : postes manquants et couverture RH.}
    \label{fig:allocation_heatmap}
\end{figure}

La vue Allocation montre ce que l'objectif ne dit pas. Elle comprend :
\begin{itemize}
    \item une \textbf{carte de chaleur} des postes manquants par service et par jour, filtrable par poste et par rôle ;
    \item un tableau de \textbf{couverture RH} par service, rôle et poste ;
    \item le \textbf{planning du personnel} du service : effectif affecté et cible de chaque poste, puis grille agent par jour. Les postes faits dans un autre service apparaissent en gris.
\end{itemize}

\begin{figure}[H]
    \centering
    \captureimg{allocation_planning}
    \caption{Vue Allocation : planning du personnel d'un service.}
    \label{fig:allocation_planning}
\end{figure}

\subsection{Simulateur de Scénario Personnalisé}
\begin{figure}[H]
    \centering
    \captureimg{simulateur}
    \caption{Vue Simulateur : réglages et écarts par rapport à $S_0$.}
    \label{fig:simulateur}
\end{figure}

L'utilisateur règle la demande ($\times 1$ à $\times 2$), les absences et les lits fermés (0 à 50\,\%), pour tout l'hôpital ou pour un service, puis choisit un profil et des méthodes. Naïve et Greedy sont cochées par défaut ; les solveurs sont facultatifs et affichent leur durée estimée. Le résultat est comparé à $S_0$, pour le même profil et la même méthode, sur quatre indicateurs.

Deux protections ont été ajoutées :
\begin{itemize}
    \item les simulations sont \textbf{isolées} : identifiant \texttt{CUSTOM\_}, stockage local ignoré par Git, lecture séparée par le paramètre \texttt{scope=custom} ;
    \item un cas \textbf{irréalisable} est signalé avant et après le calcul, et ses écarts ne sont pas présentés comme comparables.
\end{itemize}

\subsection{Vue Sensibilité}
La vue Sensibilité lit la grille de \NbCas{} cas : filtres, courbes des postes couverts selon les absences, courbes de coût selon la demande, carte de chaleur absences $\times$ demande et comparaison de Naïve et Greedy (figure~\ref{fig:sensibilite}). Elle ne lance aucun calcul.

% ------------------------------------------------------------------------------
\section*{Conclusion du Chapitre}
\addcontentsline{toc}{section}{Conclusion du Chapitre}
% ------------------------------------------------------------------------------
Le prototype couvre toute la chaîne, des données au tableau de bord. Les points délicats ont porté sur l'interprétation honnête des solveurs (statut, borne, reprise), sur un contrôle indépendant des plans et sur un affichage qui distingue contraintes respectées et postes couverts. Le chapitre suivant présente les tests et les résultats.

\chapter{Tests, Validation et Analyse des Résultats}
\label{ch:tests_validation_resultats}

La valeur d'un outil d'aide à la décision dépend de la confiance que l'on peut accorder à ses résultats, et de la clarté avec laquelle il en montre les limites.

Ce chapitre présente les tests, la validation des prévisions, la comparaison des quatre méthodes sur six scénarios et trois profils, l'arbitrage entre accueil et personnel, l'analyse de sensibilité, puis les limites. Tous les tableaux de résultats sont générés depuis le dossier \texttt{results/}.

% ------------------------------------------------------------------------------
\section{Protocole de Validation et Tests}
% ------------------------------------------------------------------------------

\subsection{Protocole Expérimental}
L'évaluation comprend trois campagnes :
\begin{enumerate}
    \item \textbf{Prévision} : cinq modèles, \NbPrevVal{} prévisions de validation et \NbPrevTest{} prévisions de test par service ;
    \item \textbf{Comparaison des méthodes} : 6 scénarios $\times$ 3 profils $\times$ 4 méthodes, soit 72 plans, avec la même instance pour les quatre méthodes d'un scénario. Chaque solveur dispose de \TempsSolveur{} secondes et de \ThreadsSolveur{} threads ;
    \item \textbf{Sensibilité} : \NbCas{} cas, croisant 5 niveaux d'absence (0, 5, 10, 15 et 20\,\%), 4 niveaux de demande ($\times 1$, $\times 1{,}25$, $\times 1{,}5$, $\times 2$) et 4 niveaux de fermeture (0, 5, 10 et 20\,\%), calculés avec Naïve et Greedy, profil équilibré.
\end{enumerate}

\subsection{Tests Automatisés}
Le projet compte \NbTestsPython{} tests Python et \NbTestsNavigateur{} parcours Playwright (tableau~\ref{tab:tests}).

% Généré par rapport/scripts/generate_tables.py : ne pas modifier à la main.
\begin{table}[H]
    \centering
    \caption{Tests automatisés du projet.}
    \label{tab:tests}
    \resizebox{\textwidth}{!}{\begin{tabular}{lcl}
        \toprule
        \textbf{Fichier} & \textbf{Tests} & \textbf{Objet} \\
        \midrule
        \texttt{test\_api\_contract.py} & 8 & Contrat de l'API, lecture des artefacts, traversée de chemin \\
        \texttt{test\_config\_contracts.py} & 5 & Validation de la configuration \\
        \texttt{test\_data\_validation.py} & 1 & Contrôles de qualité des données \\
        \texttt{test\_forecasting.py} & 5 & Prévisions chronologiques, métriques, absence de fuite \\
        \texttt{test\_generator.py} & 3 & Reproductibilité du générateur \\
        \texttt{test\_greedy\_and\_validator.py} & 10 & Greedy, validateur, occupation du premier jour \\
        \texttt{test\_occupancy.py} & 1 & Conservation des patients \\
        \texttt{test\_operational\_requirements.py} & 6 & Instance opérationnelle et empreinte de demande \\
        \texttt{test\_portability.py} & 5 & Chemins relatifs, empreinte indépendante du système \\
        \texttt{test\_profiles\_and\_naive.py} & 16 & Stratégie naïve, profils, équité sans les absents \\
        \texttt{test\_scenario\_comparison.py} & 4 & Comparaison par graine et par profil \\
        \texttt{test\_scenarios.py} & 4 & Transformations des scénarios \\
        \texttt{test\_sensitivity.py} & 5 & Grille de sensibilité et fusion des résultats \\
        \texttt{test\_simulator.py} & 11 & Scénarios personnalisés, isolement de S0--S5 \\
        \texttt{test\_solver\_settings.py} & 9 & Budget, borne, gap et statut prouvé des solveurs \\
        \texttt{test\_solver\_smoke.py} & 2 & Disponibilité des solveurs \\
        \midrule
        \textbf{Total Python (pytest)} & \textbf{95} &  \\
        \midrule
        \texttt{dashboard.spec.ts} & 1 & Ouverture des dix vues, sans calcul automatique \\
        \texttt{simulator.spec.ts} & 3 & Simulateur, cas irréalisable, comparaison et profils \\
        \midrule
        \textbf{Total navigateur (Playwright)} & \textbf{4} &  \\
        \bottomrule
    \end{tabular}}
\end{table}


Les tests du validateur lui soumettent notamment des plans volontairement incorrects. Les parcours Playwright utilisent une copie des résultats : ils ne modifient pas les expériences du projet.

% ------------------------------------------------------------------------------
\section{Validation des Prévisions}
% ------------------------------------------------------------------------------

% Généré par rapport/scripts/generate_tables.py : ne pas modifier à la main.
\begin{table}[H]
    \centering
    \caption{Erreurs de validation des cinq modèles (84 prévisions par service et par modèle).}
    \label{tab:previsions_validation}
    \resizebox{\textwidth}{!}{\begin{tabular}{llccccc}
        \toprule
        \textbf{Service} & \textbf{Modèle} & \textbf{MAE} & \textbf{RMSE} & \textbf{WAPE} & \textbf{Biais} & \textbf{Décision} \\
        \midrule
        \multirow{5}{*}{Urgences (\texttt{URG})} & Naïf & 7{,}64 & 9{,}11 & 21{,}93\,\% & $+$1{,}64 & -- \\
         & Naïf saisonnier & 7{,}20 & 9{,}23 & 20{,}66\,\% & $-$1{,}11 & -- \\
         & Moyenne mobile & 5{,}82 & 7{,}25 & 16{,}68\,\% & $-$0{,}46 & -- \\
         & Lissage exponentiel & 4{,}92 & 6{,}38 & 14{,}11\,\% & $-$1{,}46 & -- \\
         & \textbf{Régression à retards} & \textbf{4{,}91} & \textbf{6{,}36} & \textbf{14{,}08\,\%} & \textbf{$-$1{,}64} & \textbf{Retenu} \\
        \midrule
        \multirow{5}{*}{Médecine (\texttt{MED})} & Naïf & 4{,}76 & 6{,}07 & 26{,}37\,\% & $-$0{,}48 & -- \\
         & Naïf saisonnier & 4{,}86 & 6{,}03 & 26{,}90\,\% & $-$0{,}19 & -- \\
         & Moyenne mobile & 3{,}46 & 4{,}33 & 19{,}18\,\% & $-$0{,}22 & -- \\
         & Lissage exponentiel & 3{,}15 & 3{,}86 & 17{,}47\,\% & $+$0{,}17 & -- \\
         & \textbf{Régression à retards} & \textbf{3{,}12} & \textbf{3{,}83} & \textbf{17{,}25\,\%} & \textbf{$+$0{,}18} & \textbf{Retenu} \\
        \midrule
        \multirow{5}{*}{Réanimation (\texttt{REA})} & Naïf & 2{,}51 & 3{,}31 & 45{,}67\,\% & $+$0{,}08 & -- \\
         & Naïf saisonnier & 2{,}83 & 3{,}47 & 51{,}52\,\% & $+$0{,}02 & -- \\
         & Moyenne mobile & 2{,}24 & 2{,}88 & 40{,}74\,\% & $+$0{,}03 & -- \\
         & Lissage exponentiel & 2{,}09 & 2{,}73 & 37{,}97\,\% & $+$0{,}06 & -- \\
         & \textbf{Régression à retards} & \textbf{2{,}05} & \textbf{2{,}69} & \textbf{37{,}26\,\%} & \textbf{$+$0{,}07} & \textbf{Retenu} \\
        \midrule
        \multirow{5}{*}{Chirurgie (\texttt{CHI})} & Naïf & 5{,}81 & 7{,}37 & 43{,}92\,\% & $+$0{,}02 & -- \\
         & Naïf saisonnier & 3{,}98 & 4{,}77 & 30{,}06\,\% & $-$0{,}95 & -- \\
         & Moyenne mobile & 4{,}38 & 5{,}25 & 33{,}10\,\% & $-$1{,}02 & -- \\
         & Lissage exponentiel & 4{,}64 & 5{,}48 & 35{,}05\,\% & $-$0{,}98 & -- \\
         & \textbf{Régression à retards} & \textbf{3{,}41} & \textbf{4{,}29} & \textbf{25{,}78\,\%} & \textbf{$-$0{,}98} & \textbf{Retenu} \\
        \bottomrule
    \end{tabular}}
\end{table}


\begin{figure}[H]
    \centering
    \captureimg{model_comparison}
    \caption{Comparaison du WAPE de validation des cinq modèles par service.}
    \label{fig:model_comparison}
\end{figure}

La régression à retards a la plus faible erreur de validation dans les quatre services : \WapeValUrg{} aux Urgences, \WapeValMed{} en Médecine, \WapeValRea{} en Réanimation et \WapeValChi{} en Chirurgie. L'avance sur le lissage exponentiel est très faible aux Urgences et en Médecine ; elle est nette en Chirurgie, dont l'activité dépend du jour de la semaine. En Réanimation, les faibles volumes rendent l'erreur relative élevée.

% Généré par rapport/scripts/generate_tables.py : ne pas modifier à la main.
\begin{table}[H]
    \centering
    \caption{MAE de la régression à retards selon l'horizon (validation, 12 origines).}
    \label{tab:previsions_horizon}
    \small
    \begin{tabular}{lccccccc}
        \toprule
        \textbf{Service} & \textbf{$J+1$} & \textbf{$J+2$} & \textbf{$J+3$} & \textbf{$J+4$} & \textbf{$J+5$} & \textbf{$J+6$} & \textbf{$J+7$} \\
        \midrule
        Urgences & 5{,}99 & 5{,}74 & 3{,}65 & 3{,}97 & 4{,}12 & 6{,}37 & 4{,}52 \\
        Médecine & 2{,}92 & 3{,}04 & 2{,}89 & 3{,}29 & 3{,}30 & 3{,}21 & 3{,}15 \\
        Réanimation & 1{,}74 & 1{,}37 & 1{,}79 & 1{,}75 & 2{,}53 & 1{,}70 & 3{,}46 \\
        Chirurgie & 3{,}63 & 2{,}49 & 4{,}16 & 3{,}91 & 3{,}18 & 3{,}59 & 2{,}92 \\
        \bottomrule
    \end{tabular}
\end{table}


L'erreur ne croît pas régulièrement avec l'horizon (tableau~\ref{tab:previsions_horizon}), mais chaque valeur ne repose que sur 12 prévisions.

\subsection{Test Final}
% Généré par rapport/scripts/generate_tables.py : ne pas modifier à la main.
\begin{table}[H]
    \centering
    \caption{WAPE sur le test final réservé (28 prévisions par service), jamais utilisé pour choisir le modèle.}
    \label{tab:previsions_test}
    \resizebox{\textwidth}{!}{\begin{tabular}{lccccc}
        \toprule
        \textbf{Service} & \textbf{Naïf} & \textbf{Naïf saisonnier} & \textbf{Moyenne mobile} & \textbf{Lissage exponentiel} & \textbf{Régression à retards} \\
        \midrule
        Urgences & 22{,}1\,\% & 20{,}3\,\% & 17{,}1\,\% & 16{,}3\,\% & \textbf{15{,}8\,\%}$^{\ast}$ \\
        Médecine & 22{,}6\,\% & 25{,}8\,\% & \textbf{18{,}9\,\%} & 20{,}4\,\% & 19{,}9\,\%$^{\ast}$ \\
        Réanimation & 77{,}7\,\% & 55{,}3\,\% & \textbf{33{,}8\,\%} & 34{,}8\,\% & 35{,}4\,\%$^{\ast}$ \\
        Chirurgie & 41{,}8\,\% & \textbf{28{,}1\,\%} & 34{,}8\,\% & 35{,}9\,\% & 28{,}7\,\%$^{\ast}$ \\
        \bottomrule
    \end{tabular}}
    \par\vspace{3pt}
    {\footnotesize $^{\ast}$ modèle retenu en validation ; en gras, meilleure valeur du test.}
\end{table}


Sur le test final, jamais utilisé pour choisir, la régression n'est la meilleure qu'aux Urgences. La moyenne mobile fait mieux en Médecine et en Réanimation, le naïf saisonnier en Chirurgie, avec des écarts limités. La sélection n'a pas été modifiée après coup, ce qui aurait invalidé ce test. Avec \NbPrevTest{} prévisions par service, ces écarts doivent être lus avec prudence.

% ------------------------------------------------------------------------------
\section{Comparaison des Quatre Méthodes}
% ------------------------------------------------------------------------------

\subsection{Respect des Contraintes et Temps de Calcul}
Les 72 plans respectent toutes les contraintes dures selon le validateur. Le tableau~\ref{tab:performances} résume les temps et les gaps.

% Généré par rapport/scripts/generate_tables.py : ne pas modifier à la main.
\begin{table}[H]
    \centering
    \caption{Temps de calcul, gap déclaré et respect des contraintes (18 plans par méthode).}
    \label{tab:performances}
    \resizebox{\textwidth}{!}{\begin{tabular}{lcccccc}
        \toprule
        \textbf{Méthode} & \textbf{Min (s)} & \textbf{Moyenne (s)} & \textbf{Max (s)} & \textbf{Gap} & \textbf{Optimal prouvé} & \textbf{Contraintes respectées} \\
        \midrule
        Naïve & 0{,}3 & 0{,}5 & 1{,}4 & -- & 0 & 18 / 18 \\
        Greedy & 0{,}3 & 0{,}5 & 1{,}3 & -- & 0 & 18 / 18 \\
        MILP/CBC & 7{,}2 & 51{,}1 & 62{,}6 & 0{,}0\,\% à 19{,}2\,\% & 1 & 18 / 18 \\
        CP-SAT & 43{,}7 & 51{,}5 & 53{,}1 & 0{,}0\,\% à 4{,}5\,\% & 1 & 18 / 18 \\
        \bottomrule
    \end{tabular}}
\end{table}


Naïve et Greedy répondent en moins de deux secondes. Les solveurs utilisent presque toujours leur budget entier ; le temps mesuré inclut aussi la construction du modèle et le contrôle du résultat, d'où des valeurs supérieures à \TempsSolveur{} secondes. Un seul cas est prouvé optimal, par les deux solveurs : $S_3$ avec la priorité accueil. Les autres plans des solveurs sont des solutions réalisables, sans preuve d'optimalité.

\subsection{Profil Équilibré}
% Généré par rapport/scripts/generate_tables.py : ne pas modifier à la main.
\begin{table}[H]
    \centering
    \caption{Résultats des quatre méthodes, profil « Équilibré ».}
    \label{tab:resultats_equilibre}
    \resizebox{\textwidth}{!}{\begin{tabular}{llrrrrrrrr}
        \toprule
        \textbf{Scénario} & \textbf{Méthode} & \textbf{Objectif $\mathcal{Z}$} & \textbf{vs Greedy} & \textbf{Gap} & \textbf{Refus} & \textbf{Couv. demande} & \textbf{Déficit RH} & \textbf{Postes couverts} & \textbf{Équité} \\
        \midrule
        \multirow{4}{*}{$S_0$ Nominal} & Naïve & 718\,885 & $+$13{,}0\,\% & -- & 4 & 99{,}0\,\% & 1\,365 & 41{,}2\,\% & 15 \\
         & Greedy & 636\,379 & -- & -- & 4 & 99{,}0\,\% & 1\,200 & 48{,}3\,\% & 9 \\
         & MILP/CBC & 609\,870 & $-$4{,}2\,\% & 2{,}7\,\% & 4 & 99{,}0\,\% & 1\,147 & 50{,}6\,\% & 0 \\
         & CP-SAT & \textbf{604\,800} & $-$5{,}0\,\% & 2{,}8\,\% & 9 & 98{,}7\,\% & 1\,123 & 51{,}1\,\% & 0 \\
        \midrule
        \multirow{4}{*}{$S_1$ Pic aux urgences} & Naïve & 787\,965 & $+$13{,}3\,\% & -- & 31 & 95{,}9\,\% & 1\,383 & 40{,}9\,\% & 15 \\
         & Greedy & 695\,559 & -- & -- & 29 & 96{,}5\,\% & 1\,218 & 47{,}9\,\% & 9 \\
         & MILP/CBC & 666\,041 & $-$4{,}2\,\% & 2{,}3\,\% & 30 & 96{,}4\,\% & 1\,157 & 50{,}4\,\% & 1 \\
         & CP-SAT & \textbf{661\,490} & $-$4{,}9\,\% & 2{,}3\,\% & 33 & 96{,}2\,\% & 1\,138 & 50{,}8\,\% & 0 \\
        \midrule
        \multirow{4}{*}{$S_2$ Hausse générale} & Naïve & 1\,010\,265 & $+$9{,}3\,\% & -- & 74 & 90{,}7\,\% & 1\,566 & 38{,}3\,\% & 15 \\
         & Greedy & 924\,327 & -- & -- & 73 & 91{,}0\,\% & 1\,408 & 44{,}6\,\% & 7 \\
         & MILP/CBC & 903\,320 & $-$2{,}3\,\% & 1{,}1\,\% & 73 & 91{,}0\,\% & 1\,366 & 46{,}2\,\% & 0 \\
         & CP-SAT & \textbf{900\,740} & $-$2{,}6\,\% & 1{,}5\,\% & 76 & 90{,}7\,\% & 1\,345 & 46{,}6\,\% & 0 \\
        \midrule
        \multirow{4}{*}{$S_3$ Absences} & Naïve & 783\,385 & $+$9{,}9\,\% & -- & 4 & 99{,}0\,\% & 1\,494 & 35{,}7\,\% & 15 \\
         & Greedy & 712\,876 & -- & -- & 4 & 99{,}0\,\% & 1\,353 & 41{,}7\,\% & 6 \\
         & MILP/CBC & 696\,200 & $-$2{,}3\,\% & 2{,}4\,\% & 9 & 98{,}1\,\% & 1\,286 & 43{,}7\,\% & 0 \\
         & CP-SAT & \textbf{691\,790} & $-$3{,}0\,\% & 2{,}2\,\% & 9 & 98{,}6\,\% & 1\,295 & 43{,}6\,\% & 0 \\
        \midrule
        \multirow{4}{*}{$S_4$ Crise combinée} & Naïve & 1\,083\,265 & $+$8{,}1\,\% & -- & 74 & 90{,}7\,\% & 1\,712 & 32{,}5\,\% & 15 \\
         & Greedy & 1\,002\,325 & -- & -- & 73 & 91{,}0\,\% & 1\,564 & 38{,}4\,\% & 5 \\
         & MILP/CBC & \textbf{988\,260} & $-$1{,}4\,\% & 0{,}7\,\% & 73 & 90{,}8\,\% & 1\,524 & 39{,}6\,\% & 0 \\
         & CP-SAT & 988\,270 & $-$1{,}4\,\% & 1{,}2\,\% & 76 & 90{,}8\,\% & 1\,526 & 39{,}6\,\% & 0 \\
        \midrule
        \multirow{4}{*}{$S_5$ Fermeture de lits} & Naïve & 757\,455 & $+$12{,}3\,\% & -- & 27 & 94{,}0\,\% & 1\,258 & 43{,}1\,\% & 15 \\
         & Greedy & 674\,469 & -- & -- & 26 & 94{,}1\,\% & 1\,096 & 50{,}6\,\% & 9 \\
         & MILP/CBC & 650\,444 & $-$3{,}6\,\% & 3{,}0\,\% & 27 & 94{,}0\,\% & 1\,044 & 52{,}8\,\% & 4 \\
         & CP-SAT & \textbf{641\,250} & $-$4{,}9\,\% & 2{,}7\,\% & 35 & 93{,}0\,\% & 988 & 54{,}3\,\% & 0 \\
        \bottomrule
    \end{tabular}}
    \par\vspace{3pt}
    {\footnotesize En gras : meilleur objectif du scénario. Les 24 plans respectent toutes les contraintes dures.}
\end{table}


Avec le profil équilibré, la meilleure méthode améliore l'objectif de Greedy de \GainMinEq{} à \GainMaxEq{}. CP-SAT est la meilleure dans cinq scénarios ; MILP l'emporte en $S_4$, avec un écart de 10 points sur l'objectif. La stratégie naïve est plus coûteuse que Greedy dans tous les cas ; sur les trois profils, l'écart va de \NaiveMin{} à \NaiveMax{}.

Le déficit RH représente \PartRhMin{} à \PartRhMax{} de l'objectif selon le plan (tableau~\ref{tab:parts_objectif}). L'objectif équilibré mesure donc surtout la tension sur le personnel.

% Généré par rapport/scripts/generate_tables.py : ne pas modifier à la main.
\begin{table}[H]
    \centering
    \caption{Part de chaque composante dans l'objectif (24 plans par profil).}
    \label{tab:parts_objectif}
    \resizebox{\textwidth}{!}{\begin{tabular}{lcccc}
        \toprule
        \textbf{Profil} & \textbf{Déficit RH} & \textbf{Demande non soutenue} & \textbf{Refus} & \textbf{Lits ouverts} \\
        \midrule
        Équilibré & 74{,}7\,\% à 95{,}4\,\% & 2{,}3\,\% à 23{,}1\,\% & 0{,}0\,\% à 0{,}0\,\% & 1{,}9\,\% à 3{,}0\,\% \\
        Priorité accueil & 33{,}9\,\% à 91{,}1\,\% & 5{,}2\,\% à 38{,}5\,\% & 1{,}9\,\% à 26{,}8\,\% & 0{,}7\,\% à 2{,}2\,\% \\
        Priorité personnel & 54{,}2\,\% à 96{,}7\,\% & 1{,}6\,\% à 44{,}3\,\% & 0{,}0\,\% à 0{,}0\,\% & 1{,}2\,\% à 2{,}1\,\% \\
        \bottomrule
    \end{tabular}}
\end{table}


\subsection{Profil Priorité Accueil}
% Généré par rapport/scripts/generate_tables.py : ne pas modifier à la main.
\begin{table}[H]
    \centering
    \caption{Résultats des quatre méthodes, profil « Priorité accueil ».}
    \label{tab:resultats_priorite_accueil}
    \resizebox{\textwidth}{!}{\begin{tabular}{llrrrrrrrr}
        \toprule
        \textbf{Scénario} & \textbf{Méthode} & \textbf{Objectif $\mathcal{Z}$} & \textbf{vs Greedy} & \textbf{Gap} & \textbf{Refus} & \textbf{Couv. demande} & \textbf{Déficit RH} & \textbf{Postes couverts} & \textbf{Équité} \\
        \midrule
        \multirow{4}{*}{$S_0$ Nominal} & Naïve & 964\,385 & $+$13{,}9\,\% & -- & 4 & 99{,}0\,\% & 1\,365 & 41{,}2\,\% & 15 \\
         & Greedy & 846\,879 & -- & -- & 4 & 99{,}0\,\% & 1\,200 & 48{,}3\,\% & 9 \\
         & MILP/CBC & \textbf{820\,370} & $-$3{,}1\,\% & 0{,}0\,\% & 4 & 99{,}0\,\% & 1\,147 & 50{,}6\,\% & 0 \\
         & CP-SAT & \textbf{820\,370} & $-$3{,}1\,\% & 0{,}5\,\% & 4 & 99{,}0\,\% & 1\,147 & 50{,}6\,\% & 0 \\
        \midrule
        \multirow{4}{*}{$S_1$ Pic aux urgences} & Naïve & 1\,847\,465 & $+$9{,}3\,\% & -- & 31 & 95{,}9\,\% & 1\,383 & 40{,}9\,\% & 15 \\
         & Greedy & 1\,690\,059 & -- & -- & 29 & 96{,}5\,\% & 1\,218 & 47{,}9\,\% & 9 \\
         & MILP/CBC & \textbf{1\,657\,550} & $-$1{,}9\,\% & 0{,}0\,\% & 29 & 96{,}5\,\% & 1\,159 & 50{,}3\,\% & 0 \\
         & CP-SAT & \textbf{1\,657\,550} & $-$1{,}9\,\% & 0{,}2\,\% & 29 & 96{,}5\,\% & 1\,159 & 50{,}3\,\% & 0 \\
        \midrule
        \multirow{4}{*}{$S_2$ Hausse générale} & Naïve & 2\,714\,265 & $+$5{,}4\,\% & -- & 74 & 90{,}7\,\% & 1\,566 & 38{,}3\,\% & 15 \\
         & Greedy & 2\,576\,327 & -- & -- & 73 & 91{,}0\,\% & 1\,408 & 44{,}6\,\% & 7 \\
         & MILP/CBC & \textbf{2\,555\,320} & $-$0{,}8\,\% & 0{,}0\,\% & 73 & 91{,}0\,\% & 1\,366 & 46{,}2\,\% & 0 \\
         & CP-SAT & \textbf{2\,555\,320} & $-$0{,}8\,\% & 0{,}1\,\% & 73 & 91{,}0\,\% & 1\,366 & 46{,}2\,\% & 0 \\
        \midrule
        \multirow{4}{*}{$S_3$ Absences} & Naïve & 1\,040\,885 & $+$10{,}0\,\% & -- & 4 & 99{,}0\,\% & 1\,494 & 35{,}7\,\% & 15 \\
         & Greedy & 945\,876 & -- & -- & 4 & 99{,}0\,\% & 1\,353 & 41{,}7\,\% & 6 \\
         & MILP/CBC & \textbf{930\,370} & $-$1{,}6\,\% & 0{,}0\,\% & 4 & 99{,}0\,\% & 1\,322 & 43{,}1\,\% & 0 \\
         & CP-SAT & \textbf{930\,370} & $-$1{,}6\,\% & 0{,}0\,\% & 4 & 99{,}0\,\% & 1\,322 & 43{,}1\,\% & 0 \\
        \midrule
        \multirow{4}{*}{$S_4$ Crise combinée} & Naïve & 2\,800\,765 & $+$4{,}6\,\% & -- & 74 & 90{,}7\,\% & 1\,712 & 32{,}5\,\% & 15 \\
         & Greedy & 2\,676\,825 & -- & -- & 73 & 91{,}0\,\% & 1\,564 & 38{,}4\,\% & 5 \\
         & MILP/CBC & \textbf{2\,665\,320} & $-$0{,}4\,\% & 0{,}0\,\% & 73 & 91{,}0\,\% & 1\,541 & 39{,}4\,\% & 0 \\
         & CP-SAT & \textbf{2\,665\,320} & $-$0{,}4\,\% & 0{,}1\,\% & 73 & 91{,}0\,\% & 1\,541 & 39{,}4\,\% & 0 \\
        \midrule
        \multirow{4}{*}{$S_5$ Fermeture de lits} & Naïve & 1\,402\,455 & $+$11{,}3\,\% & -- & 27 & 94{,}0\,\% & 1\,258 & 43{,}1\,\% & 15 \\
         & Greedy & 1\,260\,469 & -- & -- & 26 & 94{,}1\,\% & 1\,096 & 50{,}6\,\% & 9 \\
         & MILP/CBC & \textbf{1\,233\,460} & $-$2{,}1\,\% & 0{,}1\,\% & 26 & 94{,}1\,\% & 1\,042 & 53{,}0\,\% & 0 \\
         & CP-SAT & \textbf{1\,233\,460} & $-$2{,}1\,\% & 0{,}4\,\% & 26 & 94{,}1\,\% & 1\,042 & 53{,}0\,\% & 0 \\
        \bottomrule
    \end{tabular}}
    \par\vspace{3pt}
    {\footnotesize En gras : meilleur objectif du scénario. Les 24 plans respectent toutes les contraintes dures.}
\end{table}


MILP et CP-SAT atteignent le même objectif dans les six scénarios, avec un gap d'au plus \GapMaxAccueilMilp{} pour MILP et \GapMaxAccueilCp{} pour CP-SAT. Ils refusent autant de patients que Greedy : les refus restants sont \textbf{imposés par les lits}. Le gain (\GainMinAcc{} à \GainMaxAcc{}) vient d'une meilleure affectation du personnel, à accueil égal. Ce profil est le seul où le budget suffit à conclure presque exactement.

\subsection{Profil Priorité Personnel}
% Généré par rapport/scripts/generate_tables.py : ne pas modifier à la main.
\begin{table}[H]
    \centering
    \caption{Résultats des quatre méthodes, profil « Priorité personnel ».}
    \label{tab:resultats_priorite_personnel}
    \resizebox{\textwidth}{!}{\begin{tabular}{llrrrrrrrr}
        \toprule
        \textbf{Scénario} & \textbf{Méthode} & \textbf{Objectif $\mathcal{Z}$} & \textbf{vs Greedy} & \textbf{Gap} & \textbf{Refus} & \textbf{Couv. demande} & \textbf{Déficit RH} & \textbf{Postes couverts} & \textbf{Équité} \\
        \midrule
        \multirow{4}{*}{$S_0$ Nominal} & Naïve & 1\,061\,620 & $+$13{,}3\,\% & -- & 4 & 99{,}0\,\% & 1\,365 & 41{,}2\,\% & 15 \\
         & Greedy & 937\,270 & -- & -- & 4 & 99{,}0\,\% & 1\,200 & 48{,}3\,\% & 9 \\
         & MILP/CBC & 999\,970 & $+$6{,}7\,\% & 16{,}8\,\% & 169 & 76{,}1\,\% & 723 & 59{,}6\,\% & 6 \\
         & CP-SAT & \textbf{865\,660} & $-$7{,}6\,\% & 4{,}4\,\% & 27 & 96{,}5\,\% & 1\,045 & 52{,}9\,\% & 0 \\
        \midrule
        \multirow{4}{*}{$S_1$ Pic aux urgences} & Naïve & 1\,135\,200 & $+$13{,}4\,\% & -- & 31 & 95{,}9\,\% & 1\,383 & 40{,}9\,\% & 15 \\
         & Greedy & 1\,000\,950 & -- & -- & 29 & 96{,}5\,\% & 1\,218 & 47{,}9\,\% & 9 \\
         & MILP/CBC & 963\,100 & $-$3{,}8\,\% & 6{,}8\,\% & 34 & 95{,}9\,\% & 1\,155 & 50{,}1\,\% & 4 \\
         & CP-SAT & \textbf{928\,710} & $-$7{,}2\,\% & 3{,}9\,\% & 61 & 93{,}4\,\% & 1\,045 & 52{,}9\,\% & 0 \\
        \midrule
        \multirow{4}{*}{$S_2$ Hausse générale} & Naïve & 1\,403\,250 & $+$9{,}9\,\% & -- & 74 & 90{,}7\,\% & 1\,566 & 38{,}3\,\% & 15 \\
         & Greedy & 1\,277\,020 & -- & -- & 73 & 91{,}0\,\% & 1\,408 & 44{,}6\,\% & 7 \\
         & MILP/CBC & 1\,237\,200 & $-$3{,}1\,\% & 4{,}3\,\% & 81 & 90{,}3\,\% & 1\,333 & 46{,}5\,\% & 3 \\
         & CP-SAT & \textbf{1\,223\,760} & $-$4{,}2\,\% & 3{,}5\,\% & 103 & 88{,}5\,\% & 1\,264 & 48{,}2\,\% & 0 \\
        \midrule
        \multirow{4}{*}{$S_3$ Absences} & Naïve & 1\,158\,370 & $+$10{,}1\,\% & -- & 4 & 99{,}0\,\% & 1\,494 & 35{,}7\,\% & 15 \\
         & Greedy & 1\,051\,720 & -- & -- & 4 & 99{,}0\,\% & 1\,353 & 41{,}7\,\% & 6 \\
         & MILP/CBC & 1\,187\,540 & $+$12{,}9\,\% & 19{,}2\,\% & 221 & 75{,}4\,\% & 955 & 46{,}4\,\% & 3 \\
         & CP-SAT & \textbf{997\,810} & $-$5{,}1\,\% & 4{,}3\,\% & 32 & 96{,}0\,\% & 1\,208 & 45{,}3\,\% & 0 \\
        \midrule
        \multirow{4}{*}{$S_4$ Crise combinée} & Naïve & 1\,512\,750 & $+$8{,}5\,\% & -- & 74 & 90{,}7\,\% & 1\,712 & 32{,}5\,\% & 15 \\
         & Greedy & 1\,393\,820 & -- & -- & 73 & 91{,}0\,\% & 1\,564 & 38{,}4\,\% & 5 \\
         & MILP/CBC & 1\,361\,520 & $-$2{,}3\,\% & 3{,}3\,\% & 126 & 81{,}9\,\% & 1\,251 & 44{,}0\,\% & 0 \\
         & CP-SAT & \textbf{1\,354\,380} & $-$2{,}8\,\% & 3{,}2\,\% & 113 & 87{,}9\,\% & 1\,421 & 41{,}3\,\% & 0 \\
        \midrule
        \multirow{4}{*}{$S_5$ Fermeture de lits} & Naïve & 1\,073\,440 & $+$13{,}1\,\% & -- & 27 & 94{,}0\,\% & 1\,258 & 43{,}1\,\% & 15 \\
         & Greedy & 949\,360 & -- & -- & 26 & 94{,}1\,\% & 1\,096 & 50{,}6\,\% & 9 \\
         & MILP/CBC & 887\,690 & $-$6{,}5\,\% & 5{,}5\,\% & 45 & 92{,}1\,\% & 966 & 54{,}8\,\% & 1 \\
         & CP-SAT & \textbf{874\,750} & $-$7{,}9\,\% & 4{,}5\,\% & 75 & 88{,}9\,\% & 871 & 57{,}4\,\% & 0 \\
        \bottomrule
    \end{tabular}}
    \par\vspace{3pt}
    {\footnotesize En gras : meilleur objectif du scénario. Les 24 plans respectent toutes les contraintes dures.}
\end{table}


CP-SAT est la meilleure méthode dans les six scénarios (gain de \GainMinPer{} à \GainMaxPer{}). MILP est peu fiable sous ce profil : il est moins bon que Greedy en $S_0$ (\MilpPerSz{}) et en $S_3$ (\MilpPerSt{}), avec des gaps de \GapMilpPerSz{} et \GapMilpPerSt{} et jusqu'à \RefusMilpPerSt{} refus. Dans ces deux cas, son plan est une solution intermédiaire, pas une recommandation.

\subsection{Synthèse}
% Généré par rapport/scripts/generate_tables.py : ne pas modifier à la main.
\begin{table}[H]
    \centering
    \caption{Meilleure méthode par profil et par scénario.}
    \label{tab:meilleures_methodes}
    \small
    \begin{tabular}{lclccc}
        \toprule
        \textbf{Profil} &  & \textbf{Meilleure méthode} & \textbf{Gain vs Greedy} & \textbf{Refus (Greedy)} & \textbf{Déficit RH (Greedy)} \\
        \midrule
        \multirow{6}{*}{Équilibré} & $S_0$ & CP-SAT & $-$5{,}0\,\% & 9 (4) & 1\,123 (1\,200) \\
         & $S_1$ & CP-SAT & $-$4{,}9\,\% & 33 (29) & 1\,138 (1\,218) \\
         & $S_2$ & CP-SAT & $-$2{,}6\,\% & 76 (73) & 1\,345 (1\,408) \\
         & $S_3$ & CP-SAT & $-$3{,}0\,\% & 9 (4) & 1\,295 (1\,353) \\
         & $S_4$ & MILP/CBC & $-$1{,}4\,\% & 73 (73) & 1\,524 (1\,564) \\
         & $S_5$ & CP-SAT & $-$4{,}9\,\% & 35 (26) & 988 (1\,096) \\
        \midrule
        \multirow{6}{*}{Priorité accueil} & $S_0$ & MILP/CBC = CP-SAT & $-$3{,}1\,\% & 4 (4) & 1\,147 (1\,200) \\
         & $S_1$ & MILP/CBC = CP-SAT & $-$1{,}9\,\% & 29 (29) & 1\,159 (1\,218) \\
         & $S_2$ & MILP/CBC = CP-SAT & $-$0{,}8\,\% & 73 (73) & 1\,366 (1\,408) \\
         & $S_3$ & MILP/CBC = CP-SAT & $-$1{,}6\,\% & 4 (4) & 1\,322 (1\,353) \\
         & $S_4$ & MILP/CBC = CP-SAT & $-$0{,}4\,\% & 73 (73) & 1\,541 (1\,564) \\
         & $S_5$ & MILP/CBC = CP-SAT & $-$2{,}1\,\% & 26 (26) & 1\,042 (1\,096) \\
        \midrule
        \multirow{6}{*}{Priorité personnel} & $S_0$ & CP-SAT & $-$7{,}6\,\% & 27 (4) & 1\,045 (1\,200) \\
         & $S_1$ & CP-SAT & $-$7{,}2\,\% & 61 (29) & 1\,045 (1\,218) \\
         & $S_2$ & CP-SAT & $-$4{,}2\,\% & 103 (73) & 1\,264 (1\,408) \\
         & $S_3$ & CP-SAT & $-$5{,}1\,\% & 32 (4) & 1\,208 (1\,353) \\
         & $S_4$ & CP-SAT & $-$2{,}8\,\% & 113 (73) & 1\,421 (1\,564) \\
         & $S_5$ & CP-SAT & $-$7{,}9\,\% & 75 (26) & 871 (1\,096) \\
        \bottomrule
    \end{tabular}
    \par\vspace{3pt}
    {\footnotesize Les objectifs de deux profils ne sont pas comparables entre eux.}
\end{table}


\begin{figure}[H]
    \centering
    \captureimg{comparaison_scenarios}
    \caption{Vue Comparaison scénarios : méthodes et scénarios côte à côte.}
    \label{fig:comparaison}
\end{figure}

CP-SAT est la meilleure méthode, ou ex æquo, dans \NbCpMeilleur{} cas sur 18. Greedy reste une référence solide : calculée en moins de deux secondes, elle n'est améliorée que de \GainMinAcc{} à \GainMaxPer{} selon le cas.

% ------------------------------------------------------------------------------
\section{Arbitrage entre Accueil et Personnel}
% ------------------------------------------------------------------------------

% Généré par rapport/scripts/generate_tables.py : ne pas modifier à la main.
\begin{table}[H]
    \centering
    \caption{Arbitrage accueil / personnel : refus et déficit RH de CP-SAT selon le profil.}
    \label{tab:arbitrage}
    \resizebox{\textwidth}{!}{\begin{tabular}{lcccccccc}
        \toprule
        \textbf{Scénario} & \textbf{Greedy refus} & \textbf{déficit} & \textbf{Équilibré refus} & \textbf{déficit} & \textbf{Accueil refus} & \textbf{déficit} & \textbf{Personnel refus} & \textbf{déficit} \\
        \midrule
        $S_0$ Nominal & 4 & 1\,200 & 9 & 1\,123 & 4 & 1\,147 & 27 & 1\,045 \\
        $S_1$ Pic aux urgences & 29 & 1\,218 & 33 & 1\,138 & 29 & 1\,159 & 61 & 1\,045 \\
        $S_2$ Hausse générale & 73 & 1\,408 & 76 & 1\,345 & 73 & 1\,366 & 103 & 1\,264 \\
        $S_3$ Absences & 4 & 1\,353 & 9 & 1\,295 & 4 & 1\,322 & 32 & 1\,208 \\
        $S_4$ Crise combinée & 73 & 1\,564 & 76 & 1\,526 & 73 & 1\,541 & 113 & 1\,421 \\
        $S_5$ Fermeture de lits & 26 & 1\,096 & 35 & 988 & 26 & 1\,042 & 75 & 871 \\
        \bottomrule
    \end{tabular}}
    \par\vspace{3pt}
    {\footnotesize Le plan Greedy est identique pour les trois profils.}
\end{table}


Le tableau~\ref{tab:arbitrage} est le résultat central du projet. Par rapport à Greedy, CP-SAT réduit le déficit RH :
\begin{itemize}
    \item de \DeficitCpMinAcc{} à \DeficitCpMaxAcc{} avec la priorité accueil, \textbf{sans aucun refus supplémentaire} ;
    \item de \DeficitCpMinEq{} à \DeficitCpMaxEq{} avec le profil équilibré, au prix de quelques refus (\RefusCpEqSz{} contre \RefusSzGreedy{} en $S_0$) ;
    \item de \DeficitCpMinPer{} à \DeficitCpMaxPer{} avec la priorité personnel, au prix de refus nettement plus nombreux (\RefusCpPerSz{} contre \RefusSzGreedy{} en $S_0$, \RefusCpPerSq{} contre \RefusGreedySq{} en $S_4$).
\end{itemize}

Un meilleur objectif n'est donc pas toujours un meilleur accueil. Les refus supplémentaires s'expliquent en grande partie par l'effet de palier décrit au chapitre~\ref{ch:conception_systeme} : refuser un patient permet de rester sous le seuil qui déclenche de nouveaux postes cibles. Le choix du profil relève de la direction de l'établissement, pas du solveur.

\subsection{Lecture par Scénario}
\begin{itemize}
    \item \textbf{$S_0$ -- nominal} : Greedy accepte \CouvDemSzGreedy{} de la demande (\RefusSzGreedy{} refus), mais ne couvre que \CouvSzGreedy{} des postes cibles ;
    \item \textbf{$S_1$ -- pic aux urgences} : \RefusGreedySu{} refus avec Greedy, le déficit RH augmente peu (\DeficitGreedySu{} contre \DeficitSzGreedy{}) ;
    \item \textbf{$S_2$ -- hausse générale} : les lits deviennent limitants, avec \RefusGreedySd{} refus ;
    \item \textbf{$S_3$ -- absences} : \AbsentsSt{} agents retirés. Les refus ne changent pas (\RefusGreedySt{}), mais la couverture RH baisse à \CouvGreedySt{} ;
    \item \textbf{$S_4$ -- crise combinée} : autant de refus que $S_2$ pour Greedy (\RefusGreedySq{}) et couverture RH de \CouvGreedySq{} ; \SaturesSq{} couples service-jour sont saturés ;
    \item \textbf{$S_5$ -- fermeture de lits} : \RefusGreedySc{} refus, mais le déficit RH baisse (\DeficitGreedySc{}), car moins de lits demandent moins de postes.
\end{itemize}

\subsection{Couverture RH par Service}
% Généré par rapport/scripts/generate_tables.py : ne pas modifier à la main.
\begin{table}[H]
    \centering
    \caption{Postes cibles couverts en $S_0$ par service et par poste, profil équilibré.}
    \label{tab:couverture_s0}
    \small
    \begin{tabular}{lcccc|ccc|c}
        \toprule
        \textbf{Méthode} & \textbf{URG} & \textbf{MED} & \textbf{REA} & \textbf{CHI} & \textbf{Matin} & \textbf{Soir} & \textbf{Nuit} & \textbf{Total} \\
        \midrule
        Naïve & 36\,\% & 35\,\% & 49\,\% & 51\,\% & 43\,\% & 44\,\% & 37\,\% & \textbf{41{,}2\,\%} \\
        Greedy & 48\,\% & 20\,\% & 90\,\% & 64\,\% & 61\,\% & 50\,\% & 33\,\% & \textbf{48{,}3\,\%} \\
        MILP/CBC & 30\,\% & 26\,\% & 91\,\% & 86\,\% & 53\,\% & 49\,\% & 50\,\% & \textbf{50{,}6\,\%} \\
        CP-SAT & 45\,\% & 25\,\% & 93\,\% & 70\,\% & 59\,\% & 54\,\% & 41\,\% & \textbf{51{,}1\,\%} \\
        \bottomrule
    \end{tabular}
\end{table}


La couverture est très inégale (tableau~\ref{tab:couverture_s0}). Greedy sert d'abord la Réanimation (\CouvSzGreedyRea{}) et laisse la Médecine à \CouvSzGreedyMed{}. MILP, avec le profil équilibré, descend aux Urgences à \CouvSzMilpUrg{} ; avec la priorité accueil, il y remonte à \CouvSzMilpUrgAccueil{}. Dans le plan équilibré de CP-SAT, la nuit aux Urgences ne couvre que \CouvNuitInfUrgCp{} des postes infirmiers cibles. Un taux global ne suffit donc pas : la vue Allocation est nécessaire pour lire un plan.

% ------------------------------------------------------------------------------
\section{Analyse de Sensibilité (Grille de 80 Cas)}
% ------------------------------------------------------------------------------

La grille donne \NbResSens{} résultats (\NbCas{} cas, Naïve et Greedy). \NbResSensOk{} respectent toutes les contraintes ; les \NbResSensKo{} autres correspondent tous à une fermeture de 20\,\% des lits. Le seuil vient de la Médecine : au-delà d'environ \LimiteMed{} de lits fermés, les patients déjà présents le premier jour dépassent les lits restants, quelle que soit la méthode. La liste complète figure en annexe C.

\begin{figure}[H]
    \centering
    \captureimg{sensibilite}
    \caption{Vue Sensibilité : effet des absences et de la demande.}
    \label{fig:sensibilite}
\end{figure}

\subsection{Effet sur les Postes Couverts}
% Généré par rapport/scripts/generate_tables.py : ne pas modifier à la main.
\begin{table}[H]
    \centering
    \caption{Postes cibles couverts par Greedy selon les absences et la demande (aucune fermeture de lits).}
    \label{tab:sensibilite_couverture}
    \small
    \begin{tabular}{lcccc}
        \toprule
        \textbf{Absences} & \textbf{Demande $\times$1{,}00} & \textbf{Demande $\times$1{,}25} & \textbf{Demande $\times$1{,}50} & \textbf{Demande $\times$2{,}00} \\
        \midrule
        0\,\% & 48{,}3\,\% & 44{,}6\,\% & 43{,}9\,\% & 43{,}6\,\% \\
        5\,\% & 46{,}2\,\% & 42{,}8\,\% & 42{,}1\,\% & 41{,}9\,\% \\
        10\,\% & 44{,}1\,\% & 40{,}8\,\% & 40{,}1\,\% & 39{,}9\,\% \\
        15\,\% & 41{,}7\,\% & 38{,}4\,\% & 37{,}8\,\% & 37{,}6\,\% \\
        20\,\% & 39{,}9\,\% & 36{,}6\,\% & 36{,}0\,\% & 35{,}8\,\% \\
        \bottomrule
    \end{tabular}
\end{table}


Avec Greedy, la couverture passe de \SensCouvBase{} à \SensCouvAbs{} lorsque 20\,\% du personnel s'absente, et à \SensCouvPire{} si la demande double en même temps. La demande pèse moins que les absences au-delà de $\times 1{,}25$, car les lits saturés limitent l'activité.

% Généré par rapport/scripts/generate_tables.py : ne pas modifier à la main.
\begin{table}[H]
    \centering
    \caption{Postes cibles couverts par Naïve et Greedy selon les absences (demande $\times$1, aucune fermeture).}
    \label{tab:sensibilite_methodes}
    \small
    \begin{tabular}{lccccc}
        \toprule
        \textbf{Méthode} & \textbf{0\,\%} & \textbf{5\,\%} & \textbf{10\,\%} & \textbf{15\,\%} & \textbf{20\,\%} \\
        \midrule
        Naïve & 41{,}2\,\% & 39{,}7\,\% & 37{,}1\,\% & 35{,}7\,\% & 34{,}0\,\% \\
        Greedy & 48{,}3\,\% & 46{,}2\,\% & 44{,}1\,\% & 41{,}7\,\% & 39{,}9\,\% \\
        \bottomrule
    \end{tabular}
\end{table}


La stratégie naïve suit la même pente, plus bas : de \SensNaiveBase{} à \SensNaiveAbs{}.

\subsection{Effet sur les Refus}
% Généré par rapport/scripts/generate_tables.py : ne pas modifier à la main.
\begin{table}[H]
    \centering
    \caption{Admissions refusées par Greedy selon la demande et les fermetures (aucune absence ajoutée).}
    \label{tab:sensibilite_refus}
    \small
    \begin{tabular}{lcccc}
        \toprule
        \textbf{Lits fermés} & \textbf{Demande $\times$1{,}00} & \textbf{Demande $\times$1{,}25} & \textbf{Demande $\times$1{,}50} & \textbf{Demande $\times$2{,}00} \\
        \midrule
        0\,\% & 4 & 73 & 188 & 434 \\
        5\,\% & 14 & 105 & 227 & 473 \\
        10\,\% & 26 & 139 & 262 & 508 \\
        20\,\% & 79$^{\dagger}$ & 206$^{\dagger}$ & 329$^{\dagger}$ & 575$^{\dagger}$ \\
        \bottomrule
    \end{tabular}
    \par\vspace{3pt}
    {\footnotesize $^{\dagger}$ contraintes non respectées : les patients déjà présents dépassent les lits restants ; la valeur ne décrit pas un plan réalisable.}
\end{table}


Les refus de Greedy passent de \SensRefusBase{} à \SensRefusVc{} ($\times 1{,}25$), \SensRefusCq{} ($\times 1{,}5$) et \SensRefusDx{} ($\times 2$). Les absences n'ont aucun effet sur ce nombre, puisque Greedy accepte selon les seuls lits : c'est la couverture RH qui se dégrade. À l'inverse, fermer 10\,\% des lits porte les refus à \SensRefusFermee{} et fait \emph{monter} la couverture RH à \SensCouvFermee{}. Cette hausse ne traduit pas un meilleur fonctionnement : moins de patients sont accueillis.

Ces résultats ne concernent que Naïve et Greedy ; les solveurs n'ont pas été exécutés sur la grille.

% ------------------------------------------------------------------------------
\section{Discussion, Limites et Perspectives}
% ------------------------------------------------------------------------------

\subsection{Apports du Prototype}
Le prototype relie prévision, lits et personnel dans un même modèle ; compare quatre méthodes sur des instances identiques ; contrôle chaque plan de façon indépendante ; et rend visible l'arbitrage entre accueil et personnel, par les profils, la vue Allocation et le simulateur.

\subsection{Limites}
\label{sec:limites}
\begin{enumerate}
    \item \textbf{Sous-effectif structurel} : en $S_0$, Greedy couvre \CouvSzGreedy{} des postes cibles. Les lits ouverts par Greedy demandent \CibleSz{} postes sur la semaine, alors que \NbAgents{} agents à \PostesMax{} postes en offrent au plus \CapacitePostes{}, soit \PlafondSz{} de cette cible, avant toute absence. Il faudrait environ \AgentsCent{} agents pour la couvrir. CP-SAT dépasse légèrement ce taux (\CouvSzCp{}) parce qu'il ouvre moins de lits et abaisse donc sa propre cible. Les écarts entre méthodes portent donc sur la répartition d'une pénurie. Nous assumons cette situation comme celle d'un hôpital en tension ; un recalibrage des effectifs ou des cibles est une perspective, étudiée dans \texttt{docs/staffing\_calibration.md} et \textbf{non appliquée} ;
    \item \textbf{Qualifications en deux viviers} : \VivierA{} agents ne peuvent travailler qu'aux Urgences ou en Chirurgie, \VivierB{} qu'en Réanimation ou en Médecine. Le générateur donne à un agent sur deux une seule qualification, toujours Chirurgie ou Réanimation : la Médecine représente \PartCibleMed{} de la cible en $S_0$, mais \QualifMed{} seulement des agents y sont qualifiés. Aucun renfort entre les deux viviers n'est possible ;
    \item \textbf{Frontières de l'horizon} : les postes de la veille ne sont pas connus, donc le repos avant le premier jour et les nuits déjà enchaînées ne sont pas contrôlés. En fin d'horizon, refuser un patient le dernier jour ne coûte qu'une journée-patient, quelle que soit sa durée de séjour ;
    \item \textbf{Pas de contrainte d'heures} : la charge est limitée en nombre de postes ; les heures contractuelles et supplémentaires ne sont pas modélisées ;
    \item \textbf{Poids choisis par l'équipe} : les trois profils sont des hypothèses, non validées par des professionnels. D'autres poids donneraient d'autres plans ;
    \item \textbf{Non-déterminisme de CP-SAT} : à graine égale, un nouveau calcul peut donner un plan différent. Les écarts faibles entre solveurs ne doivent pas être surinterprétés. CBC, lui aussi limité en temps, dépend de la machine ;
    \item \textbf{Optimalité rarement prouvée} : un seul cas sur 18 l'est, pour chaque solveur ; avec la priorité personnel, MILP peut être moins bon que Greedy ;
    \item \textbf{Simplifications} : durée de séjour fixe, occupation journalière, cibles identiques dans tous les services et à toute heure, aucun transfert entre services ;
    \item \textbf{Données synthétiques} : les erreurs de prévision et les résultats ne se transposent pas à un hôpital réel ;
    \item \textbf{Prévision ponctuelle} : le plan ne tient pas compte de l'incertitude ;
    \item \textbf{Sécurité et usage} : pas d'authentification, pas d'intégration au SIH, utilisabilité non mesurée.
\end{enumerate}

\subsection{Perspectives d'Évolution}
\begin{itemize}
    \item recalibrer les effectifs ou les cibles, et corriger la génération des qualifications ;
    \item enchaîner les semaines, pour contrôler les frontières de l'horizon ;
    \item ajouter heures contractuelles et préférences du personnel ;
    \item prendre en compte l'incertitude des prévisions dans le plan ;
    \item fixer les poids avec des professionnels et tester l'outil avec eux ;
    \item authentifier les utilisateurs et lire les données par HL7 FHIR.
\end{itemize}

% ------------------------------------------------------------------------------
\section*{Conclusion du Chapitre}
\addcontentsline{toc}{section}{Conclusion du Chapitre}
% ------------------------------------------------------------------------------
Les \NbTestsPython{} tests Python et les \NbTestsNavigateur{} parcours de navigateur encadrent le prototype, et les 72 plans respectent les contraintes dures. CP-SAT est la méthode la plus régulière, MILP l'égale avec la priorité accueil, et Greedy fournit une bonne réponse immédiate. Surtout, les résultats montrent que le profil choisi déplace l'effort entre les patients et le personnel, dans un hôpital dont l'effectif ne couvre qu'environ la moitié de sa cible.

\chapter*{Conclusion Générale et Perspectives}
\addcontentsline{toc}{chapter}{Conclusion Générale et Perspectives}
\markboth{CONCLUSION GÉNÉRALE ET PERSPECTIVES}{CONCLUSION GÉNÉRALE ET PERSPECTIVES}

\section*{Synthèse des Réalisations du Projet}
Ce projet, mené au sein de la filière INDIA de l'ENSAM Rabat dans le module \textbf{Santé Digitale}, a conçu, implémenté et évalué un \textbf{prototype d'aide à la décision pour la planification conjointe des lits et du personnel hospitalier}. Ses contributions sont les suivantes :
\begin{enumerate}
    \item \textbf{Une chaîne de données synthétiques contrôlée} : \NbJours{} jours d'activité pour \NbLits{} lits et \NbAgents{} professionnels, sans aucune donnée personnelle réelle.
    \item \textbf{Une prévision évaluée chronologiquement} : cinq modèles comparés, avec un WAPE de validation de \WapeValUrg{} aux Urgences pour la régression à retards, et un test final qui nuance ce choix dans les autres services.
    \item \textbf{Un modèle unique des lits et du personnel} : cohortes de patients, dotation minimale, repos, qualifications, déficit de personnel et équité, avec une fonction objectif réglée par trois profils.
    \item \textbf{Une comparaison de quatre méthodes} : naïve, Greedy, MILP/CBC et CP-SAT, sur six scénarios et trois profils. Les 72 plans respectent les contraintes dures ; CP-SAT est la meilleure méthode, ou ex æquo, dans \NbCpMeilleur{} cas sur 18.
    \item \textbf{Une analyse des cas dégradés} : saturation, absences, crise combinée et fermeture de lits, complétées par une grille de sensibilité de \NbCas{} cas.
    \item \textbf{Un outil utilisable} : une API de \NbEndpoints{} points d'accès et un tableau de bord de dix vues, dont la couverture RH et un simulateur, vérifiés par \NbTestsPython{} tests Python et \NbTestsNavigateur{} parcours de navigateur.
    \item \textbf{Un cadre d'usage} : loi 09-08 et CNDP, HL7 FHIR, ISO/IEC 27001 et ISO 27799, rôles et processus de décision hebdomadaire.
\end{enumerate}

\section*{Enseignements}
\begin{itemize}
    \item \textbf{Un meilleur objectif n'est pas toujours un meilleur accueil} : avec la priorité personnel, CP-SAT réduit le déficit RH de \DeficitCpMinPer{} à \DeficitCpMaxPer{} par rapport à Greedy, mais refuse \RefusCpPerSq{} patients au lieu de \RefusGreedySq{} en $S_4$.
    \item \textbf{Respecter les contraintes ne suffit pas} : en situation nominale, \CouvDemSzGreedy{} de la demande est acceptée, mais \CouvSzGreedy{} seulement des postes cibles sont couverts.
    \item \textbf{L'optimisation répartit une pénurie, elle ne la supprime pas} : \NbAgents{} agents offrent au plus \CapacitePostes{} postes par semaine, pour une cible de \CibleSz{} postes dans le plan nominal de Greedy.
    \item \textbf{Une heuristique simple reste utile} : Greedy répond en moins de deux secondes et sert de référence et de repli.
    \item \textbf{Un contrôle indépendant est nécessaire} : il a conduit à renforcer la vérification des cohortes de patients et à ne plus qualifier un plan de « valide » sans préciser sa couverture.
\end{itemize}

\section*{Clôture et Perspectives}
Le prototype n'est pas prêt pour un usage hospitalier : les données sont synthétiques, les poids ont été choisis par l'équipe, l'optimalité est rarement prouvée et la sécurité reste à construire. Les suites naturelles sont le recalibrage des effectifs et des qualifications, l'enchaînement des semaines, la prise en compte de l'incertitude, la définition des poids avec des professionnels et l'intégration au système d'information hospitalier. Dans tous les cas, l'outil propose et explique ; la décision reste aux équipes de soins.


% Bibliographie et Webographie enrichie
\clearpage
\begin{thebibliography}{99}
\addcontentsline{toc}{chapter}{Bibliographie et Webographie}
\markboth{BIBLIOGRAPHIE ET WEBOGRAPHIE}{BIBLIOGRAPHIE ET WEBOGRAPHIE}

% ==============================================================================
% I. RECHERCHE OPÉRATIONNELLE ET PLANIFICATION HOSPITALIÈRE
% ==============================================================================
\subsubsection*{I. Recherche Opérationnelle et Planification Hospitalière}

\bibitem{hulshof2012}
P.~J. Hulshof, N.~Kortbeek, R.~J. Boucherie, E.~W. Hans et P.~J. Bakker,
\newblock \emph{Taxonomic classification of planning decisions in health care: a structured review of the state of the art in OR/MS},
\newblock Health Systems, vol.~1, n°~2, pp.~129--175, 2012.
\newblock DOI : \url{https://doi.org/10.1057/hs.2012.18}.

\bibitem{hallbeds}
R.~Hall,
\newblock \emph{Bed Assignment and Bed Management},
\newblock dans Handbook of Healthcare System Scheduling, Springer, pp.~177--200, 2012.
\newblock DOI : \url{https://doi.org/10.1007/978-1-4614-1734-7_8}.

\bibitem{burke2004}
E.~K. Burke, P.~De Causmaecker, G.~V. Berghe et H.~Van Landeghem,
\newblock \emph{The state of the art of nurse rostering},
\newblock Journal of Scheduling, vol.~7, n°~6, pp.~441--499, 2004.
\newblock DOI : \url{https://doi.org/10.1023/B:JOSH.0000046076.75950.0b}.

\bibitem{koukaras2026}
Koukaras et al.,
\newblock \emph{Operations research: a comprehensive literature review on nurse rostering optimisation},
\newblock Operational Research, vol.~26, 2026.
\newblock DOI : \url{https://doi.org/10.1007/s12351-026-01090-9}.

% ==============================================================================
% II. OPTIMISATION COMBINATOIRE ET PROGRAMMATION PAR CONTRAINTES
% ==============================================================================
\subsubsection*{II. Optimisation Combinatoire et Solveurs Mathématiques}

\bibitem{mitchell2011}
S.~Mitchell, M.~O'Sullivan et I.~Dunning,
\newblock \emph{PuLP: A Linear Programming Toolkit for Python},
\newblock The University of Auckland, 2011.
\newblock \url{https://github.com/coin-or/pulp}.

\bibitem{ortools}
Google,
\newblock \emph{OR-Tools : CP-SAT Solver},
\newblock Documentation officielle.
\newblock \url{https://developers.google.com/optimization/cp}.

\bibitem{cbc}
COIN-OR,
\newblock \emph{CBC : COIN-OR Branch and Cut},
\newblock Dépôt et documentation du projet.
\newblock \url{https://github.com/coin-or/Cbc}.

\bibitem{wolsey2020}
L.~A. Wolsey,
\newblock \emph{Integer Programming},
\newblock 2\textsuperscript{e} édition, John Wiley \& Sons, 2020.

\bibitem{rossi2006}
F.~Rossi, P.~van Beek et T.~Walsh (dir.),
\newblock \emph{Handbook of Constraint Programming},
\newblock Elsevier, 2006.

% ==============================================================================
% III. SÉRIES TEMPORELLES ET MODÉLISATION PRÉDICTIVE
% ==============================================================================
\subsubsection*{III. Science des Données et Séries Temporelles}

\bibitem{hyndman2021}
R.~J. Hyndman et G.~Athanasopoulos,
\newblock \emph{Forecasting: Principles and Practice},
\newblock 3\textsuperscript{e} édition, OTexts, 2021.
\newblock \url{https://otexts.com/fpp3/}.

% ==============================================================================
% IV. CADRE RÉGLEMENTAIRE ET NORMATIF
% ==============================================================================
\subsubsection*{IV. Cadre Réglementaire et Normatif}

\bibitem{loi0908}
Royaume du Maroc,
\newblock \emph{Loi n° 09-08 relative à la protection des personnes physiques à l'égard du traitement des données à caractère personnel},
\newblock Dahir n° 1-09-15 du 18 février 2009, Bulletin officiel n° 5714 du 5 mars 2009. Texte publié par la CNDP.
\newblock \url{https://www.cndp.ma/wp-content/uploads/2023/11/Loi-09-08-Fr.pdf}.

\bibitem{directive2003}
Parlement européen et Conseil,
\newblock \emph{Directive 2003/88/CE concernant certains aspects de l'aménagement du temps de travail},
\newblock Journal officiel de l'Union européenne, L~299, 2003, article~3.
\newblock \url{https://eur-lex.europa.eu/eli/dir/2003/88/oj}.

\bibitem{fhir}
HL7 International,
\newblock \emph{FHIR R4 (v4.0.1) : ressources Location, Encounter, Practitioner et Schedule}.
\newblock \url{https://hl7.org/fhir/R4/}.

\bibitem{iso27001}
ISO/IEC,
\newblock \emph{ISO/IEC 27001:2022 -- Systèmes de management de la sécurité de l'information -- Exigences},
\newblock notice de la norme.
\newblock \url{https://www.iso.org/standard/27001}.

\bibitem{iso27799}
ISO,
\newblock \emph{ISO 27799:2025 -- Informatique de santé -- Mesures de sécurité de l'information fondées sur ISO/IEC 27002},
\newblock notice de la norme.
\newblock \url{https://www.iso.org/standard/84647.html}.

% ==============================================================================
% V. WEBOGRAPHIE ET DOCUMENTATIONS TECHNIQUES OFFICIELLES
% ==============================================================================
\subsubsection*{V. Webographie et Documentations Logicielles Officielles}

\bibitem{fastapi}
FastAPI,
\newblock \emph{Documentation officielle}.
\newblock \url{https://fastapi.tiangolo.com/}.

\bibitem{pydantic}
Pydantic,
\newblock \emph{Documentation officielle (v2)}.
\newblock \url{https://docs.pydantic.dev/}.

\bibitem{nextjs}
Vercel,
\newblock \emph{Next.js : documentation officielle}.
\newblock \url{https://nextjs.org/docs}.

\bibitem{react}
Meta Platforms,
\newblock \emph{React : documentation officielle}.
\newblock \url{https://react.dev/}.

\bibitem{tailwindcss}
Tailwind Labs,
\newblock \emph{Tailwind CSS : documentation officielle}.
\newblock \url{https://tailwindcss.com/docs}.

\end{thebibliography}


% Annexes techniques
\clearpage
\chapter*{Annexes Techniques}
\addcontentsline{toc}{chapter}{Annexes Techniques}

\section*{Annexe A : Dictionnaire des Données et Configuration}
\markboth{ANNEXES TECHNIQUES}{ANNEXE A : DICTIONNAIRE DES DONNÉES}

\subsection*{Configuration de l'Établissement (\texttt{configs/hospital.json})}
L'extrait suivant est produit depuis le fichier de configuration :
\lstinputlisting[language=json, caption={Extrait de la configuration de l'établissement.}]{tables/hospital_extrait.json}

\subsection*{Structure de la Table des Admissions}
Le tableau~\ref{tab:dictionnaire} décrit les colonnes de \texttt{data/synthetic/historical\_demand.csv}.

% Généré par rapport/scripts/generate_tables.py : ne pas modifier à la main.
\begin{table}[H]
    \centering
    \caption{Colonnes de \texttt{historical\_demand.csv}.}
    \label{tab:dictionnaire}
    \small
    \begin{tabular}{lll}
        \toprule
        \textbf{Colonne} & \textbf{Type} & \textbf{Description} \\
        \midrule
        \texttt{date} & Date ISO & Jour observé \\
        \texttt{service} & Texte & Code du service : URG, MED, REA ou CHI \\
        \texttt{admissions} & Entier & Admissions du jour \\
        \texttt{discharges} & Entier & Sorties issues de la dynamique des séjours \\
        \texttt{emergency\_arrivals} & Entier ou vide & Arrivées aux urgences, renseignées pour URG seulement \\
        \texttt{is\_holiday} & Booléen & Jour férié synthétique \\
        \texttt{is\_special\_event} & Booléen & Événement exceptionnel tiré au hasard \\
        \texttt{data\_source} & Texte & Toujours synthetic \\
        \bottomrule
    \end{tabular}
\end{table}


\newpage
\section*{Annexe B : Guide d'Exploitation et Commandes}
\markboth{ANNEXES TECHNIQUES}{ANNEXE B : GUIDE D'EXPLOITATION}

Les commandes sont données pour Windows PowerShell, depuis la racine du projet.

\subsection*{Installation}
\begin{lstlisting}[caption={Installation de l'environnement.}]
python -m venv .venv
.\.venv\Scripts\python.exe -m pip install . pytest pytest-cov ruff httpx
$env:PYTHONPATH = "src"
cd frontend
npm.cmd install
\end{lstlisting}

\subsection*{Chaîne de Calcul}
\begin{lstlisting}[caption={Commandes de la chaîne de calcul.}]
# 1. Donnees synthetiques et rapport de qualite
.\.venv\Scripts\python.exe -m hospital_planning.data.generator

# 2. Comparaison des modeles et previsions J+1 a J+7
.\.venv\Scripts\python.exe -m hospital_planning.forecasting.evaluation

# 3. Instance operationnelle S0
.\.venv\Scripts\python.exe -m hospital_planning.operational.requirements

# 4. Comparaison complete (environ 32 minutes)
.\.venv\Scripts\python.exe -m hospital_planning.experiments.runner --methods naive greedy milp cp_sat --profiles equilibre priorite_accueil priorite_personnel

# 5. Grille de sensibilite, Naive et Greedy (environ 100 secondes)
.\.venv\Scripts\python.exe -m hospital_planning.experiments.sensitivity

# 6. Tableaux et chiffres de ce rapport
.\.venv\Scripts\python.exe rapport/scripts/generate_tables.py
\end{lstlisting}

\subsection*{API et Tableau de Bord}
\begin{lstlisting}[caption={Lancement de l'API et du tableau de bord (deux terminaux).}]
# Terminal 1 : API, documentation sur http://127.0.0.1:8000/docs
.\.venv\Scripts\python.exe -m uvicorn hospital_planning.api.app:app --host 127.0.0.1 --port 8000

# Terminal 2 : tableau de bord sur http://localhost:3011
cd frontend
npm.cmd run dev
\end{lstlisting}

\subsection*{Contrôles}
\begin{lstlisting}[caption={Tests et analyse statique.}]
.\.venv\Scripts\python.exe -m pytest -q
.\.venv\Scripts\python.exe -m ruff check src tests
cd frontend
npm.cmd run typecheck
npm.cmd run lint
npm.cmd run test
\end{lstlisting}

\newpage
\section*{Annexe C : Grille de Sensibilité Complète (80 Cas)}
\markboth{ANNEXES TECHNIQUES}{ANNEXE C : GRILLE DE SENSIBILITÉ}

Le tableau~\ref{tab:sensibilite_complete} donne les résultats de Greedy. Les cas dont les contraintes ne sont pas respectées ne décrivent pas un plan réalisable.

% Généré par rapport/scripts/generate_tables.py : ne pas modifier à la main.
{\footnotesize
\begin{longtable}{lcccrrrrc}
\caption{Les 80 cas de la grille de sensibilité, heuristique Greedy, profil équilibré.}\label{tab:sensibilite_complete}\\
\toprule
\textbf{Cas} & \textbf{Abs.} & \textbf{Dem.} & \textbf{Lits} & \textbf{Objectif} & \textbf{Refus} & \textbf{Déficit RH} & \textbf{Postes couv.} & \textbf{Contraintes} \\
\midrule\endfirsthead
\toprule
\textbf{Cas} & \textbf{Abs.} & \textbf{Dem.} & \textbf{Lits} & \textbf{Objectif} & \textbf{Refus} & \textbf{Déficit RH} & \textbf{Postes couv.} & \textbf{Contraintes} \\
\midrule\endhead
\bottomrule\endfoot
\texttt{A00\_D100\_C00} & 0\,\% & $\times$1{,}00 & -- & 636\,379 & 4 & 1\,200 & 48{,}3\,\% & respectées \\
\texttt{A00\_D100\_C05} & 0\,\% & $\times$1{,}00 & $-$5\,\% & 650\,419 & 14 & 1\,137 & 49{,}7\,\% & respectées \\
\texttt{A00\_D100\_C10} & 0\,\% & $\times$1{,}00 & $-$10\,\% & 674\,469 & 26 & 1\,096 & 50{,}6\,\% & respectées \\
\texttt{A00\_D100\_C20} & 0\,\% & $\times$1{,}00 & $-$20\,\% & 726\,459 & 79 & 901 & 55{,}3\,\% & non \\
\texttt{A00\_D125\_C00} & 0\,\% & $\times$1{,}25 & -- & 924\,327 & 73 & 1\,408 & 44{,}6\,\% & respectées \\
\texttt{A00\_D125\_C05} & 0\,\% & $\times$1{,}25 & $-$5\,\% & 967\,887 & 105 & 1\,309 & 46{,}4\,\% & respectées \\
\texttt{A00\_D125\_C10} & 0\,\% & $\times$1{,}25 & $-$10\,\% & 997\,969 & 139 & 1\,187 & 48{,}6\,\% & respectées \\
\texttt{A00\_D125\_C20} & 0\,\% & $\times$1{,}25 & $-$20\,\% & 1\,079\,619 & 206 & 984 & 53{,}2\,\% & non \\
\texttt{A00\_D150\_C00} & 0\,\% & $\times$1{,}50 & -- & 1\,246\,627 & 188 & 1\,450 & 43{,}9\,\% & respectées \\
\texttt{A00\_D150\_C05} & 0\,\% & $\times$1{,}50 & $-$5\,\% & 1\,284\,547 & 227 & 1\,312 & 46{,}3\,\% & respectées \\
\texttt{A00\_D150\_C10} & 0\,\% & $\times$1{,}50 & $-$10\,\% & 1\,318\,589 & 262 & 1\,190 & 48{,}6\,\% & respectées \\
\texttt{A00\_D150\_C20} & 0\,\% & $\times$1{,}50 & $-$20\,\% & 1\,403\,209 & 329 & 987 & 53{,}1\,\% & non \\
\texttt{A00\_D200\_C00} & 0\,\% & $\times$2{,}00 & -- & 1\,906\,827 & 434 & 1\,468 & 43{,}6\,\% & respectées \\
\texttt{A00\_D200\_C05} & 0\,\% & $\times$2{,}00 & $-$5\,\% & 1\,948\,707 & 473 & 1\,330 & 46{,}1\,\% & respectées \\
\texttt{A00\_D200\_C10} & 0\,\% & $\times$2{,}00 & $-$10\,\% & 1\,984\,219 & 508 & 1\,205 & 48{,}2\,\% & respectées \\
\texttt{A00\_D200\_C20} & 0\,\% & $\times$2{,}00 & $-$20\,\% & 2\,076\,279 & 575 & 1\,005 & 52{,}7\,\% & non \\
\texttt{A05\_D100\_C00} & 5\,\% & $\times$1{,}00 & -- & 660\,879 & 4 & 1\,249 & 46{,}2\,\% & respectées \\
\texttt{A05\_D100\_C05} & 5\,\% & $\times$1{,}00 & $-$5\,\% & 674\,919 & 14 & 1\,186 & 47{,}5\,\% & respectées \\
\texttt{A05\_D100\_C10} & 5\,\% & $\times$1{,}00 & $-$10\,\% & 698\,469 & 26 & 1\,144 & 48{,}4\,\% & respectées \\
\texttt{A05\_D100\_C20} & 5\,\% & $\times$1{,}00 & $-$20\,\% & 748\,959 & 79 & 946 & 53{,}1\,\% & non \\
\texttt{A05\_D125\_C00} & 5\,\% & $\times$1{,}25 & -- & 947\,327 & 73 & 1\,454 & 42{,}8\,\% & respectées \\
\texttt{A05\_D125\_C05} & 5\,\% & $\times$1{,}25 & $-$5\,\% & 990\,887 & 105 & 1\,355 & 44{,}5\,\% & respectées \\
\texttt{A05\_D125\_C10} & 5\,\% & $\times$1{,}25 & $-$10\,\% & 1\,020\,969 & 139 & 1\,233 & 46{,}6\,\% & respectées \\
\texttt{A05\_D125\_C20} & 5\,\% & $\times$1{,}25 & $-$20\,\% & 1\,103\,119 & 206 & 1\,031 & 51{,}0\,\% & non \\
\texttt{A05\_D150\_C00} & 5\,\% & $\times$1{,}50 & -- & 1\,269\,627 & 188 & 1\,496 & 42{,}1\,\% & respectées \\
\texttt{A05\_D150\_C05} & 5\,\% & $\times$1{,}50 & $-$5\,\% & 1\,307\,547 & 227 & 1\,358 & 44{,}5\,\% & respectées \\
\texttt{A05\_D150\_C10} & 5\,\% & $\times$1{,}50 & $-$10\,\% & 1\,341\,589 & 262 & 1\,236 & 46{,}6\,\% & respectées \\
\texttt{A05\_D150\_C20} & 5\,\% & $\times$1{,}50 & $-$20\,\% & 1\,426\,709 & 329 & 1\,034 & 50{,}9\,\% & non \\
\texttt{A05\_D200\_C00} & 5\,\% & $\times$2{,}00 & -- & 1\,929\,327 & 434 & 1\,513 & 41{,}9\,\% & respectées \\
\texttt{A05\_D200\_C05} & 5\,\% & $\times$2{,}00 & $-$5\,\% & 1\,971\,207 & 473 & 1\,375 & 44{,}2\,\% & respectées \\
\texttt{A05\_D200\_C10} & 5\,\% & $\times$2{,}00 & $-$10\,\% & 2\,007\,219 & 508 & 1\,251 & 46{,}3\,\% & respectées \\
\texttt{A05\_D200\_C20} & 5\,\% & $\times$2{,}00 & $-$20\,\% & 2\,099\,779 & 575 & 1\,052 & 50{,}5\,\% & non \\
\texttt{A10\_D100\_C00} & 10\,\% & $\times$1{,}00 & -- & 685\,877 & 4 & 1\,299 & 44{,}1\,\% & respectées \\
\texttt{A10\_D100\_C05} & 10\,\% & $\times$1{,}00 & $-$5\,\% & 699\,917 & 14 & 1\,236 & 45{,}3\,\% & respectées \\
\texttt{A10\_D100\_C10} & 10\,\% & $\times$1{,}00 & $-$10\,\% & 723\,467 & 26 & 1\,194 & 46{,}1\,\% & respectées \\
\texttt{A10\_D100\_C20} & 10\,\% & $\times$1{,}00 & $-$20\,\% & 772\,957 & 79 & 994 & 50{,}7\,\% & non \\
\texttt{A10\_D125\_C00} & 10\,\% & $\times$1{,}25 & -- & 972\,825 & 73 & 1\,505 & 40{,}8\,\% & respectées \\
\texttt{A10\_D125\_C05} & 10\,\% & $\times$1{,}25 & $-$5\,\% & 1\,016\,385 & 105 & 1\,406 & 42{,}4\,\% & respectées \\
\texttt{A10\_D125\_C10} & 10\,\% & $\times$1{,}25 & $-$10\,\% & 1\,046\,967 & 139 & 1\,285 & 44{,}4\,\% & respectées \\
\texttt{A10\_D125\_C20} & 10\,\% & $\times$1{,}25 & $-$20\,\% & 1\,127\,617 & 206 & 1\,080 & 48{,}6\,\% & non \\
\texttt{A10\_D150\_C00} & 10\,\% & $\times$1{,}50 & -- & 1\,295\,125 & 188 & 1\,547 & 40{,}1\,\% & respectées \\
\texttt{A10\_D150\_C05} & 10\,\% & $\times$1{,}50 & $-$5\,\% & 1\,333\,045 & 227 & 1\,409 & 42{,}4\,\% & respectées \\
\texttt{A10\_D150\_C10} & 10\,\% & $\times$1{,}50 & $-$10\,\% & 1\,367\,587 & 262 & 1\,288 & 44{,}3\,\% & respectées \\
\texttt{A10\_D150\_C20} & 10\,\% & $\times$1{,}50 & $-$20\,\% & 1\,451\,207 & 329 & 1\,083 & 48{,}6\,\% & non \\
\texttt{A10\_D200\_C00} & 10\,\% & $\times$2{,}00 & -- & 1\,955\,825 & 434 & 1\,566 & 39{,}9\,\% & respectées \\
\texttt{A10\_D200\_C05} & 10\,\% & $\times$2{,}00 & $-$5\,\% & 1\,997\,705 & 473 & 1\,428 & 42{,}1\,\% & respectées \\
\texttt{A10\_D200\_C10} & 10\,\% & $\times$2{,}00 & $-$10\,\% & 2\,032\,716 & 508 & 1\,302 & 44{,}1\,\% & respectées \\
\texttt{A10\_D200\_C20} & 10\,\% & $\times$2{,}00 & $-$20\,\% & 2\,124\,277 & 575 & 1\,101 & 48{,}2\,\% & non \\
\texttt{A15\_D100\_C00} & 15\,\% & $\times$1{,}00 & -- & 712\,876 & 4 & 1\,353 & 41{,}7\,\% & respectées \\
\texttt{A15\_D100\_C05} & 15\,\% & $\times$1{,}00 & $-$5\,\% & 726\,916 & 14 & 1\,290 & 42{,}9\,\% & respectées \\
\texttt{A15\_D100\_C10} & 15\,\% & $\times$1{,}00 & $-$10\,\% & 750\,466 & 26 & 1\,248 & 43{,}7\,\% & respectées \\
\texttt{A15\_D100\_C20} & 15\,\% & $\times$1{,}00 & $-$20\,\% & 800\,956 & 79 & 1\,050 & 47{,}9\,\% & non \\
\texttt{A15\_D125\_C00} & 15\,\% & $\times$1{,}25 & -- & 1\,002\,325 & 73 & 1\,564 & 38{,}4\,\% & respectées \\
\texttt{A15\_D125\_C05} & 15\,\% & $\times$1{,}25 & $-$5\,\% & 1\,045\,885 & 105 & 1\,465 & 40{,}0\,\% & respectées \\
\texttt{A15\_D125\_C10} & 15\,\% & $\times$1{,}25 & $-$10\,\% & 1\,073\,965 & 139 & 1\,339 & 42{,}0\,\% & respectées \\
\texttt{A15\_D125\_C20} & 15\,\% & $\times$1{,}25 & $-$20\,\% & 1\,155\,616 & 206 & 1\,136 & 46{,}0\,\% & non \\
\texttt{A15\_D150\_C00} & 15\,\% & $\times$1{,}50 & -- & 1\,324\,625 & 188 & 1\,606 & 37{,}8\,\% & respectées \\
\texttt{A15\_D150\_C05} & 15\,\% & $\times$1{,}50 & $-$5\,\% & 1\,362\,545 & 227 & 1\,468 & 40{,}0\,\% & respectées \\
\texttt{A15\_D150\_C10} & 15\,\% & $\times$1{,}50 & $-$10\,\% & 1\,394\,585 & 262 & 1\,342 & 42{,}0\,\% & respectées \\
\texttt{A15\_D150\_C20} & 15\,\% & $\times$1{,}50 & $-$20\,\% & 1\,479\,206 & 329 & 1\,139 & 45{,}9\,\% & non \\
\texttt{A15\_D200\_C00} & 15\,\% & $\times$2{,}00 & -- & 1\,984\,825 & 434 & 1\,624 & 37{,}6\,\% & respectées \\
\texttt{A15\_D200\_C05} & 15\,\% & $\times$2{,}00 & $-$5\,\% & 2\,026\,705 & 473 & 1\,486 & 39{,}7\,\% & respectées \\
\texttt{A15\_D200\_C10} & 15\,\% & $\times$2{,}00 & $-$10\,\% & 2\,060\,215 & 508 & 1\,357 & 41{,}7\,\% & respectées \\
\texttt{A15\_D200\_C20} & 15\,\% & $\times$2{,}00 & $-$20\,\% & 2\,152\,276 & 575 & 1\,157 & 45{,}5\,\% & non \\
\texttt{A20\_D100\_C00} & 20\,\% & $\times$1{,}00 & -- & 733\,875 & 4 & 1\,395 & 39{,}9\,\% & respectées \\
\texttt{A20\_D100\_C05} & 20\,\% & $\times$1{,}00 & $-$5\,\% & 747\,915 & 14 & 1\,332 & 41{,}0\,\% & respectées \\
\texttt{A20\_D100\_C10} & 20\,\% & $\times$1{,}00 & $-$10\,\% & 771\,465 & 26 & 1\,290 & 41{,}8\,\% & respectées \\
\texttt{A20\_D100\_C20} & 20\,\% & $\times$1{,}00 & $-$20\,\% & 821\,455 & 79 & 1\,091 & 45{,}9\,\% & non \\
\texttt{A20\_D125\_C00} & 20\,\% & $\times$1{,}25 & -- & 1\,025\,824 & 73 & 1\,611 & 36{,}6\,\% & respectées \\
\texttt{A20\_D125\_C05} & 20\,\% & $\times$1{,}25 & $-$5\,\% & 1\,069\,384 & 105 & 1\,512 & 38{,}1\,\% & respectées \\
\texttt{A20\_D125\_C10} & 20\,\% & $\times$1{,}25 & $-$10\,\% & 1\,095\,464 & 139 & 1\,382 & 40{,}2\,\% & respectées \\
\texttt{A20\_D125\_C20} & 20\,\% & $\times$1{,}25 & $-$20\,\% & 1\,176\,115 & 206 & 1\,177 & 44{,}0\,\% & non \\
\texttt{A20\_D150\_C00} & 20\,\% & $\times$1{,}50 & -- & 1\,348\,124 & 188 & 1\,653 & 36{,}0\,\% & respectées \\
\texttt{A20\_D150\_C05} & 20\,\% & $\times$1{,}50 & $-$5\,\% & 1\,386\,044 & 227 & 1\,515 & 38{,}0\,\% & respectées \\
\texttt{A20\_D150\_C10} & 20\,\% & $\times$1{,}50 & $-$10\,\% & 1\,416\,084 & 262 & 1\,385 & 40{,}1\,\% & respectées \\
\texttt{A20\_D150\_C20} & 20\,\% & $\times$1{,}50 & $-$20\,\% & 1\,499\,705 & 329 & 1\,180 & 44{,}0\,\% & non \\
\texttt{A20\_D200\_C00} & 20\,\% & $\times$2{,}00 & -- & 2\,008\,823 & 434 & 1\,672 & 35{,}8\,\% & respectées \\
\texttt{A20\_D200\_C05} & 20\,\% & $\times$2{,}00 & $-$5\,\% & 2\,050\,703 & 473 & 1\,534 & 37{,}8\,\% & respectées \\
\texttt{A20\_D200\_C10} & 20\,\% & $\times$2{,}00 & $-$10\,\% & 2\,081\,714 & 508 & 1\,400 & 39{,}9\,\% & respectées \\
\texttt{A20\_D200\_C20} & 20\,\% & $\times$2{,}00 & $-$20\,\% & 2\,172\,775 & 575 & 1\,198 & 43{,}6\,\% & non \\
\end{longtable}
}



\end{document}
